%========================================
%. CONFIGURATION ARTICLE  COURS  SEC 101
%========================================
% 				CONTRIBUTIONS
%========================================

%************************************************************
% Chargement des variables du modèle
%************************************************************

%-------------------------------------------------------------
%               FR CYBERDEF SECOPS COURSE
%                          CONFIG NOTES File
%          SET BEAMER or Article/book Mode
%                    PRZMODE OR DOCMODE
%                             2020 eduf@ction
%-------------------------------------------------------------

% insert InterTitlle for each file compose main "Notes"
%\newcommand{\INTERTITLE}{OK}

 %------------------------------------------
 % Use global project CYBERDEF 
 %            Configuration file
 %------------------------------------------
%-----------------------------------------------
%               FR CYBERDEF SECOPS COURSE
%                         GLOBAL CONFIG  File
%                             2020 eduf@ction
%-----------------------------------------------

% just used for quick access to samples on project editing tool
%%-------------------------------------------------------------
%               FR CYBERDEF SECOPS COURSE
%              SAMPLES for edx & u  commands
%                             2020 eduf@ction
%-------------------------------------------------------------


% 												I M A G E S   
%-------------------------------------------------------------------------
%\upicture{../Latex/Sources/EDU/Pictures/img-liovar-incidents}{Incidents}{01}{lb:mtbf}
%-------------------------------------------------------------------------


%\mode<all>{\picframe{../Latex/Sources/EDU/Pictures/img-cycle}{Cycle de vie de gouvernance Cyberdef}{0.3}{lbl:cycle}}


% Planche PRZ 
% begin frame \mode<all>{\texframe{Déroulement}{présentation Gestion de la menace}
%{

%}} % end frame



%\begin{table}[!hbtp]
%\resizebox{\textwidth}{!}{%
%\begin{tabular}{ l|l l}
%%\begin{tabular}{l*{3}{S[table-format=2.1]}}
%\toprule
%\textbf{Package\_name }                 & \textbf{C\&C }                                        & \textbf{sha256}                                                            \\ \hline
%caracal.raceinspace.astronaut & http://api.lulquid.xyz                       & f1d32c17a169574369088...\\ \hline
%
%\bottomrule
%\end{tabular}
%}
%\caption{Exemple de  données de TI (Threat Intelligence)}
%\label{tab:ExempleTI}
%\end{table}

%----------------------------
% Path
\def\upath{../Latex}

%-------------------------------------------------------------
%			 FR CYBERDEF SECOPS COURSE
% 					INIT Commands File
% 		Select BEAMER or Article/book Mode
% 			with sub extension .doc or .prz
%                           2020 eduf@ction
%-------------------------------------------------------------
\RequirePackage{xstring}
\RequirePackage[realmainfile]{currfile}

\newcommand{\udocname}{\jobname}

\IfSubStr*{\currfilebase}{.prz}{\newcommand{\PRZMODE}{YES}}{\newcommand{\DOCMODE}{YES}}
   
\IfSubStr*{\currfilebase}{.edt}{\newcommand{\EDTMODE}{YES}}{\newcommand{\NOMODE}{YES}}

\input{../Latex/Templates/Magic/commons.inc/edx.library.tex}

\def\utemplatepath{../Latex/Templates/Magic}

\def\edxlibspath{\utemplatepath/edxlibs}
%\input{\edxlibspath/load.tex} à mettre dans les class article, book ...

\def\uauthorphoto{\utemplatepath/commons.inc/CommonsPictures/missionauthorphoto.png}


\def\utemplatetheme{cnam}
\def\ubeamermodel{edx}
\def\uproject{CNAMSEC101}


\newcommand{\upictureext}{pdf}

 %----------------------------
 % Basic Var
 
 %\setcounter{cntx}{\year+1}
\newcommand{\INFODistrib}{Notes de cours SECOPS 2022-2023}
\newcommand{\uJournalInfo}{Orange Cyber et Cloud avec CNAM Bretagne, Cybersécurité SEC101, edufaction}

\newcommand{\udescription}{La sécurité opérationnelle au coeur de la cyberdéfense d'entreprise}

\newcommand{\uUniversity}{Conservatoire National des Arts et Métiers}
\newcommand{\uCompany}{Orange Campus Cyber}
 %\newcommand{\uJournalInfo}{CNAM Bretagne, Cybersécurité SEC101, eduf@ction}
\newcommand{\uinstituteshort}{CNAM}
\newcommand{\uinstitute}{\uUniversity}
\newcommand{\uchaire}{Chaire de Cybersécurité}
\newcommand{\uversion}{Cnam Bretagne}





%----------------------------
% All other variables are defined in
 % file : *.lesson.tex
\newcommand{\ushorttitle}{CYBERDEF SEC101}
\newcommand{\uCoursetittle}{Eléments de Sécurité Opérationnelle}
\newcommand{\uMinilogo}{Cyberdéfense d'entreprise}
\newcommand{\uCourseLongName} {Eléments de sécurité opérationnelle en cyberdéfense d'entreprise}
%\newcommand{\utitle}{Cyberdéfense d'entreprise}

%----------------------------
% Book versus Course Notes
%----------------------------
% Use "e document' "e cours"
\newcommand{\edoc}{e document\xspace}
\newcommand{\ecours}{e cours\xspace} 
\newcommand{\etitle}{éléments de sécurité opérationnelle en cyberdéfense d'entreprise} 
\newcommand{\fichetech}{une fiche TECHNO}


%%----------------------------
%% Spécifique BOOK
% %----------------------------
\newcommand{\printer}{Publication limitée pour le \uUniversity \xspace et Orange Campus Cyber}

%%----------------------------
% Attention : Use hard coded Pictures Path
% (PicframeS)
\newcommand{\picpath}{../Latex/Sources/EDU/SRC1/Pictures}
%-------------------------------------------------------------

%%----------------------------
%% Spécifique COURSE
% %----------------------------

\newcommand {\parttitlecourse}{TITRE PARTIE (config)}
\newcommand{\coursenumber}{1}

\newcommand{\uGlossaryAcronyms}{
\loadglsentries{../Latex/Sources/commons/acronymes.tex}
\loadglsentries{../Latex/Sources/commons/glossaire.tex}
}

\newcommand{\ucourseresume}{\normalsize Aborder la sécurité des systèmes d’information sous l’angle d'une sécurité dynamique est un axe qui depuis quelques années apporte de nouvelle manière d’aborder la protection, la défense, et la résilience des systèmes d’information. La transformation digitale de l'entreprise modifie et rend plus flous les périmètres des systèmes d’informations. Cela nécessite une approche élargie du risque numérique et des nouvelles architectures de systèmes d'information intégrant des technologies de cloud et de cyberdéfense. Malgré la mise en place de mesures et de technologies de protection de plus en plus élaborées, l'impact d'une attaque ayant franchi ces frontières poreuses à considérablement augmenté.  Cette compilation des notes de cours élaborée dans le cadre d'un cours d'introduction à la cybersécurité et la cyberdéfense d'entreprise aborde les grands éléments fondamentaux permettant d'appréhender le domaine. Il permet d'en comprendre les enjeux, les codes, et acteurs. Protéger l'ensemble de l'entreprise alors qu'il est complexe de définir ses frontières est illusoire. Identifier les actifs essentiels ou vitaux et mettre en place les moyens adaptés à leur protection et leur défense est une démarche tactique qui permet de graduellement réduire ses cyber-risques tout en assurant un pilotage globale de la gouvernance de la sécurité sur ses volets de conformité.}

 %------------------------------------------
 
\renewcommand{\uMinilogo}{\includegraphics[width=0.04\paperwidth]{\utemplatepath/commons.inc/CommonsPictures/\utemplatetheme.theme/shield-20.pdf}}
 
\newcommand {\Ucontribute}
{
%%-------------------------------------------------------------
%               FR CYBERDEF SECOPS COURSE
%                            Contributions Rules
%                              2020 eduf@ction
%-------------------------------------------------------------

% \section{Contributions}

% \subsection{Comment contribuer}

Les notes et les présentations sont réalisées sous \hologo{LaTeX}. 

Vous pouvez contribuer au projet du cours CYBERDEF101. Les contributions peuvent se faire sous deux formes :

\begin{itemize}
  \item Corriger, amender, améliorer les notes publiées. A chaque session des modifications et évolutions sont apportées pour tenir compte des corrections de fond et de formes. 
  \item Ajouter, compléter, modifier des parties de notes sur la base de votre lecture du cours et de vos expertises dans chacun des domaines évoqués.
\end{itemize}

 Les fichiers sources sont publiés sur GITHUB dans l'espace : \ulink{https://github.com/edufaction/CYBERDEF101}{(edufaction/CYBERDEF101)}.
 
%%-------------------------------------------------------------
%               FR CYBERDEF SECOPS COURSE
%                Contributions DATAS & NAMES
%                              2020 eduf@ction
%-------------------------------------------------------------

\subsection{Les contributeurs/auteurs du cours}

\subsubsection{co-auteurs}

%-------------------------------------------------------------
%               FR CYBERDEF SECOPS COURSE
%                Contributions DATAS & NAMES
%                              2020 eduf@ction
%-------------------------------------------------------------
%           CO AUTEURS - Rédacteur de chapitre
%-------------------------------------------------------------

\coauthor{David  BATANY}{2019-2020}{Architecture et fonctionnement des Botnets}{Cnam SEC101}

\subsubsection{contributeurs}

%-------------------------------------------------------------
%               FR CYBERDEF SECOPS COURSE
%                Contributions DATAS & NAMES
%                              2020 eduf@ction
%-------------------------------------------------------------
%                 CONTRIBUTIONS AU COURS
%-------------------------------------------------------------


\coauthor{Céline JUBY}{2020}{Contributions d'amélioration et relectures}{Orange Cyberdefense}


} 

%------------------------------------------
 %     BEAMER OU ARTICLE MODEL
 %------------------------------------------
  
\ifdefined\PRZMODE 
	\providecommand{\umainload}{%-------------------------------------------------------------
%               FR CYBERDEF SECOPS COURSE
%                          BEAMER MAIN FILE
%                             2020 eduf@ction
%-------------------------------------------------------------

%\begin{frame}
%\frametitle<presentation>{TITRE DIAPO}
% frame text
%\end{frame}

%-------------------------------------------------------------
\documentclass[ignorenonframetext, 8pt, allowframebreaks]{beamer}

%-----------------------------------------
% Load Edxlibs in MAGIC path
%-----------------------------------------
% Need : \edxlibspath defined

\input{\edxlibspath/edxlib-placeholder.tex}
\input{\edxlibspath/edxlib-epigraphe.tex}
\input{\edxlibspath/edxlib-quizzmoodle.tex}


\usepackage[utf8]{inputenc}
\usepackage[T1]{fontenc}
\usepackage{csquotes}
\usepackage[french]{babel}
\usepackage{lmodern}
\usepackage{multicol}

\input{\utemplatepath/commons.inc/CommonsPictures/\utemplatetheme.theme/colors.inc} % set uMainColor

%-------------------------------------------------------------
%               FR CYBERDEF SECOPS COURSE
%                                  EDX PACKAGES
%.    Commons packages for BOOK, NOTES (.doc) 
%                            and BEAMER (.prz)
%                              2020 eduf@ction
%-------------------------------------------------------------

%\usepackage[hyphens,spaces,obeyspaces]{url}
%\usepackage[hyphens,spaces,obeyspaces]{url}



\PassOptionsToPackage{hyphens,spaces,obeyspaces}{url}

\usepackage[utf8]{inputenc}
\usepackage[french]{babel}
\usepackage[T1]{fontenc}
%\usepackage{markdown}



\usepackage{ccicons}

\usepackage{tikz} % Required for drawing custom shapes
\usetikzlibrary{tikzmark} 
\usetikzlibrary{mindmap,trees, backgrounds}


\usepackage{verbatim}
\usepackage{wrapfig}
\usepackage{xpatch}

\usepackage{hologo}

\usepackage[figure]{totalcount}

\usepackage{fontawesome}

\usepackage[skins]{tcolorbox} % don't forget SKINS option

\usepackage{listings}
\usepackage{upquote}

\lstset{
upquote=true,columns=flexible, basicstyle=\ttfamily,
language=HTML, 
frameround=tttt,
commentstyle=\color{gray},
identifierstyle=\color{blue},
keywordstyle=\color{uMainColor}\bfseries,
xleftmargin=2em,
xrightmargin=2em,
aboveskip=\topsep,
belowskip=\topsep, 
frame=single,
rulecolor=\color{uMainColor},
backgroundcolor=\color{uMainColor!5},
breaklines,
breakindent=1.5em,
showspaces=false,
showstringspaces=false,
showtabs=false,
xleftmargin=2em,
xrightmargin=2em,
aboveskip=1em,
belowskip=1em,
}



\usepackage{graphicx} % Required for including pictures
%\graphicspath{{Pictures/}} % Specifies the directory where pictures are stored

\usepackage{booktabs} % Required for nicer horizontal rules in tables

\usepackage{setspace}
\setstretch{1,2}

\usepackage{csquotes}

\usepackage{fancyhdr}

\usepackage{makeidx} 

\usepackage{xspace} 

\usepackage{floatflt} 

\usepackage{hyperref}
\usepackage{cleveref}

%pdfdisplaydoctitle={\utitle}, Attention le titre ne doit pas contenir de saut de ligne 

\hypersetup{
pdftitle={\uCoursetittle},
pdfauthor={\uauthor},
hidelinks,
backref=true,
pagebackref=true,
hyperindex=true,
pdfcenterwindow = true,
colorlinks=true,
breaklinks=true,
bookmarks=true,
allcolors=uMainColor,
pdfsubject = {\udescription},
pdfcreator = {MACTEX on TEXPAD},
bookmarksopen=true}


% urlcolor=cnam,
% runcolor=uMainColor,
%\pdfstringdefDisableCommands{%
%  \def\\{}%
%  \def\texttt#1{<#1>}%
%}

\usepackage[style=numeric,citestyle=numeric,sorting=nyt,sortcites=true,autopunct=true,hyperref=true,abbreviate=false,backref=true,backend=biber]{biblatex}

%\usepackage{cleveref} old position

\usepackage{multicol}

%--------------------------------------------------------
% pour la génération du QRCODE en page de garde d'article
\usepackage{qrcode} 

%--------------------------------------------------------
% POur la date d'impression en premiere page d'article
\usepackage[useregional]{datetime2}

%--------------------------------------------------------
% Gestion des glossaires et Sources/EDU/acronymes
% Un fichier acroymes.tex centralisé
\usepackage[acronym]{glossaries}

%\makeglossaries
\makenoidxglossaries 

\loadglsentries{Latex/Sources/commons/acronymes.tex}
\loadglsentries{Latex/Sources/commons/glossaire.tex}


%\renewcommand{\bibname}{R\'ef\'erences}

%--------------------------------------------------------
% Pour exclure des "environnement", ici techworkbox par exemple, travaux pour étudiants, non affiché dans le BOOK.
\usepackage{comment}


%--------------------------------------------------------
% pour filigramme sur les pages
\usepackage{eso-pic}

%--------------------------------------------------------
% for \gappto{\UrlBreaks}{\UrlOrds}
% rather than 
%\makeatletter
% \g@addto@macro{\UrlBreaks}{\UrlOrds}
% \makeatother

\usepackage{etoolbox}

\usepackage{breakurl}


\newcommand{\ucite}[1]{\cite{#1}}
%\newcommand{\ucite}[1]{REF BIB #1}
\newcommand{\ugls}[1]{\gls{#1}}
%\newcommand{\ugls}[1]{GLOSSAIRE (#1)}
%\usepackage[Rejne]{fncychap}

%\usepackage{pythontex} % need shell escape

%
%\makeatletter
%\DeclareRobustCommand{\afilename}[1]{%
% \begingroup
%  % \lstname seems to change hyphens into \textendash
%  \def\textendash{-}%
%  \afilename@parse{#1}%
%  \edef\afilename@base{\detokenize\expandafter{\afilename@base}}%
%  \texttt{\afilename@base.\afilename@ext}%
% \endgroup
%}
%\makeatother

%-------------------------------------------------------------
%               FR CYBERDEF SECOPS COURSE
%                                  EDX PACKAGES
%.    Commons COMMANDS for BOOK, NOTES (.doc) 
%                            and BEAMER (.prz)
%                              2020 eduf@ction
%-------------------------------------------------------------

\newcommand{\ucite}[1]{\cite{#1}}
%\newcommand{\ucite}[1]{REF BIB #1}
\newcommand{\ugls}[1]{\gls{#1}}
%\newcommand{\ugls}[1]{GLOSSAIRE (#1)}
\newcommand{\uac}[1]{\acrshort{#1}}



% \newcommand{\ExtractBeforeLastSlash}[1]{\directlua{extract_before_last_slash("#1")}}
% \newcommand{\ExtractAfterLastSlash}[1]{\directlua{extract_after_last_slash("#1")}}
% \newcommand{\RemoveExtension}[1]{\directlua{remove_extension("#1")}}
% \newcommand{\ExtractExtension}[1]{\directlua{extract_extension("#1")}}


\newcommand{\includepicturesubpath}[1]
{#1}

%{\ExtractBeforeLastSlash{#1}/\upicturethemepath\RemoveExtension{\ExtractAfterLastSlash{#1}}.\ExtractExtension{#1}}
    

%---------------------------------------------------
%										 	AUDIO NOTE 
%---------------------------------------------------

\newcommand\audionote[2]{%
  \note{#2}
  \immediate\write\audioexport{==AudioSlide=============}
  \immediate\write\audioexport{#1}
  \immediate\write\audioexport{\insertframenumber}
  \immediate\write\audioexport{---------------------------------}
  \immediate\write\audioexport{#2}
  \immediate\write\audioexport{---------------------------------}
}%

%---------------------------------------------------
%										 	TEXT STYLE SHORTCUT 
%---------------------------------------------------

\newcommand{\tcb}[1]{\textbf{\textcolor{uMainColor}{#1}}}

\newcommand{\voir}[1]{(Schéma~\ref{#1},~page~\pageref{#1}}

\newcommand{\graphicxextension}{pdf}

\newcommand\myemptypage{
   
    \thispagestyle{empty}
    \addtocounter{page}{-1}
    \newpage
  \null
    %CECI EST UNE PAGE VIDE
       \vfill
    }


\newcommand{\uNA}[2]
{
\newacronym{a#1}{#1}{#2}
}

\newcommand{\coauthor}[4]
{
\noindent (#2) \textbf{#1} - #4 : \textit{#3}
}


\newcommand\UKword[1]{\emph{#1}\xspace}
\newcommand\textUK[1]{\textit{\textbf{#1}}\xspace}
\newcommand\Ucenter[1]{\begin{center}{#1}\end{center}}

\newcommand\head[1]{\textbf{#1}}
\newcommand\tb[1]{\textbf{#1}}
\newcommand{\g}[1]{\og #1 \fg{}}
\newcommand{\uload}[1] {\input{../Latex/#1}}
\newcommand{\ucurcolor}{black}

\newcommand{\mak}[1]{\faicon{\aArrowCircleORight} #1}


%-----------------------------------------------------
%										 	\INFOGitHub 
%---------------------------------------------------


\newcommand{\GITfilename}{https://github.com/edufaction/CYBERDEF101/raw/main/Release/Builder/\jobname.pdf}

%\newcommand{\safeqrcode}[2][]{%
%\qrcode[#1]{\detokenize{#2}} }

\newcommand{\safeqrcode}[2][]{\qrcode[#1]{#2}} 

\newcommand{\INFOGitHub}
{
\begin{center} 
\setstretch{1.5}
 {Vérifiez la disponibilité d'une version plus récente de} \\
 {\ulink{\GITfilename}{\textbf{\jobname.pdf }sur GITHUB CYBERDEF} }   \\ 
{{\huge\ccbyncndeu}}  \\  
{Publication en \href{https://creativecommons.org/licenses/by-nc-nd/2.0/fr/}{Creative Common BY-NC-ND} by eduf@ction}    \\  % Copyright notice
{\safeqrcode[height=4 cm,padding]{\GITfilename}}  
\end{center} 
}

\newcommand{\INFOGitHubS}
{
\noindent{\safeqrcode[height=3 cm]{\GITfilename}} \\
\noindent {\tiny Téléchargez une version à jour} 
}



%-----------------------------------------------------
%										 	\upspicture 
%---------------------------------------------------
\newcommand{\upspicture}[3]
{
\renewcommand{\ucurcolor}{#3}
\resizebox{#2\textwidth}{!}{\input{../#1}}
}


%-----------------------------------------------------
%										 	\uindex 
%---------------------------------------------------


\newcommand{\uindex}[1]
{
#1\index{#1}\xspace
}

\newcommand{\ulindex}[1]{\index{#1}}

%-----------------------------------------------------
%										 	\includer 
%---------------------------------------------------

\newcommand{\includer}[1]
{
\input{../Latex/Sources/EDU/SRC1/Chapters/#1}
}

%-----------------------------------------------------
%										 	\edxdico 
%---------------------------------------------------

\newcommand{\edxdico}[2] {\textbf{#1}:\textit{#2} \index{#1}}

%-----------------------------------------------------
%  command DEF : InTextx
%-----------------------------------------------------
\newcommand\InTexta{no text}
\newcommand\InTextb{no text}
\newcommand\InTextc{no text}
\newcommand\rnc{\renewcommand}


%-----------------------------------------------------
%										 	ENV warningbox
%---------------------------------------------------

\newtcolorbox{warningbox}[2][]
{
  colframe = uMainColor!25,
  colback  = black!05,
  coltitle = uMainColor!20!black,
  fonttitle = \bfseries, 
  title    = #2,
  #1,
}


\newtcolorbox{techworkbox}[2][]{sharp corners,colframe = uMainColor!25,colback  = black!05,coltitle = uMainColor!20!black,fonttitle = \bfseries, title={\textcolor{uMainColor}{\faFile~\small {Sujet de mémoire SECOPS}}} : #2,#1,}


\newtcolorbox{notebox2}[2][]
{
sharp corners, 
colback = uMainColor!5!white, 
colframe = uMainColor!75!black,
fonttitle = \bfseries, 
colbacktitle= uMainColor!85!black,
title=#2,#1
}

\newtcolorbox{notebox}[2][]{colback=uMainColor!5!white,
colframe=uMainColor!75!black,fonttitle=\bfseries,
colbacktitle=uMainColor!85!black,enhanced,
attach boxed title to top right={yshift=+1mm},
title=#2,#1}


\newtcolorbox{toolsboxin}[2][]{sharp corners, colback=uMainColor!0!white,
colframe= black!25,fonttitle=\bfseries,
colbacktitle=black!60,enhanced,
attach boxed title to top left={yshift=+1mm},
title=#2,#1}


\newtcolorbox{techworkin}[2][]{sharp corners, colback=uMainColor!5!white,
colframe= uMainColor!100, fonttitle=\bfseries,
colbacktitle=black!5, enhanced,
 coltitle=black!50,
attach boxed title to top right ={yshift=+1mm},
title=#2,#1}



%-----------------------------------------------------
%										 	ENV ToolsBox
%---------------------------------------------------


\newcommand\tooleditor{No EDITOR in Datas.tex}
\newcommand\toolurl{No URL INSERTED in Datas.tex}
\newcommand\toolclass{No EDITOR in Datas.tex}
\newcommand\toolname{No NAME in Datas.tex}
\newcommand\toolanalyst{No ANALYST in Datas.tex}

\newcommand{\toolsbox}[2]
{
\input{../Latex/Sources/EDU/SRC1/SecComponents/#2.#1/datas.tex} 
\begin{toolsboxin}[sidebyside,righthand width=.7\textwidth]{#1}
\includegraphics[width=0.9\textwidth]{../Latex/Sources/EDU/SRC1/SecComponents/#2.#1/icon.png} \\
\tcblower
\input{../Latex/Sources/EDU/SRC1/SecComponents/#2.#1/description.tex}
\end{toolsboxin}
\noindent \small {\faCogs~Classe : \uppercase{\textbf{#2}}, Site de référence : \textbf{\ulink{\toolurl}{#1}}}
\newline
\small {\faGlobe~Editeur : \textbf{\tooleditor}}
\small {\faEye~Analyste : \textbf{\toolanalyst}}
}

%\renewcommand\toolsclass{SIEM}
%\renewcommand\toolsname{RSA Netwitness}
%\renewcommand\toolseditor{RSA}
%\renewcommand\toolsurl{https://www.rsa.com/fr-fr/products/threat-detection-response}

%-----------------------------------------------------
%										 	ENV  TechWork
%---------------------------------------------------


\newcommand{\techwork}[2]
{
\begin{techworkin} [righthand width=.8\textwidth]{\faCogs~\small {Fiche TECHNO} : #1}#2
\end{techworkin}
}

%-----------------------------------------------------
%										 	\uchap
%---------------------------------------------------
\ifdefined\INTERTITLE
\providecommand{\uchap}[1]
	{
	\begin{LARGE}
	\textbf{\Ucenter{#1 \\ -oOo-}}	
	\end{LARGE}
	}
\else
	\providecommand{\uchap}[1]{}
\fi

%-----------------------------------------------------
%											\uexpand
%---------------------------------------------------
\newcommand {\uexpand}[1]{
\mode<all>\input{../Latex/Sources/EDU/SRC1/Chapters/#1}
}


%-----------------------------------------------------
%											\spicture 
% direct include graphic with Scrivener
%-----------------------------------------------------

\newcommand{\spicture}[3]{
\begin{figure}[!hbtp] %h !hbtp
  \begin{center}
	 \includegraphics[width=0.8\textwidth]{#1}
  \end{center}
\caption{#3 - (slbl-#2)\label{slbl-#2}}
\end{figure}
}

%\StrBefore{#2}{.pdf}

%-----------------------------------------------------
%											\upicture 
%---------------------------------------------------

\ifdefined\PRZMODE

\providecommand{\upicture}[4]{
\framesubtitle{#2}
%\begin{figure}[!h] %h
  \begin{center}
	 \includegraphics[height=#3\textheight, width=#3\textwidth, keepaspectratio ]{\includepicturesubpath{#1.pdf}}
  \end{center}
%  \end{figure}
}

\else

\providecommand{\upicture}[4]{
\begin{figure}[!hbtp] %h !hbtp
  \begin{center}
	 \includegraphics[width=#3\textwidth]{\includepicturesubpath{#1.pdf}}
  \end{center}
\caption{#2 \label{#4}}
\end{figure}
}
	
\fi


\ifdefined\PRZMODE

\providecommand{\uwpicture}[4]{
\framesubtitle{#2}
%\begin{figure}[!h] %h
  \begin{center}
	 \includegraphics[height=0.8\textheight, width=0.8\textwidth, keepaspectratio ]{\includepicturesubpath{#1.pdf}}
  \end{center}
%  \end{figure}
}
\else
\providecommand{\uwpicture}[4]
{
\begin{wrapfigure}{O}{#3\textwidth} % inside / Outside
  \begin{center}
	 \includegraphics[width=#3\textwidth]{\includepicturesubpath{#1.pdf}}
\end{center}
\caption{#2 \label{#4}}
\end{wrapfigure}
}
\fi

%-----------------------------------------------------
%	\picframe for BEAMER (Image included in Beamer Slide)
%---------------------------------------------------

\ifdefined\PRZMODE

\providecommand{\picframe}[4]
{
\begin{frame}
\frametitle<presentation>{#2}
  \begin{center}

   \begin{tikzpicture}[remember picture, overlay]
    \node at (current page.center) {\includegraphics[width=#3\paperwidth, height=#3\paperheight,keepaspectratio]{\includepicturesubpath{#1.pdf}}};
    \end{tikzpicture}

  \end{center}
\end{frame}

}
\else
\providecommand{\picframe}[4]
{
\begin{figure}[!hbtp] % !hbtp ou !h 
\centering 
%{\includegraphics[width=#3\paperwidth, height=#3\paperheight,keepaspectratio]{\includepicturesubpath{#1.pdf}}}
\includegraphics[width=#3\textwidth]{\includepicturesubpath{#1.pdf}}
\caption{#2 \label{#4}}
\end{figure}
}
\fi

%-----------------------------------------------------
%.  for BEAMER (Image included in Beamer Slide)
%	\picframeS Simplified with Name and Label from One Param
%---------------------------------------------------


\ifdefined\PRZMODE

\providecommand{\picframeS}[2]
{
\begin{frame}
\frametitle<presentation>{#2}
  \begin{center}
	 \includegraphics[height=0.95\textheight, width=0.95\textwidth, keepaspectratio ]{\picpath/img-#1.\graphicxextension}
  \end{center}
\end{frame}

}
\else
\providecommand{\picframeS}[2]
{
\begin{figure}[!hbtp] % !hbtp ou !h 
\centering \includegraphics[width=0.9\textwidth]{\picpath/img-#1.\graphicxextension}
\caption{(s) #2 \label{lbl:#1}}
\end{figure}
}

\fi


%{figure}[!h]
%{wrapfigure}{O}{#3\textwidth}
%-----------------------------------------------------
%	FRAME			\wpicframe 
%---------------------------------------------------

\ifdefined\PRZMODE

\providecommand{\wpicframe}[4]
{
\begin{frame}
\frametitle<presentation>{#2}
  \begin{center}
	 \includegraphics[height=0.95\textheight, width=0.95\textwidth, keepaspectratio ]{#1.\graphicxextension}
  \end{center}
\end{frame}

}
\else
\providecommand{\wpicframe}[4]
{
\begin{wrapfigure}{O}{#3\textwidth}
\centering \includegraphics[width=#3\textwidth]{#1.\graphicxextension}
\caption{#2 \label{#4}}
\end{wrapfigure}
}
\fi

%-----------------------------------------------------
%	FRAME										\texframe 
%---------------------------------------------------
\ifdefined\PRZMODE

\providecommand{\texframe}[3]
{
\begin{frame}
\frametitle<presentation>{#1}
\framesubtitle<presentation>{#2}
#3
\end{frame}
}
\else
\providecommand{\texframe}[3]
{
#3
}

\fi

%-----------------------------------------------------
%											\uref 
%---------------------------------------------------
%\newcommand{\uref}[2]{(Voir~#1~\cref{#2} page~\pageref{#2})}

\newcommand{\uref}[2]{(Voir~#1~\ref{#2} page~\pageref{#2})}
%-----------------------------------------------------
%											\ubg
%---------------------------------------------------
\newcommand{\ubg}[1]{
\textbf{\g{\uppercase{#1}}}
}
%-----------------------------------------------------
%											\updfimage 
%---------------------------------------------------

\newcommand{\updfimage}[3]
{
\begin{wrapfigure}{R}{#3\textwidth}
  \begin{center}
	 \resizebox{#3\textwidth}{!}{\includegraphics{\includepicturesubpath{#1.pdf}}}
  \end{center}
\caption{#2}
\end{wrapfigure}
}


%************************PACKAGE***

%\usepackage{tikz}% Required for drawing custom shapes
\usetikzlibrary{shadows}


%-----------------------------------------------------
%											 		\rem 
% insertion d'une remarque avec R
%---------------------------------------------------
\newcommand\rem[1]{%
   \marginpar{%
   \tikzpicture[baseline={(title.base)}]
      \node[inner sep=5pt,text width=4cm,drop shadow={shadow yshift=-5pt,shadow xshift=5pt,uMainColor},fill=white] (box) {\vskip5pt \nointerlineskip #1};
      \node[right=10pt,inner sep=0pt,outer sep=10pt] (title) at (box.north west) {\bfseries\color{uMainColor}Remarque};
      \draw[draw=uMainColor,very thick](title.west)--(box.north west)--(box.south west)--(box.south east)--(box.north east)--(title.east);
      \fill[uMainColor]([yshift=-10pt]box.north west)--+(-5pt,-5pt)--+(0pt,-10pt);
   \endtikzpicture}%
}
%-----------------------------------------------------
%											 \utikzimage 
% insertion d'une image sur fichier tkz.tex
%---------------------------------------------------

\newcommand{\utikzimage}[3]
{
\begin{wrapfigure}{R}{#3\textwidth}
  \begin{center}
	 \resizebox{#3\textwidth}{!}{\input{../#1.tkz.tex}}
  \end{center}
\caption{#2}
\end{wrapfigure}
}

%-----------------------------------------------------
%														\link 
% Url et footnote du lien
%-----------------------------------------------------

% TT every character and hyphenate after it
\def\hyphenateAndTtWholeString #1{\xHyphenate#1$\wholeString\unskip}

\def\xHyphenate#1#2\wholeString {\if#1$%
    \else\transform{#1}%
    \takeTheRest#2\ofTheString\fi}

\def\takeTheRest#1\ofTheString\fi
{\fi \xHyphenate#1\wholeString}

\def\transform#1{\url{#1}\hskip 0pt plus 1pt}

% Define the \urlx command which works like \url, but with line brakes
\def\urlx #1{\href{#1}{\hyphenateAndTtWholeString{#1}}}
  

 %-----------------------------------------------------
%														\link 
%-----------------------------------------------------  
  
  \newcommand{\ulink}[2]
  {\href{#1}
  {#2~\raisebox{-0.2ex}{\faExternalLink}\footnote{\url{#1}}}\xspace}
  
  \newcommand{\wikipedia}[1]
  { \href{https://fr.wikipedia.org/wiki/#1}{\textcolor{grey}{#1}}\raisebox{0.5ex}{\tiny{\textcolor{uMainColor}{\faWikipediaW}}}}
  


 %-----------------------------------------------------
%														\utodo 
%-----------------------------------------------------

  \newcommand{\utodo}{\begin{center}{\noindent \textcolor{cyan}{\faGears~{En cours de rédaction, DRAFT non publiable} \faGears}}\end{center}}
  
    \newcommand{\utocomplete}{\begin{center}{\noindent \textcolor{blue}{\faClipboard~{à compléter, les éléments ne donnent qu'une vue trop réduite ou trop parcellaire du sujet}}}\end{center}}
  

 





%-------------------------------------------------------
\usepackage{\utemplatepath/LessonModel/Beamer/\ubeamermodel.model/edxbeamer}
%-------------------------------------------------------


% a charger après hyperref, pour avoir les guillemets sous beamer
\usepackage{ae}
\usepackage{aecompl}


\usepackage{environ}
\newwrite\audioexport



%\renewcommand{\familydefault}{pag}

%-------------------------------------------------------
% Optional: a subtitle to be displayed on the title slide

\title[\ushorttitle]{\lessonID~-~\utitle}
\subtitle{\tiny\uCourseLongName}

%-------------------------------------------------------
% The author(s) of the presentation:
%  - again first a short version to be displayed at the bottom;
%  - next the full list of authors, which may include contact information;

\author[\uauthorwriter]{ \LARGE{\uauthor} \\ \medskip {\small \url{\umaila} \and \url{\umailb}\\ {\small \url{\usitea}}}}

%-------------------------------------------------------
% The institute:
%  - to start the name of the university as displayed on the top of each slide
%    this can be adjusted such that you can also create a Dutch version
%  - next the institute information as displayed on the title slide

\institute[\uinstituteshort]{{\uinstitute \\ \uchaire}}


%-------------------------------------------------------
% Add a date and possibly the name of the event to the slides
%  - again first a short version to be shown at the bottom of each slide
%  - second the full date and event name for the title slide
\date[\uversion]{\tiny{Publication \INFODistrib ~du \\ \DTMnow}}


\setbeamertemplate{navigation symbols}{%
\insertslidenavigationsymbol
\insertframenavigationsymbol
\insertsubsectionnavigationsymbol
\insertsectionnavigationsymbol
\insertdocnavigationsymbol
\insertbackfindforwardnavigationsymbol
}

%\mode<all>{
%%-------------------------------------------------------------
%               FR CYBERDEF SECOPS COURSE
%
%                 DEFINITION DES Sources/EDU/acronymes

%                      $File :  Sources/EDU/acronymes.tex
%                             2020 eduf@ction
%-------------------------------------------------------------
% Require : glossaries package
%-------------------------------------------------------------
%\newacronym{XXX}{XXX}{XXX}
%-----------------------------------------------------------

%  Organismes

\newacronym{aCNIL}{CNIL}{Commission National Informatique et Liberté}
\newacronym{aOTAN}{OTAN}{Organisation du Traité de l'Atlantique Nord}
\newacronym{aANSSI}{ANSSI}{Agence Nationale de la Sécurité des systèmes d'information}
\newacronym{aCERT}{CERT}{Computer Emergency Response Team}
\newacronym{aAFNOR}{AFNOR}{Association Française de Normalisation}
\newacronym{aSIEM}{SIEM}{System Incident and Event Management}
\newacronym{aSIM}{SIM}{System Incident Management}
\newacronym{aSEM}{SEM}{System Event Management}
\newacronym{aIRP}{IRP}{Incident Response Plan, Plan de Réponse à Incident}
\newacronym{aISO}{ISO}{International Organization for Standardization}
\newacronym{aDICT}{DICT}{Disponibilité, Intégrité, Confidentialité, Traçabilité}
\newacronym{aPDCA}{PDCA}{Plan, Do, Check, Act (Roue de Deming)}
\newacronym{aGDPR}{GDPR}{General Data Protection Regulation ou RGPD Règlement général sur la protection des données}

%-----------------------------------------------------------

% ROLES ET METIERS

\newacronym{aOS}{OS}{Officier de sécurité}
\newacronym{aRSSI}{RSSI}{Responsable de la sécurité des systèmes d'information}
\newacronym{aDPO}{DPO}{Data Protection Officier}
\newacronym{aCIL}{CIL}{ Correspondant Informatique et Liberté}
\newacronym{aSSI}{SSI}{Sécurité des Systèmes d’information}
\newacronym{aGRC}{GRC}{Gouvernance Risques et Conformité}
\newacronym{aDTF}{DTF}{Defend The Flag}
\newacronym{aCTF}{CTF}{Capture The Flag}
\newacronym{aOIV}{OIV}{Opérateur d'Infrastructures Vitales}
\newacronym{aOSE}{OSE}{Opérateur de Services Essentiels}
\newacronym{aLPM}{LPM}{Loi de Programmation Militaire}
\newacronym{aEBIOS}{EBIOS}{Expression des Besoins et Identification des Objectifs de Sécurité}
\newacronym{aDLL}{DLL}{Dynamic Loads Library}
\newacronym{aACID}{ACID}{Authentification, Confidentialité, Intégrité, Disponibilité}

%-------------------------------------------------------
% Sources/EDU/acronymes DATABASE XLS Generation (cf. RESSOURCES DIR)
%-------------------------------------------------------
% DO NOT MODIFY DIRECTLY, USE XLS File (BEGIN)

\newacronym{a2FA}{2FA}{\textit{Two-Factor Authentication} : Authentification à deux facteurs}
\newacronym{aAAA}{AAA}{\textit{Authentication, Authorization, Accounting} : Authentification, autorisation, comptabilisation}
\newacronym{aAC}{AC}{\textit{Access Control} : Accès conditionnel}
\newacronym{aAES}{AES}{\textit{Advanced Encryption Standard} : Algorithme de chiffrement avancé}
\newacronym{aAPT}{APT}{\textit{Advanced Persistent Threat Group} : Groupe de cyber-espionnage avancé}
\newacronym{aASLR}{ASLR}{\textit{Address Space Layout Randomization} : Randomisation de l'espace d'adressage}
\newacronym{aAV}{AV}{\textit{Antivirus Software} : Logiciel antivirus}
\newacronym{aBOT}{BOT}{\textit{Robot (Malicious)} : Robot (malveillant)}
\newacronym{aCC}{C\&C}{\textit{Command And Control} : Commande et contrôle}
\newacronym{aCSRF}{CSRF}{\textit{Cross-Site Request Forgery Attack} : Attaque de confiance de site à site}
\newacronym{aDDOS}{DDOS}{\textit{Distributed Denial Of Service Attack} : Attaque de déni de service distribué}
\newacronym{aDES}{DES}{\textit{Data Encryption Standard} : Algorithme de chiffrement à clé de 64 bits}
\newacronym{aDLP}{DLP}{\textit{Data Loss Prevention} : Gestion de la protection des données sensibles}
\newacronym{aEDR}{EDR}{\textit{Endpoint Detection And Response} : Détection et réponse aux comportements anormaux d'un terminal}
\newacronym{aEK}{EK}{\textit{Exploit Kit} : Kit d'exploit}
\newacronym{aFTP}{FTP}{\textit{File Transfer Protocol} : Protocole de transfert de fichier}
\newacronym{aHTTPS}{HTTPS}{\textit{Protocol For Securing Communications On The Internet} : Protocole de sécurisation HTTP}
\newacronym{aIAM}{IAM}{\textit{Identity And Access Management} : Identité et Gestion des Accés}
\newacronym{aIDS}{IDS}{\textit{Intrusion Detection System} : Système de détection d'intrusion}
\newacronym{aIIS}{IIS}{\textit{Internet Information Services Server}}
\newacronym{aIPS}{IPS}{\textit{Intrusion Prevention System} : Système de prévention d'incident}
\newacronym{aISA}{ISA}{\textit{Internet Security And Acceleration Server} : Serveur de sécurité et d'accélération Internet (Microsoft)}
\newacronym{aLDAP}{LDAP}{\textit{Lightweight Directory Access Protocol} : Protocole de gestion d'annuaire léger}
\newacronym{aMD5}{MD5}{\textit{Message-Digest Algorithm 5}}
\newacronym{aMFA}{MFA}{\textit{Multi-Factor Authentication} : Authentification multi-facteurs}
\newacronym{aNAC}{NAC}{\textit{Network Access Control}}
\newacronym{aOAuth}{OAuth}{\textit{Open Authorization Protocol}}
\newacronym{aPAM}{PAM}{\textit{Privilege Access Management}}
\newacronym{aPKI}{PKI}{\textit{Public Key Infrastructure} : Infrastructure de clefs publiques}
\newacronym{aRSA}{RSA}{\textit{Rivest-Shamir-Adleman Encryption Algorithm}}
\newacronym{aSAML}{SAML}{\textit{Security Assertion Markup Language}}
\newacronym{aSMS}{SMS}{\textit{Short Message Service}}
\newacronym{aSQLI}{SQLI}{\textit{Sql Injection} : Injection de script de commande SQL}
\newacronym{aSSH}{SSH}{\textit{Secure Shell Protocol}}
\newacronym{aSSL}{SSL}{\textit{Protocol For Securing Communications On The Internet} : Protocole de sécurisation des communications sur Internet}
\newacronym{aTKIP}{TKIP}{\textit{Temporal Key Integrity Protocol}}
\newacronym{aTLS}{TLS}{\textit{Transport Layer Security}}
\newacronym{aURL}{URL}{\textit{Uniform Resource Locator}}
\newacronym{aVPN}{VPN}{\textit{Virtual Private Network} : Réseau privé virtuel}
\newacronym{aWAF}{WAF}{\textit{Web Application Firewall}}
\newacronym{aWPA}{WPA}{\textit{Wifi Protected Access}}
\newacronym{aXSS}{XSS}{\textit{Cross-Site Scripting Attack}}

% DO NOT MODIFY (END)
%-------------------------------------------------------




%%-------------------------------------------------------------
%          FR CYBERDEF SECOPS COURSE
%
%                        GLOSSAIRE
%
%                 $File :  Sources/EDU/glossaire.tex
%                    2020 eduf@ction
%-------------------------------------------------------------
% Require : glossaries package
%-------------------------------------------------------------

%\newglossaryentry{xxx}
%{
%name={xxx},
%description={xxx}
%}

\newglossaryentry{gLPM}
{
name={Loi de programmation militaire},
description={Une loi de programmation militaire (LPM), est une loi visant à établir une programmation pluriannuelle des dépenses et donc du budget que l'État français consacre à ses forces armées. Les lois de programmation militaire ont des durées d'application de quatre, cinq ou six ans. La cyberdéfense fait partie des axes importants des lois de programmation}
}

\newglossaryentry{gCTF}
{
name={Capture The Flag},
description={Challenge consistant à exploiter ses connaissance et les vulnérabilités affectant des logiciels  et systèmes numériques de manière à s’introduire sur des réseaux , des ordinateurs, pour récupérer (capturer) des preuves de l’intrusion réussie dénommées G\g{des drapeaux (\textit{Flags} en anglais)}.}
}

\newglossaryentry{gDTF}
{
name={Defend The Flag},
description={Challenge consistant à détecter, résister à des attaquants qui tentent de s'introduire ou sont introduits dans des systèmes numériques de manière à défendre le bien considéré comme le plus sensible convoité par l'attaquant dénommé G\g{drapeau (\textit{Flags} en anglais)}.}
}


\newglossaryentry{gROOTKIT}
{
name={Rootkit},
description={ou simplement « kit » (est un « outil de dissimulation d'activité » ou  « code furtif » qui est une sorte de trousse à outils pour un pirate informatique). Il est constitué  de techniques mises en œuvre par un ou plusieurs logiciels, dont le but est d'obtenir et de pérenniser un accès ceci avec en déployant des techniques de furtivité et de dissimulation.}
}

%-------------------------------------------------------
% GLS DATABASE XLS Generation (cf. RESSOURCES DIR)
%-------------------------------------------------------
% DO NOT MODIFY DIRECTLY, USE XLS File (BEGIN)

\newglossaryentry{gAffiliate}{name={\textit{Affiliates}},description={(Affiliés): Cybercriminels qui mènent des cyberattaques, parfois sophistiquées, par le biais de rançongiciels qui sont mis à disposition par des développeurs ou concepteurs. Les gains générés par les affiliés sont partagés avec les développeurs via un système de commissions.}}

% DO NOT MODIFY (END)
%-------------------------------------------------------


%}

%++++++++++++++++++++++++++++++++
\begin{document}



  %-------------------------------------------------------
  % AudioNote Open File
  %-------------------------------------------------------
\mode<all>{\immediate\openout\audioexport=\jobname.txt}

%================================
\begin{frame}[plain]
 					\begin{figure}
  					 \centering
  					 \begin{columns}[c]
  \begin{column}{0.4\textwidth}
  
    \begin{figure}
        \begin{flushleft}
  			\includegraphics[width=0.9\textwidth]{\utemplatepath/commons.inc/CommonsPictures/\utemplatetheme.theme/maincourselogo.pdf}
  \end{flushleft}
  \end{figure}
      \end{column}
  
  \begin{column}{0.6\textwidth}
    \begin{figure}
        \begin{flushright}
 		    \includegraphics[width=0.7\textwidth]{\utemplatepath/commons.inc/CommonsPictures/\utemplatetheme.theme/maininstitutelogo.pdf}
    \end{flushright}
  \end{figure}
  \end{column}
\end{columns}
   							  
   							  \end{figure}
							
				  				\titlepage
\end{frame}
%================================



%================================
\begin{frame}
	\frametitle{Abstract}
	\begin{techworkin} [righthand width=.8\textwidth]{\faCogs~\small {Hashtags} : \ukeywords}\uabstract
	\end{techworkin}
	\end{frame}
%================================

%================================
\begin{frame}
 \frametitle{Sommaire}
      \begin{multicols}{2}
         \tableofcontents[hideallsubsections]
     \end{multicols}
 \end{frame}
 %================================

%-------------------------------------------------------
%  BEAMER MAIN TEX configured with ubody
%-------------------------------------------------------
\mode<all>{\ubody}

%-------------------------------------------------------
% Last Frame
%-------------------------------------------------------
\begin{frame}
 \frametitle{des questions ?}
  \framesubtitle{contacter \umaila}
  					\begin{figure}
  					 \centering
   						 \includegraphics[width=0.2\textwidth]{\utemplatepath/commons.inc/CommonsPictures/\utemplatetheme.theme/mainlogo.pdf}\\
				  \end{figure}
				 	 \centering \textit{Tous les documents publiés dans le cadre de ce cours sont perfectibles, \\ne pas hésiter à m'envoyer vos remarques !}
 \end{frame}

%-------------------------------------------------------
% Contributions
%-------------------------------------------------------
\begin{frame}
 \frametitle{Contributions}
\begin{columns}
		\begin{column}{0.1\textwidth}

 					\begin{figure}
  						  \centering
   						 \includegraphics[width=0.9\textwidth]{../Latex/Sources/EDU/SRC1/Contribs/contribute/ContribPictures/github.pdf}\\
   						 	 \includegraphics[width=0.9\textwidth]{../Latex/Sources/EDU/SRC1/Contribs/contribute/ContribPictures/overleaf.pdf}
				  \end{figure}
		\end{column}
		\begin{column}{0.9\textwidth}
							%-------------------------------------------------------------
%               FR CYBERDEF SECOPS COURSE
%                            Contributions Rules
%                              2020 eduf@ction
%-------------------------------------------------------------

% \section{Contributions}

% \subsection{Comment contribuer}

Les notes et les présentations sont réalisées sous \hologo{LaTeX}. 

Vous pouvez contribuer au projet du cours CYBERDEF101. Les contributions peuvent se faire sous deux formes :

\begin{itemize}
  \item Corriger, amender, améliorer les notes publiées. A chaque session des modifications et évolutions sont apportées pour tenir compte des corrections de fond et de formes. 
  \item Ajouter, compléter, modifier des parties de notes sur la base de votre lecture du cours et de vos expertises dans chacun des domaines évoqués.
\end{itemize}

 Les fichiers sources sont publiés sur GITHUB dans l'espace : \ulink{https://github.com/edufaction/CYBERDEF101}{(edufaction/CYBERDEF101)}.
 
		\end{column}
\end{columns}
 \end{frame}
 
 %-------------------------------------------------------
% Link to download file and copyright
%-------------------------------------------------------
\begin{frame}
\frametitle{Mises à jour régulières}
\framesubtitle{Eduf@ction~\umaila}
%\INFOGitHub
{\begin{center} 
\setstretch{2}
 {Vérifiez la disponibilité d'une version plus récente de} \\
 {\ulink{\GITfilename}{\textbf{\udocname.pdf }sur GITHUB CYBERDEF} }   \\ 
{{\huge\ccbyncndeu}}  \\  
{\the\year~eduf@ction - Publication en Creative Common BY-NC-ND }    \\  % Copyright notice
{\safeqrcode[padding]{\GITfilename}}  
\end{center} }
\end{frame}


  %-------------------------------------------------------
  % AudioNote Open File
  %-------------------------------------------------------
  \mode<all>{\immediate\closeout\audioexport}


\end{document}




}
\else

\ifdefined\EDTMODE 
	\providecommand{\umainload}{%-------------------------------------------------------------
%               FR CYBERDEF SECOPS COURSE
%                          ARTICLE MAIN FILE
%                             2020 eduf@ction
%-------------------------------------------------------------

\documentclass[10pt,fleqn,twoside]{\utemplatepath/LessonModel/Article/edxstyle} 
%edxstyle
\usepackage{beamerarticle}

\usepackage{ae}
\usepackage[french]{babel}
\usepackage[utf8]{inputenc}
\usepackage[T1]{fontenc}
\usepackage{lmodern}




\usepackage{titletoc} % Required for manipulating the table of contents

\input{\utemplatepath/commons.inc/CommonsPictures/\utemplatetheme.theme/colors.inc}

%-------------------------------------------------------------
%               FR CYBERDEF SECOPS COURSE
%                                  EDX PACKAGES
%.    Commons packages for BOOK, NOTES (.doc) 
%                            and BEAMER (.prz)
%                              2020 eduf@ction
%-------------------------------------------------------------

%\usepackage[hyphens,spaces,obeyspaces]{url}
%\usepackage[hyphens,spaces,obeyspaces]{url}



\PassOptionsToPackage{hyphens,spaces,obeyspaces}{url}

\usepackage[utf8]{inputenc}
\usepackage[french]{babel}
\usepackage[T1]{fontenc}
%\usepackage{markdown}



\usepackage{ccicons}

\usepackage{tikz} % Required for drawing custom shapes
\usetikzlibrary{tikzmark} 
\usetikzlibrary{mindmap,trees, backgrounds}


\usepackage{verbatim}
\usepackage{wrapfig}
\usepackage{xpatch}

\usepackage{hologo}

\usepackage[figure]{totalcount}

\usepackage{fontawesome}

\usepackage[skins]{tcolorbox} % don't forget SKINS option

\usepackage{listings}
\usepackage{upquote}

\lstset{
upquote=true,columns=flexible, basicstyle=\ttfamily,
language=HTML, 
frameround=tttt,
commentstyle=\color{gray},
identifierstyle=\color{blue},
keywordstyle=\color{uMainColor}\bfseries,
xleftmargin=2em,
xrightmargin=2em,
aboveskip=\topsep,
belowskip=\topsep, 
frame=single,
rulecolor=\color{uMainColor},
backgroundcolor=\color{uMainColor!5},
breaklines,
breakindent=1.5em,
showspaces=false,
showstringspaces=false,
showtabs=false,
xleftmargin=2em,
xrightmargin=2em,
aboveskip=1em,
belowskip=1em,
}



\usepackage{graphicx} % Required for including pictures
%\graphicspath{{Pictures/}} % Specifies the directory where pictures are stored

\usepackage{booktabs} % Required for nicer horizontal rules in tables

\usepackage{setspace}
\setstretch{1,2}

\usepackage{csquotes}

\usepackage{fancyhdr}

\usepackage{makeidx} 

\usepackage{xspace} 

\usepackage{floatflt} 

\usepackage{hyperref}
\usepackage{cleveref}

%pdfdisplaydoctitle={\utitle}, Attention le titre ne doit pas contenir de saut de ligne 

\hypersetup{
pdftitle={\uCoursetittle},
pdfauthor={\uauthor},
hidelinks,
backref=true,
pagebackref=true,
hyperindex=true,
pdfcenterwindow = true,
colorlinks=true,
breaklinks=true,
bookmarks=true,
allcolors=uMainColor,
pdfsubject = {\udescription},
pdfcreator = {MACTEX on TEXPAD},
bookmarksopen=true}


% urlcolor=cnam,
% runcolor=uMainColor,
%\pdfstringdefDisableCommands{%
%  \def\\{}%
%  \def\texttt#1{<#1>}%
%}

\usepackage[style=numeric,citestyle=numeric,sorting=nyt,sortcites=true,autopunct=true,hyperref=true,abbreviate=false,backref=true,backend=biber]{biblatex}

%\usepackage{cleveref} old position

\usepackage{multicol}

%--------------------------------------------------------
% pour la génération du QRCODE en page de garde d'article
\usepackage{qrcode} 

%--------------------------------------------------------
% POur la date d'impression en premiere page d'article
\usepackage[useregional]{datetime2}

%--------------------------------------------------------
% Gestion des glossaires et Sources/EDU/acronymes
% Un fichier acroymes.tex centralisé
\usepackage[acronym]{glossaries}

%\makeglossaries
\makenoidxglossaries 

\loadglsentries{Latex/Sources/commons/acronymes.tex}
\loadglsentries{Latex/Sources/commons/glossaire.tex}


%\renewcommand{\bibname}{R\'ef\'erences}

%--------------------------------------------------------
% Pour exclure des "environnement", ici techworkbox par exemple, travaux pour étudiants, non affiché dans le BOOK.
\usepackage{comment}


%--------------------------------------------------------
% pour filigramme sur les pages
\usepackage{eso-pic}

%--------------------------------------------------------
% for \gappto{\UrlBreaks}{\UrlOrds}
% rather than 
%\makeatletter
% \g@addto@macro{\UrlBreaks}{\UrlOrds}
% \makeatother

\usepackage{etoolbox}

\usepackage{breakurl}


\newcommand{\ucite}[1]{\cite{#1}}
%\newcommand{\ucite}[1]{REF BIB #1}
\newcommand{\ugls}[1]{\gls{#1}}
%\newcommand{\ugls}[1]{GLOSSAIRE (#1)}
%\usepackage[Rejne]{fncychap}

%\usepackage{pythontex} % need shell escape

%
%\makeatletter
%\DeclareRobustCommand{\afilename}[1]{%
% \begingroup
%  % \lstname seems to change hyphens into \textendash
%  \def\textendash{-}%
%  \afilename@parse{#1}%
%  \edef\afilename@base{\detokenize\expandafter{\afilename@base}}%
%  \texttt{\afilename@base.\afilename@ext}%
% \endgroup
%}
%\makeatother

%-------------------------------------------------------------
%               FR CYBERDEF SECOPS COURSE
%                                  EDX PACKAGES
%.    Commons COMMANDS for BOOK, NOTES (.doc) 
%                            and BEAMER (.prz)
%                              2020 eduf@ction
%-------------------------------------------------------------

\newcommand{\ucite}[1]{\cite{#1}}
%\newcommand{\ucite}[1]{REF BIB #1}
\newcommand{\ugls}[1]{\gls{#1}}
%\newcommand{\ugls}[1]{GLOSSAIRE (#1)}
\newcommand{\uac}[1]{\acrshort{#1}}



% \newcommand{\ExtractBeforeLastSlash}[1]{\directlua{extract_before_last_slash("#1")}}
% \newcommand{\ExtractAfterLastSlash}[1]{\directlua{extract_after_last_slash("#1")}}
% \newcommand{\RemoveExtension}[1]{\directlua{remove_extension("#1")}}
% \newcommand{\ExtractExtension}[1]{\directlua{extract_extension("#1")}}


\newcommand{\includepicturesubpath}[1]
{#1}

%{\ExtractBeforeLastSlash{#1}/\upicturethemepath\RemoveExtension{\ExtractAfterLastSlash{#1}}.\ExtractExtension{#1}}
    

%---------------------------------------------------
%										 	AUDIO NOTE 
%---------------------------------------------------

\newcommand\audionote[2]{%
  \note{#2}
  \immediate\write\audioexport{==AudioSlide=============}
  \immediate\write\audioexport{#1}
  \immediate\write\audioexport{\insertframenumber}
  \immediate\write\audioexport{---------------------------------}
  \immediate\write\audioexport{#2}
  \immediate\write\audioexport{---------------------------------}
}%

%---------------------------------------------------
%										 	TEXT STYLE SHORTCUT 
%---------------------------------------------------

\newcommand{\tcb}[1]{\textbf{\textcolor{uMainColor}{#1}}}

\newcommand{\voir}[1]{(Schéma~\ref{#1},~page~\pageref{#1}}

\newcommand{\graphicxextension}{pdf}

\newcommand\myemptypage{
   
    \thispagestyle{empty}
    \addtocounter{page}{-1}
    \newpage
  \null
    %CECI EST UNE PAGE VIDE
       \vfill
    }


\newcommand{\uNA}[2]
{
\newacronym{a#1}{#1}{#2}
}

\newcommand{\coauthor}[4]
{
\noindent (#2) \textbf{#1} - #4 : \textit{#3}
}


\newcommand\UKword[1]{\emph{#1}\xspace}
\newcommand\textUK[1]{\textit{\textbf{#1}}\xspace}
\newcommand\Ucenter[1]{\begin{center}{#1}\end{center}}

\newcommand\head[1]{\textbf{#1}}
\newcommand\tb[1]{\textbf{#1}}
\newcommand{\g}[1]{\og #1 \fg{}}
\newcommand{\uload}[1] {\input{../Latex/#1}}
\newcommand{\ucurcolor}{black}

\newcommand{\mak}[1]{\faicon{\aArrowCircleORight} #1}


%-----------------------------------------------------
%										 	\INFOGitHub 
%---------------------------------------------------


\newcommand{\GITfilename}{https://github.com/edufaction/CYBERDEF101/raw/main/Release/Builder/\jobname.pdf}

%\newcommand{\safeqrcode}[2][]{%
%\qrcode[#1]{\detokenize{#2}} }

\newcommand{\safeqrcode}[2][]{\qrcode[#1]{#2}} 

\newcommand{\INFOGitHub}
{
\begin{center} 
\setstretch{1.5}
 {Vérifiez la disponibilité d'une version plus récente de} \\
 {\ulink{\GITfilename}{\textbf{\jobname.pdf }sur GITHUB CYBERDEF} }   \\ 
{{\huge\ccbyncndeu}}  \\  
{Publication en \href{https://creativecommons.org/licenses/by-nc-nd/2.0/fr/}{Creative Common BY-NC-ND} by eduf@ction}    \\  % Copyright notice
{\safeqrcode[height=4 cm,padding]{\GITfilename}}  
\end{center} 
}

\newcommand{\INFOGitHubS}
{
\noindent{\safeqrcode[height=3 cm]{\GITfilename}} \\
\noindent {\tiny Téléchargez une version à jour} 
}



%-----------------------------------------------------
%										 	\upspicture 
%---------------------------------------------------
\newcommand{\upspicture}[3]
{
\renewcommand{\ucurcolor}{#3}
\resizebox{#2\textwidth}{!}{\input{../#1}}
}


%-----------------------------------------------------
%										 	\uindex 
%---------------------------------------------------


\newcommand{\uindex}[1]
{
#1\index{#1}\xspace
}

\newcommand{\ulindex}[1]{\index{#1}}

%-----------------------------------------------------
%										 	\includer 
%---------------------------------------------------

\newcommand{\includer}[1]
{
\input{../Latex/Sources/EDU/SRC1/Chapters/#1}
}

%-----------------------------------------------------
%										 	\edxdico 
%---------------------------------------------------

\newcommand{\edxdico}[2] {\textbf{#1}:\textit{#2} \index{#1}}

%-----------------------------------------------------
%  command DEF : InTextx
%-----------------------------------------------------
\newcommand\InTexta{no text}
\newcommand\InTextb{no text}
\newcommand\InTextc{no text}
\newcommand\rnc{\renewcommand}


%-----------------------------------------------------
%										 	ENV warningbox
%---------------------------------------------------

\newtcolorbox{warningbox}[2][]
{
  colframe = uMainColor!25,
  colback  = black!05,
  coltitle = uMainColor!20!black,
  fonttitle = \bfseries, 
  title    = #2,
  #1,
}


\newtcolorbox{techworkbox}[2][]{sharp corners,colframe = uMainColor!25,colback  = black!05,coltitle = uMainColor!20!black,fonttitle = \bfseries, title={\textcolor{uMainColor}{\faFile~\small {Sujet de mémoire SECOPS}}} : #2,#1,}


\newtcolorbox{notebox2}[2][]
{
sharp corners, 
colback = uMainColor!5!white, 
colframe = uMainColor!75!black,
fonttitle = \bfseries, 
colbacktitle= uMainColor!85!black,
title=#2,#1
}

\newtcolorbox{notebox}[2][]{colback=uMainColor!5!white,
colframe=uMainColor!75!black,fonttitle=\bfseries,
colbacktitle=uMainColor!85!black,enhanced,
attach boxed title to top right={yshift=+1mm},
title=#2,#1}


\newtcolorbox{toolsboxin}[2][]{sharp corners, colback=uMainColor!0!white,
colframe= black!25,fonttitle=\bfseries,
colbacktitle=black!60,enhanced,
attach boxed title to top left={yshift=+1mm},
title=#2,#1}


\newtcolorbox{techworkin}[2][]{sharp corners, colback=uMainColor!5!white,
colframe= uMainColor!100, fonttitle=\bfseries,
colbacktitle=black!5, enhanced,
 coltitle=black!50,
attach boxed title to top right ={yshift=+1mm},
title=#2,#1}



%-----------------------------------------------------
%										 	ENV ToolsBox
%---------------------------------------------------


\newcommand\tooleditor{No EDITOR in Datas.tex}
\newcommand\toolurl{No URL INSERTED in Datas.tex}
\newcommand\toolclass{No EDITOR in Datas.tex}
\newcommand\toolname{No NAME in Datas.tex}
\newcommand\toolanalyst{No ANALYST in Datas.tex}

\newcommand{\toolsbox}[2]
{
\input{../Latex/Sources/EDU/SRC1/SecComponents/#2.#1/datas.tex} 
\begin{toolsboxin}[sidebyside,righthand width=.7\textwidth]{#1}
\includegraphics[width=0.9\textwidth]{../Latex/Sources/EDU/SRC1/SecComponents/#2.#1/icon.png} \\
\tcblower
\input{../Latex/Sources/EDU/SRC1/SecComponents/#2.#1/description.tex}
\end{toolsboxin}
\noindent \small {\faCogs~Classe : \uppercase{\textbf{#2}}, Site de référence : \textbf{\ulink{\toolurl}{#1}}}
\newline
\small {\faGlobe~Editeur : \textbf{\tooleditor}}
\small {\faEye~Analyste : \textbf{\toolanalyst}}
}

%\renewcommand\toolsclass{SIEM}
%\renewcommand\toolsname{RSA Netwitness}
%\renewcommand\toolseditor{RSA}
%\renewcommand\toolsurl{https://www.rsa.com/fr-fr/products/threat-detection-response}

%-----------------------------------------------------
%										 	ENV  TechWork
%---------------------------------------------------


\newcommand{\techwork}[2]
{
\begin{techworkin} [righthand width=.8\textwidth]{\faCogs~\small {Fiche TECHNO} : #1}#2
\end{techworkin}
}

%-----------------------------------------------------
%										 	\uchap
%---------------------------------------------------
\ifdefined\INTERTITLE
\providecommand{\uchap}[1]
	{
	\begin{LARGE}
	\textbf{\Ucenter{#1 \\ -oOo-}}	
	\end{LARGE}
	}
\else
	\providecommand{\uchap}[1]{}
\fi

%-----------------------------------------------------
%											\uexpand
%---------------------------------------------------
\newcommand {\uexpand}[1]{
\mode<all>\input{../Latex/Sources/EDU/SRC1/Chapters/#1}
}


%-----------------------------------------------------
%											\spicture 
% direct include graphic with Scrivener
%-----------------------------------------------------

\newcommand{\spicture}[3]{
\begin{figure}[!hbtp] %h !hbtp
  \begin{center}
	 \includegraphics[width=0.8\textwidth]{#1}
  \end{center}
\caption{#3 - (slbl-#2)\label{slbl-#2}}
\end{figure}
}

%\StrBefore{#2}{.pdf}

%-----------------------------------------------------
%											\upicture 
%---------------------------------------------------

\ifdefined\PRZMODE

\providecommand{\upicture}[4]{
\framesubtitle{#2}
%\begin{figure}[!h] %h
  \begin{center}
	 \includegraphics[height=#3\textheight, width=#3\textwidth, keepaspectratio ]{\includepicturesubpath{#1.pdf}}
  \end{center}
%  \end{figure}
}

\else

\providecommand{\upicture}[4]{
\begin{figure}[!hbtp] %h !hbtp
  \begin{center}
	 \includegraphics[width=#3\textwidth]{\includepicturesubpath{#1.pdf}}
  \end{center}
\caption{#2 \label{#4}}
\end{figure}
}
	
\fi


\ifdefined\PRZMODE

\providecommand{\uwpicture}[4]{
\framesubtitle{#2}
%\begin{figure}[!h] %h
  \begin{center}
	 \includegraphics[height=0.8\textheight, width=0.8\textwidth, keepaspectratio ]{\includepicturesubpath{#1.pdf}}
  \end{center}
%  \end{figure}
}
\else
\providecommand{\uwpicture}[4]
{
\begin{wrapfigure}{O}{#3\textwidth} % inside / Outside
  \begin{center}
	 \includegraphics[width=#3\textwidth]{\includepicturesubpath{#1.pdf}}
\end{center}
\caption{#2 \label{#4}}
\end{wrapfigure}
}
\fi

%-----------------------------------------------------
%	\picframe for BEAMER (Image included in Beamer Slide)
%---------------------------------------------------

\ifdefined\PRZMODE

\providecommand{\picframe}[4]
{
\begin{frame}
\frametitle<presentation>{#2}
  \begin{center}

   \begin{tikzpicture}[remember picture, overlay]
    \node at (current page.center) {\includegraphics[width=#3\paperwidth, height=#3\paperheight,keepaspectratio]{\includepicturesubpath{#1.pdf}}};
    \end{tikzpicture}

  \end{center}
\end{frame}

}
\else
\providecommand{\picframe}[4]
{
\begin{figure}[!hbtp] % !hbtp ou !h 
\centering 
%{\includegraphics[width=#3\paperwidth, height=#3\paperheight,keepaspectratio]{\includepicturesubpath{#1.pdf}}}
\includegraphics[width=#3\textwidth]{\includepicturesubpath{#1.pdf}}
\caption{#2 \label{#4}}
\end{figure}
}
\fi

%-----------------------------------------------------
%.  for BEAMER (Image included in Beamer Slide)
%	\picframeS Simplified with Name and Label from One Param
%---------------------------------------------------


\ifdefined\PRZMODE

\providecommand{\picframeS}[2]
{
\begin{frame}
\frametitle<presentation>{#2}
  \begin{center}
	 \includegraphics[height=0.95\textheight, width=0.95\textwidth, keepaspectratio ]{\picpath/img-#1.\graphicxextension}
  \end{center}
\end{frame}

}
\else
\providecommand{\picframeS}[2]
{
\begin{figure}[!hbtp] % !hbtp ou !h 
\centering \includegraphics[width=0.9\textwidth]{\picpath/img-#1.\graphicxextension}
\caption{(s) #2 \label{lbl:#1}}
\end{figure}
}

\fi


%{figure}[!h]
%{wrapfigure}{O}{#3\textwidth}
%-----------------------------------------------------
%	FRAME			\wpicframe 
%---------------------------------------------------

\ifdefined\PRZMODE

\providecommand{\wpicframe}[4]
{
\begin{frame}
\frametitle<presentation>{#2}
  \begin{center}
	 \includegraphics[height=0.95\textheight, width=0.95\textwidth, keepaspectratio ]{#1.\graphicxextension}
  \end{center}
\end{frame}

}
\else
\providecommand{\wpicframe}[4]
{
\begin{wrapfigure}{O}{#3\textwidth}
\centering \includegraphics[width=#3\textwidth]{#1.\graphicxextension}
\caption{#2 \label{#4}}
\end{wrapfigure}
}
\fi

%-----------------------------------------------------
%	FRAME										\texframe 
%---------------------------------------------------
\ifdefined\PRZMODE

\providecommand{\texframe}[3]
{
\begin{frame}
\frametitle<presentation>{#1}
\framesubtitle<presentation>{#2}
#3
\end{frame}
}
\else
\providecommand{\texframe}[3]
{
#3
}

\fi

%-----------------------------------------------------
%											\uref 
%---------------------------------------------------
%\newcommand{\uref}[2]{(Voir~#1~\cref{#2} page~\pageref{#2})}

\newcommand{\uref}[2]{(Voir~#1~\ref{#2} page~\pageref{#2})}
%-----------------------------------------------------
%											\ubg
%---------------------------------------------------
\newcommand{\ubg}[1]{
\textbf{\g{\uppercase{#1}}}
}
%-----------------------------------------------------
%											\updfimage 
%---------------------------------------------------

\newcommand{\updfimage}[3]
{
\begin{wrapfigure}{R}{#3\textwidth}
  \begin{center}
	 \resizebox{#3\textwidth}{!}{\includegraphics{\includepicturesubpath{#1.pdf}}}
  \end{center}
\caption{#2}
\end{wrapfigure}
}


%************************PACKAGE***

%\usepackage{tikz}% Required for drawing custom shapes
\usetikzlibrary{shadows}


%-----------------------------------------------------
%											 		\rem 
% insertion d'une remarque avec R
%---------------------------------------------------
\newcommand\rem[1]{%
   \marginpar{%
   \tikzpicture[baseline={(title.base)}]
      \node[inner sep=5pt,text width=4cm,drop shadow={shadow yshift=-5pt,shadow xshift=5pt,uMainColor},fill=white] (box) {\vskip5pt \nointerlineskip #1};
      \node[right=10pt,inner sep=0pt,outer sep=10pt] (title) at (box.north west) {\bfseries\color{uMainColor}Remarque};
      \draw[draw=uMainColor,very thick](title.west)--(box.north west)--(box.south west)--(box.south east)--(box.north east)--(title.east);
      \fill[uMainColor]([yshift=-10pt]box.north west)--+(-5pt,-5pt)--+(0pt,-10pt);
   \endtikzpicture}%
}
%-----------------------------------------------------
%											 \utikzimage 
% insertion d'une image sur fichier tkz.tex
%---------------------------------------------------

\newcommand{\utikzimage}[3]
{
\begin{wrapfigure}{R}{#3\textwidth}
  \begin{center}
	 \resizebox{#3\textwidth}{!}{\input{../#1.tkz.tex}}
  \end{center}
\caption{#2}
\end{wrapfigure}
}

%-----------------------------------------------------
%														\link 
% Url et footnote du lien
%-----------------------------------------------------

% TT every character and hyphenate after it
\def\hyphenateAndTtWholeString #1{\xHyphenate#1$\wholeString\unskip}

\def\xHyphenate#1#2\wholeString {\if#1$%
    \else\transform{#1}%
    \takeTheRest#2\ofTheString\fi}

\def\takeTheRest#1\ofTheString\fi
{\fi \xHyphenate#1\wholeString}

\def\transform#1{\url{#1}\hskip 0pt plus 1pt}

% Define the \urlx command which works like \url, but with line brakes
\def\urlx #1{\href{#1}{\hyphenateAndTtWholeString{#1}}}
  

 %-----------------------------------------------------
%														\link 
%-----------------------------------------------------  
  
  \newcommand{\ulink}[2]
  {\href{#1}
  {#2~\raisebox{-0.2ex}{\faExternalLink}\footnote{\url{#1}}}\xspace}
  
  \newcommand{\wikipedia}[1]
  { \href{https://fr.wikipedia.org/wiki/#1}{\textcolor{grey}{#1}}\raisebox{0.5ex}{\tiny{\textcolor{uMainColor}{\faWikipediaW}}}}
  


 %-----------------------------------------------------
%														\utodo 
%-----------------------------------------------------

  \newcommand{\utodo}{\begin{center}{\noindent \textcolor{cyan}{\faGears~{En cours de rédaction, DRAFT non publiable} \faGears}}\end{center}}
  
    \newcommand{\utocomplete}{\begin{center}{\noindent \textcolor{blue}{\faClipboard~{à compléter, les éléments ne donnent qu'une vue trop réduite ou trop parcellaire du sujet}}}\end{center}}
  

 




%-------------------------------------------------------------
%               FR CYBERDEF SECOPS COURSE
%
%                                    EDX STYLE
%
%                              2020 eduf@ction
%-------------------------------------------------------------
\usepackage{enumitem}

\usepackage{parskip}
\setlength{\parindent}{0em}
\setlength{\parskip}{5pt} % 1ex plus 0.5ex minus 0.2ex}

\def\uCoverColor{}

%---------------------------------------------------------------
%	MARGINS
%---------------------------------------------------------------

\usepackage{geometry}

\geometry{
	paper=a4paper, 
%	top=1cm, 
%	bottom=1cm, 
%	left=1cm, 
%	right=1cm, 
%	headheight=20pt, 
%	footskip=1.4cm, 
%	headsep=10pt, 
%	showframe, 
	showcrop 
}

%---------------------------------------------------------------
%	FOOTNOTE (format)
%---------------------------------------------------------------

\makeatletter
\long\def\@makefntextFB#1{%
    \ifx\thefootnote\ftnISsymbol
        \@makefntextORI{#1}%
    \else
        \rule\z@\footnotesep
        \setbox\@tempboxa\hbox{\@thefnmark}%
            \ifdim\wd\@tempboxa>\z@
                \kern2em\llap{\@thefnmark.\kern0.5em}%
            \fi
        \hangindent2em\hangafter\@ne#1
    \fi}
\makeatother

%-----------------------------------------------------
%	COLOR
%-----------------------------------------------------
%\usepackage{xcolor} % Required for specifying colors by name
%\definecolor{uMainColor}{RGB}{160,0,0} 
\colorlet{edxcolorcover}{uCoverColor}
%\definecolor{edxcolorcover}{RGB}{160,0,0} 

%\definecolor{cnam}{RGB}{128, 0, 32} 
                 

  \definecolor{comments}{rgb}{0.7,0,0}    % rouge foncé
  \definecolor{link}{rgb}{0,0.4,0.6}      % ~RoyalBlue de dvips
  \definecolor{url}{rgb}{0.6,0,0}         % rouge-brun
  \definecolor{citation}{rgb}{0,0.5,0}    % vert foncé
  \definecolor{ULlinkcolor}{rgb}{0,0,0.3} % de ulthese.cls
  \definecolor{rouge}{rgb}{0.85,0,0.07}   % rouge bandeau identitaire
  \definecolor{or}{rgb}{1,0.8,0}          % or bandeau identitaire

%-----------------------------------------------------
%	FONTS
%-----------------------------------------------------

%%\usepackage{avant} % Use the Avantgarde font for headings
%\usepackage{times} % Use the Times font for headings
%\usepackage{mathptmx} % Use the Adobe Times Rgman as the default text font together with math symbols from the Sym­bol, Chancery and Com­puter Modern fonts

%%\usepackage{microtype} % Slightly tweak font spacing for aesthetics
%%\usepackage[utf8]{inputenc} % Required for including letters with accents
%\usepackage[T1]{fontenc} % Use 8-bit encoding that has 256 glyphs

% Using AVANT GARDE Family

%\renewcommand{\familydefault}{\sfdefault}

%\fontfamily{pag}\selectfont
%\renewcommand{\familydefault}{pag} % sfdefault pag lmss

%bch         Charter
%lmr         Latin Modern Roman
%lmss        Latin Modern Sans Serif
%lmssq       Latin Modern Sans Serif extended
%lmtt        Latin Modern Typewriter
%lmvtt       Latin Modern Typewriter proportional
%pag         Avant Garde
%pbk         Bookman
%pcr         Courier
%phv         Helvetica
%pnc         New Century Schoolbook
%ppl         Palatino
%ptm         Times
%put         Utopia
 
%\renewcommand{\rmdefault}{pag} % text 
%\renewcommand{\sfdefault}{pag} % titre


%\renewcommand{\rmdefault}{pag} % text 
%\renewcommand{\sfdefault}{phv} % titre


%-----------------------------------------------------
%	BIBLIOGRAPHY AND INDEX
%-----------------------------------------------------


\usepackage{calc} % For simpler calculation - used for spacing the index letter headings correctly

\usepackage{imakeidx}
%\makeindex % Tells LaTeX to create the files required for indexing

\makeindex[columns=2,intoc=true,options={-s \upath/indexstyle.ist}]

%-----------------------------------------------------
%	HEADERS AND FOOTERS
%-----------------------------------------------------


\usepackage{fancyhdr} % Required for header and footer configuration
\pagestyle{fancy} % Enable the custom headers and footers

\ifcsname chapter \endcsname   % verify if command name exists in the CLASS (Book, report, article ...)
% BEGIN IF
\renewcommand{\chaptermark}[1]{\markboth{\sffamily\normalsize\bfseries\chaptername\ \thechapter.\ #1}{}} % Styling for the current chapter in the header
  \else
%nothing
\fi
% END IF

\renewcommand{\sectionmark}[1]{\markright{\sffamily\normalsize\thesection\hspace{5pt}#1}{}} % Styling for the current section in the header

\fancyhf{} % Clear default headers and footers
\fancyhead[LE,RO]{\sffamily\normalsize\thepage} % Styling for the page number in the header
\fancyhead[LO]{\sffamily\normalsize\rightmark} % Print the nearest section name on the left side of odd pages
\fancyhead[RE]{\sffamily\normalsize\leftmark} % Print the current chapter name on the right side of even pages
\fancyfoot[RE,LO]{\sffamily\normalsize\uCoursetittle} % Uncomment to include a footer
\fancyfoot[LE,RO]{\hyperlink{toc}{\sffamily\normalsize\uMinilogo}} % Uncomment to include a footer
% \hyperlink{toc}
\renewcommand{\headrulewidth}{1.0pt} % Thickness of the rule under the header

\fancypagestyle{plain}{% Style for when a plain pagestyle is specified
	\fancyhead{}\renewcommand{\headrulewidth}{0pt}} 

% Removes the header from odd empty pages at the end of chapters
\makeatletter
\renewcommand{\cleardoublepage}{
\clearpage\ifodd\c@page\else
\hbox{}
\vspace*{\fill}
\thispagestyle{empty}
\newpage
\fi}


%----------------------------------------------------------------------------------------
%	MAIN TABLE OF CONTENTS
%----------------------------------------------------------------------------------------
%\usepackage{titletoc} % Required for manipulating the table of contents

\contentsmargin{0cm} % Removes the default margin

% Part text styling (this is mostly taken care of in the PART HEADINGS section of this file)
\titlecontents{part}
	[0cm] % Left indentation
	{\addvspace{20pt}\bfseries} % Spacing and font options for parts
	{}
	{}
	{}

% Chapter text styling
\titlecontents{chapter}
	[1.25cm] % Left indentation
	{\addvspace{12pt}\large\sffamily\bfseries} % Spacing and font options for chapters
	{\color{uMainColor!60}\contentslabel[\Large\thecontentslabel]{1.25cm}\color{uMainColor}} % Formatting of numbered sections of this type
	{\color{uMainColor}} % Formatting of numberless sections of this type
	{\color{uMainColor!60}\normalsize\;\titlerule*[.5pc]{.}\;\thecontentspage} % Formatting of the filler to the right of the heading and the page number

% Section text styling
\titlecontents{section}
	[1.25cm] % Left indentation
	{\addvspace{3pt}\sffamily\bfseries} % Spacing and font options for sections
	{\contentslabel[\thecontentslabel]{1.25cm}} % Formatting of numbered sections of this type
	{} % Formatting of numberless sections of this type
	{\hfill\color{black}\thecontentspage} % Formatting of the filler to the right of the heading and the page number

% Subsection text styling
\titlecontents{subsection}
	[1.25cm] % Left indentation
	{\addvspace{1pt}\sffamily\small} % Spacing and font options for subsections
	{\contentslabel[\thecontentslabel]{1.25cm}} % Formatting of numbered sections of this type
	{} % Formatting of numberless sections of this type
	{\ \titlerule*[.5pc]{.}\;\thecontentspage} % Formatting of the filler to the right of the heading and the page number

%\titlecontents{subsubsection}
%	[1.25cm] % Left indentation
%	{\addvspace{1pt}\sffamily\tiny} % Spacing and font options for subsections
%	{\contentslabel[\thecontentslabel]{1.25cm}} % Formatting of numbered sections of this type
%	{} % Formatting of numberless sections of this type
%	{\ \titlerule*[.5pc]{.}\;\thecontentspage} % Formatting of the filler to the right of the heading and the page number

\titlecontents{subsubsection}
[1.25cm]
  {\addvspace{1pt}\sffamily\tiny}
 {\contentslabel[\thecontentslabel]{1.25cm}}
  {}
  {}

% Figure text styling
\titlecontents{figure}
	[1.25cm] % Left indentation
	{\addvspace{1pt}\sffamily\small} % Spacing and font options for figures
	{\thecontentslabel\hspace*{1em}} % Formatting of numbered sections of this type
	{} % Formatting of numberless sections of this type
	{\ \titlerule*[.5pc]{.}\;\thecontentspage} % Formatting of the filler to the right of the heading and the page number

% Table text styling
\titlecontents{table}
	[1.25cm] % Left indentation
	{\addvspace{1pt}\sffamily\small} % Spacing and font options for tables
	{\thecontentslabel\hspace*{1em}} % Formatting of numbered sections of this type
	{} % Formatting of numberless sections of this type
	{\ \titlerule*[.5pc]{.}\;\thecontentspage} % Formatting of the filler to the right of the heading and the page number

%----------------------------------------------------------------------------------------
%	MINI TABLE OF CONTENTS IN PART HEADS
%----------------------------------------------------------------------------------------

% Chapter text styling
\titlecontents{lchapter}
	[0em] % Left indentation
	{\addvspace{15pt}\large\sffamily\bfseries} % Spacing and font options for chapters
	{\color{uMainColor}\contentslabel[\Large\thecontentslabel]{1.25cm}\color{uMainColor}} % Chapter number
	{}  
	{\color{uMainColor}\normalsize\sffamily\bfseries\;\titlerule*[.5pc]{.}\;\thecontentspage} % Page number

% Section text styling
\titlecontents{lsection}
	[0em] % Left indentation
	{\sffamily\small} % Spacing and font options for sections
	{\contentslabel[\thecontentslabel]{1.25cm}} % Section number
	{}
	{}

% Subsection text styling (note these aren't shown by default, display them by searchings this file for tgcdepth and reading the commented text)
%\titlecontents{lsubsection}
%	[.5em] % Left indentation
%	{\sffamily\footnotesize} % Spacing and font options for subsections
%	{\contentslabel[\thecontentslabel]{1.25cm}}
%	{}
%	{}


\usepackage[]{ccicons}


%----------------------------------------------------------------------------------------
%	THEOREM STYLES
%----------------------------------------------------------------------------------------

\usepackage{amsmath,amsfonts,amssymb,amsthm} % For math equations, theorems, symbols, etc

\newcommand{\intoo}[2]{\mathopen{]}#1\,;#2\mathclose{[}}
\newcommand{\ud}{\mathop{\mathrm{{}d}}\mathopen{}}
\newcommand{\intff}[2]{\mathopen{[}#1\,;#2\mathclose{]}}
\renewcommand{\qedsymbol}{$\blacksquare$}

\ifcsname chapter \endcsname   % verify if command name exists in the CLASS (Book, report, article ...)
	\newtheorem{notation}{Notation}[chapter]
 \else
	\newtheorem{notation}{Notation}[section]
\fi


% Boxed/framed environments
\newtheoremstyle{ocrenumbox}% Theorem style name
{0pt}% Space above
{0pt}% Space below
{\normalfont}% Body font
{}% Indent amount
{\small\bf\sffamily\color{uMainColor}}% Theorem head font
{\;}% Punctuation after theorem head
{0.25em}% Space after theorem head
{\small\sffamily\color{uMainColor}\thmname{#1}\nobreakspace\thmnumber{\@ifnotempty{#1}{}\@upn{#2}}% Theorem text (e.g. Theorem 2.1)
\thmnote{\nobreakspace\the\thm@notefont\sffamily\bfseries\color{black}---\nobreakspace#3.}} % Optional theorem note

\newtheoremstyle{ocrenumbox}% Theorem style name
{0pt}% Space above
{0pt}% Space below
{\normalfont}% Body font
{}% Indent amount
{\small\bf\sffamily\color{uMainColor}}% Theorem head font
{\;}% Punctuation after theorem head
{0.25em}% Space after theorem head
{\small\sffamily\color{uMainColor}\thmname{#1}\nobreakspace\thmnumber{\@ifnotempty{#1}{}\@upn{#2}}% Theorem text (e.g. Theorem 2.1)
\thmnote{\nobreakspace\the\thm@notefont\sffamily\bfseries\color{black}---\nobreakspace#3.}} % Optional theorem note

\newtheoremstyle{blacknumex}% Theorem style name
{5pt}% Space above
{5pt}% Space below
{\normalfont}% Body font
{} % Indent amount
{\small\bf\sffamily}% Theorem head font
{\;}% Punctuation after theorem head
{0.25em}% Space after theorem head
{\small\sffamily{\tiny\ensuremath{\blacksquare}}\nobreakspace\thmname{#1}\nobreakspace\thmnumber{\@ifnotempty{#1}{}\@upn{#2}}% Theorem text (e.g. Theorem 2.1)
\thmnote{\nobreakspace\the\thm@notefont\sffamily\bfseries---\nobreakspace#3.}}% Optional theorem note


\newtheoremstyle{blacknumboxS}% name of the style to be used
{5pt}% measure of space to leave above the theorem. E.g.: 3pt
{5pt}% measure of space to leave below the theorem. E.g.: 3pt
{\normalfont}% name of font to use in the body of the theorem
{}% measure of space to indent
{\small\bf\sffamily}% name of head font
{}% punctuation between head and body
{ }% space after theorem head; " " = normal interword space
{\color{black}\faEye~  \color{uMainColor}\small\sffamily\thmnote{#3 : }}

% Non-boxed/non-framed environments
\newtheoremstyle{ocrenum}% Theorem style name
{5pt}% Space above
{5pt}% Space below
{\normalfont}% Body font
{}% Indent amount
{\small\bf\sffamily\color{uMainColor}}% Theorem head font
{\;}% Punctuation after theorem head
{0.25em}% Space after theorem head
{\small\sffamily\color{uMainColor}\thmname{#1}\nobreakspace\thmnumber{\@ifnotempty{#1}{}\@upn{#2}}% Theorem text (e.g. Theorem 2.1)
\thmnote{\nobreakspace\the\thm@notefont\sffamily\bfseries\color{black}---\nobreakspace#3.}} % Optional theorem note
\makeatother

% Defines the theorem text style for each type of theorem to one of the three styles above



\ifcsname chapter \endcsname

\newcounter{dummy} 

\numberwithin{dummy}{section}
	\theoremstyle{ocrenumbox}
		\newtheorem{theoremeT}[dummy]{Proposition}
		\newtheorem{problem}{Problem}[chapter]
		\newtheorem{exerciseT}{Outillage}[chapter]
	\theoremstyle{blacknumex}
		\newtheorem{exampleT}{Example}[chapter]
	\theoremstyle{blacknumbox}
		\newtheorem{vocabulary}{Vocabulary}[chapter]
		\newtheorem{definitionT}{Definition}[section]
		\newtheorem{corollaryT}[dummy]{Concept}
	\theoremstyle{blacknumboxS}
		\newtheorem{remarqueST}[]{}
		\newtheorem{remarqueT}[dummy]{Remarque}
	\theoremstyle{ocrenum}
	\newtheorem{proposition}[dummy]{Proposition}
  
 \else
 
\newcounter{dummy} 
\numberwithin{dummy}{section}

	\theoremstyle{ocrenumbox}
		\newtheorem{theoremeT}[dummy]{Proposition}
		%\newtheorem{problem}{Problem}[section]
		\newtheorem{exerciseT}{Outillage}[section]
	\theoremstyle{blacknumex}
		\newtheorem{exampleT}{Example}[section]
	\theoremstyle{blacknumbox}
		\newtheorem{vocabulary}{Vocabulary}[section]
		\newtheorem{definitionT}{Definition}[section]
		\newtheorem{corollaryT}[dummy]{concept}
	\theoremstyle{blacknumboxS}
		\newtheorem{remarqueST}[]{}
		\newtheorem{remarqueT}[dummy]{Remarque}
	\theoremstyle{ocrenum}
		\newtheorem{proposition}[dummy]{Proposition}
\fi



%----------------------------------------------------------------------------------------
%	DEFINITION OF COLORED BOXES
%----------------------------------------------------------------------------------------

\RequirePackage[framemethod=default]{mdframed} % Required for creating the theorem, definition, exercise and corollary boxes

% Theorem box
\newmdenv[skipabove=7pt,
skipbelow=7pt,
backgroundcolor=black!5,
linecolor=uMainColor,
innerleftmargin=5pt,
innerrightmargin=5pt,
innertopmargin=5pt,
leftmargin=0cm,
rightmargin=0cm,
innerbottommargin=5pt]{tBox}

% Exercise box	  
\newmdenv[skipabove=7pt,
skipbelow=7pt,
rightline=false,
leftline=true,
topline=false,
bottomline=false,
backgroundcolor=uMainColor!10,
linecolor=uMainColor,
innerleftmargin=5pt,
innerrightmargin=5pt,
innertopmargin=5pt,
innerbottommargin=5pt,
leftmargin=0cm,
rightmargin=0cm,
linewidth=4pt]{eBox}	

% Definition box
\newmdenv[skipabove=7pt,
skipbelow=7pt,
rightline=false,
leftline=true,
topline=false,
bottomline=false,
linecolor=uMainColor,
innerleftmargin=5pt,
innerrightmargin=5pt,
innertopmargin=0pt,
leftmargin=0cm,
rightmargin=0cm,
linewidth=4pt,
innerbottommargin=0pt]{dBox}	

% Corollary box
\newmdenv[skipabove=7pt,
skipbelow=7pt,
rightline=false,
leftline=true,
topline=false,
bottomline=false,
linecolor=gray,
backgroundcolor=black!5,
innerleftmargin=5pt,
innerrightmargin=5pt,
innertopmargin=5pt,
leftmargin=0cm,
rightmargin=0cm,
linewidth=4pt,
innerbottommargin=5pt]{cBox}

% Creates an environment for each type of theorem and assigns it a theorem text style from the "Theorem Styles" section above and a colored box from above
%\newenvironment{theorem}{\begin{tBox}\begin{theoremeT}}{\end{theoremeT}\end{tBox}}
\newenvironment{exercise}{\begin{eBox}\begin{exerciseT}}{\hfill{\color{uMainColor}\tiny\ensuremath{\blacksquare}}\end{exerciseT}\end{eBox}}				  
%\newenvironment{definition}{\begin{dBox}\begin{definitionT}}{\end{definitionT}\end{dBox}}	
%\newenvironment{example}{\begin{exampleT}}{\hfill{\tiny\ensuremath{\blacksquare}}\end{exampleT}}		
%\newenvironment{corollary}{\begin{cBox}\begin{corollaryT}}{\end{corollaryT}\end{cBox}}	
\newenvironment{nota}{\begin{cBox}\begin{remarqueT}}{\end{remarqueT}\end{cBox}}	

%----------------------------------------------------------------------------------------
%	REMARK ENVIRONMENT
%----------------------------------------------------------------------------------------

\newenvironment{remark}{\par\vspace{10pt}\small % Vertical white space above the remark and smaller font size
\begin{list}{}{
\leftmargin=35pt % Indentation on the left
\rightmargin=25pt}\item\ignorespaces % Indentation on the right
\makebox[-2.5pt]{\begin{tikzpicture}[overlay]
\node[draw=uMainColor!60,line width=1pt,circle,fill=uMainColor!25,font=\sffamily\bfseries,inner sep=2pt,outer sep=0pt] at (-15pt,0pt){\textcolor{uMainColor}{i}};\end{tikzpicture}} % Orange R in a circle
\advance\baselineskip -1pt}{\end{list}\vskip5pt} % Tighter line spacing and white space after remark

%----------------------------------------------------------------------------------------
%	SECTION NUMBERING IN THE MARGIN
%----------------------------------------------------------------------------------------

\makeatletter
% \renewcommand{\@seccntformat}[1]{\textcolor{uMainColor}{\csname the#1\endcsname}\quad}   

%\makeatother

%----------------------------------------------------------------------------------------
%	PART HEADINGS
%----------------------------------------------------------------------------------------

% Numbered part in the table of contents
\newcommand{\@mypartnumtocformat}[2]{%
	\setlength\fboxsep{0pt}%
	\noindent\colorbox{uMainColor!20}{\strut\parbox[c][.7cm]{\ecart}{\color{uMainColor!70}\Large\sffamily\bfseries\centering#1}}\hskip\esp\colorbox{uMainColor!40}{\strut\parbox[c][.7cm]{\linewidth-\ecart-\esp}{\Large\sffamily\centering#2}}%
}

% Unnumbered part in the table of contents
\newcommand{\@myparttocformat}[1]{%
	\setlength\fboxsep{0pt}%
	\noindent\colorbox{uMainColor!40}{\strut\parbox[c][.7cm]{\linewidth}{\Large\sffamily\centering#1}}%
}

\newlength\esp
\setlength\esp{4pt}
\newlength\ecart
\setlength\ecart{1.2cm-\esp}
\newcommand{\thepartimage}{}%
\newcommand{\partimage}[1]{\renewcommand{\thepartimage}{#1}}%
\def\@part[#1]#2{%
\ifnum \c@secnumdepth >-2\relax%
\refstepcounter{part}%
\addcontentsline{toc}{part}{\texorpdfstring{\protect\@mypartnumtocformat{\thepart}{#1}}{\partname~\thepart\ ---\ #1}}
\else%
\addcontentsline{toc}{part}{\texorpdfstring{\protect\@myparttocformat{#1}}{#1}}%
\fi%
\startcontents%
\markboth{}{}%
{\thispagestyle{empty}%
\begin{tikzpicture}[remember picture,overlay]%
\node at (current page.north west){\begin{tikzpicture}[remember picture,overlay]%	
\fill[uMainColor!20](0cm,0cm) rectangle (\paperwidth,-\paperheight);
\node[anchor=north] at (4cm,-3.25cm){\color{uMainColor!40}\fontsize{220}{100}\sffamily\bfseries\thepart}; 
\node[anchor=south east] at (\paperwidth-1cm,-\paperheight+1cm){\parbox[t][][t]{8.5cm}{
\printcontents{l}{0}{\setcounter{tocdepth}{1}}% The depth to which the Part mini table of contents displays headings; 0 for chapters only, 1 for chapters and sections and 2 for chapters, sections and subsections
}};
\node[anchor=north east] at (\paperwidth-1.5cm,-3.25cm){\parbox[t][][t]{15cm}{\strut\raggedleft\color{white}\fontsize{30}{30}\sffamily\bfseries#2}};
\end{tikzpicture}};
\end{tikzpicture}}%
\@endpart}
\def\@spart#1{%
\startcontents%
\phantomsection
{\thispagestyle{empty}%
\begin{tikzpicture}[remember picture,overlay]%
\node at (current page.north west){\begin{tikzpicture}[remember picture,overlay]%	
\fill[uMainColor!20](0cm,0cm) rectangle (\paperwidth,-\paperheight);
\node[anchor=north east] at (\paperwidth-1.5cm,-3.25cm){\parbox[t][][t]{15cm}{\strut\raggedleft\color{white}\fontsize{30}{30}\sffamily\bfseries#1}};
\end{tikzpicture}};
\end{tikzpicture}}
\addcontentsline{toc}{part}{\texorpdfstring{%
\setlength\fboxsep{0pt}%
\noindent\protect\colorbox{uMainColor!40}{\strut\protect\parbox[c][.7cm]{\linewidth}{\Large\sffamily\protect\centering #1\quad\mbox{}}}}{#1}}%
\@endpart}
\def\@endpart{\vfil\newpage
\if@twoside
\if@openright
\null
\thispagestyle{empty}%
\newpage
\fi
\fi
\if@tempswa
\twocolumn
\fi}

%----------------------------------------------------------------------------------------
%	CHAPTER HEADINGS
%-------------------------------------------------%--------------------------------------

%
%
%% A switch to conditionally include a picture, implemented by Christian Hupfer
\newif\ifusechapterimage
\usechapterimagetrue
\newcommand{\thechapterimage}{}%
\newcommand{\chapterimage}[1]{\ifusechapterimage\renewcommand{\thechapterimage}{#1}\fi}%
\newcommand{\autodot}{.}

\def\@makechapterhead#1{%
{\parindent \z@ \raggedright \normalfont
\ifnum \c@secnumdepth >\m@ne
\if@mainmatter
\begin{tikzpicture}[remember picture,overlay]
	
\node at (current page.north west)
			{\begin{tikzpicture}[remember picture,overlay]
			\node[anchor=north west,inner sep=0pt] at (0,0) {\ifusechapterimage\includegraphics[width=\paperwidth]{\thechapterimage}\fi};
			\draw[anchor=west] (\Gm@lmargin,-9cm) node [line width=2pt,rounded corners=15pt,draw=uMainColor,fill=white,fill opacity=1,inner sep=15pt]{\strut\makebox[22cm]{}};
			
			\draw[anchor=west] (\Gm@lmargin+.3cm,-9cm) node {\fontsize{20}{30}\selectfont\sffamily\bfseries\color{uMainColor}~#1\strut};
			%{\Huge\sffamily\bfseries\color{uMainColor}\thechapter\autodot~#1\strut};
			\node at (current page.north east)  [xshift=-3.95cm, yshift=-4.05cm, text opacity=1]  {{\color{white}\centering\fontsize{200}{30}\selectfont \bfseries\sffamily\thechapter\strut} };
			\node at (current page.north east)  [xshift=-4cm, yshift=-4cm, text opacity=1]  {{\color{uMainColor}\centering\fontsize{200}{30}\selectfont \bfseries\sffamily\thechapter\strut} };
			
			\end{tikzpicture}};
	\end{tikzpicture}

\else

\begin{tikzpicture}[remember picture,overlay]
\node at (current page.north west)
{\begin{tikzpicture}[remember picture,overlay]
\node[anchor=north west,inner sep=0pt] at (0,0) {\ifusechapterimage\includegraphics[width=\paperwidth]{\thechapterimage}\fi};
\draw[anchor=west] (\Gm@lmargin,-9cm) node [line width=2pt,rounded corners=15pt,draw=uMainColor,fill=white,fill opacity=1,inner sep=15pt]{\strut\makebox[22cm]{}};

\draw[anchor=east] (\Gm@lmargin,-9cm) node [line width=2pt,rounded corners=15pt,draw=uMainColor,fill=white,fill opacity=1,inner sep=15pt]{\strut\makebox[22cm]{}};
\draw[anchor=west] (\Gm@lmargin+.3cm,-9cm) node {\huge\sffamily\bfseries\color{uMainColor}#1\strut};

\end{tikzpicture}};
 
\end{tikzpicture}
\fi\fi\par\vspace*{270\p@}}}

\makeatother
%-------------------------------------------

%\def\@makeschapterhead#1{%
%\begin{tikzpicture}[remember picture,overlay]
%\node at (current page.north west)
%{\begin{tikzpicture}[remember picture,overlay]
%\node[anchor=north west,inner sep=0pt] at (0,0) {\ifusechapterimage\includegraphics[width=\paperwidth]{\thechapterimage}\fi};
%\draw[anchor=west] (\Gm@lmargin,-9cm) node [line width=2pt,rounded corners=15pt,draw=uMainColor,fill=white, fill opacity=0.8,inner sep=15pt]{\strut\makebox[22cm]{}};
%\draw[anchor=west] (\Gm@lmargin+.3cm,-9cm) node {\huge\sffamily\bfseries\color{uMainColor}#1\strut};
%\end{tikzpicture}};
%\end{tikzpicture}
%\par\vspace*{270\p@}}




%----------------------------------------------------------------------------------------
%	LINKS
%----------------------------------------------------------------------------------------


\usepackage{bookmark}
\bookmarksetup{
open,
numbered,
addtohook={%
\ifnum\bookmarkget{level}=0 % chapter
\bookmarksetup{bold}%
\fi
\ifnum\bookmarkget{level}=-1 % part
\bookmarksetup{color=uMainColor,bold}%
\fi
}
}

%----------------------------------------------------------------------------------------
%	FANCY CHAPTER Header
%----------------------------------------------------------------------------------------





%\definecolor{uMainColor}{RGB}{255,128,0}   

%-------------------------------------------------
%	BIBLIOGRAPHY
%-------------------------------------------------

\addbibresource{Latex/Sources/commons/bibliography.bib} % BibTeX bibliography file
\defbibheading{bibempty}{}

%-------------------------------------------------
%	COLORS & BORDERS
%-------------------------------------------------
\setlength{\fboxrule}{0.75pt} % Width of the border around the abstract
%\definecolor{uMainColor}{RGB}{255,128,0} % Color of the article title and section
\definecolor{color2}{RGB}{220,220,220} % Color of the boxes behind the abstract and headings

%-------------------------------------------------
% ITEMS DEFINITION
%-------------------------------------------------

\setlist [itemize,1]{label=\color{uMainColor}\faCaretRight }

\begin{document}


\addtocontents{toc}{\protect\hypertarget{toc}{}}

%-------------------------------------------------
%	ARTICLE INFORMATION
%-------------------------------------------------

\Abstract {\uabstract%===========================
% COURS "INTRO CYBERDEF"
% Abstract général des articles
%===========================


Il fait partie du cours introductif aux fondamentaux de la sécurité des systèmes d'information vue sous deux prismes quelques fois opposés dans la littérature : la gouvernance et la gestion opérationnelle de la sécurité.
Le cours est constitué d'un ensemble de notes de synthèse indépendantes compilées en un document unique, mais édité par chapitre dans le cadre de ce cours.\\
C\edoc ne constitue pas à lui seul le référentiel du cours CYBERDEF101 (SEC101 du Cnam). Il compile des notes de cours mises à disposition de l'auditeur comme support pédagogique partiel à ce cours introductif à la cyberdéfense d'entreprise. }
\JournalInfo{\uJournalInfo} % Journal information
\Archive{Notes de cours éditées le  \DTMnow} % Additional notes (e.g. copyright, DOI, review/research article)
\PaperTitle{\utitle} % Article title
\Authors{\uauthor\textsuperscript{1,}\textsuperscript{2}*} % Authors
\affiliation{\textsuperscript{1}\textit{\uproa}} % Author affiliation
\affiliation{\textsuperscript{2}\textit{\uprob}} % Author affiliation
\affiliation{*\textbf{email}: \umaila\ -- \umailb} % Corresponding author
\Keywords{\ukeywords} % Keywords 
\newcommand{\keywordname}{Mots clefs} % Defines the keywords heading name

%-------------------------------------------------
%	MAKE TITLE
%-------------------------------------------------

%\maketitle

%\begin{tikzpicture}[remember picture,overlay]\node at (current page.north west)[xshift=4cm, yshift=-4cm, text opacity=1.0]{\includegraphics[width=0.1\paperwidth]{\utemplatepath/commons.inc/CommonsPictures/\utemplatetheme.theme/shield-20.pdf}}; 
%\end{tikzpicture} 
%
%
%\AddToShipoutPicture{%
%\AtPageCenter{%
%\begin{tikzpicture}[remember picture,overlay]
%\node [rotate=60,scale=5,text opacity=0.1] at (current page.center) { \textit{\color{uAlertColor}\INFODistrib} };
%\end{tikzpicture} 
%%\begin{tikzpicture}[remember picture,overlay]
%%\node [rotate=60,scale=10,text opacity=0.3]
%%at (current page.center) {moi};
%%\end{tikzpicture}
%}%
%}


%-------------------------------------------------	
%\begin{center}
%	\textcolor{uMainColor}{\textbf{\INFODistrib}} % Distribution information
%\end{center}
%
%\INFOGitHub % Link to download file and copyright
%%-------------------------------------------------

%%-------------------------------------------------------------
%               FR CYBERDEF SECOPS COURSE
%
%                 DEFINITION DES Sources/EDU/acronymes

%                      $File :  Sources/EDU/acronymes.tex
%                             2020 eduf@ction
%-------------------------------------------------------------
% Require : glossaries package
%-------------------------------------------------------------
%\newacronym{XXX}{XXX}{XXX}
%-----------------------------------------------------------

%  Organismes

\newacronym{aCNIL}{CNIL}{Commission National Informatique et Liberté}
\newacronym{aOTAN}{OTAN}{Organisation du Traité de l'Atlantique Nord}
\newacronym{aANSSI}{ANSSI}{Agence Nationale de la Sécurité des systèmes d'information}
\newacronym{aCERT}{CERT}{Computer Emergency Response Team}
\newacronym{aAFNOR}{AFNOR}{Association Française de Normalisation}
\newacronym{aSIEM}{SIEM}{System Incident and Event Management}
\newacronym{aSIM}{SIM}{System Incident Management}
\newacronym{aSEM}{SEM}{System Event Management}
\newacronym{aIRP}{IRP}{Incident Response Plan, Plan de Réponse à Incident}
\newacronym{aISO}{ISO}{International Organization for Standardization}
\newacronym{aDICT}{DICT}{Disponibilité, Intégrité, Confidentialité, Traçabilité}
\newacronym{aPDCA}{PDCA}{Plan, Do, Check, Act (Roue de Deming)}
\newacronym{aGDPR}{GDPR}{General Data Protection Regulation ou RGPD Règlement général sur la protection des données}

%-----------------------------------------------------------

% ROLES ET METIERS

\newacronym{aOS}{OS}{Officier de sécurité}
\newacronym{aRSSI}{RSSI}{Responsable de la sécurité des systèmes d'information}
\newacronym{aDPO}{DPO}{Data Protection Officier}
\newacronym{aCIL}{CIL}{ Correspondant Informatique et Liberté}
\newacronym{aSSI}{SSI}{Sécurité des Systèmes d’information}
\newacronym{aGRC}{GRC}{Gouvernance Risques et Conformité}
\newacronym{aDTF}{DTF}{Defend The Flag}
\newacronym{aCTF}{CTF}{Capture The Flag}
\newacronym{aOIV}{OIV}{Opérateur d'Infrastructures Vitales}
\newacronym{aOSE}{OSE}{Opérateur de Services Essentiels}
\newacronym{aLPM}{LPM}{Loi de Programmation Militaire}
\newacronym{aEBIOS}{EBIOS}{Expression des Besoins et Identification des Objectifs de Sécurité}
\newacronym{aDLL}{DLL}{Dynamic Loads Library}
\newacronym{aACID}{ACID}{Authentification, Confidentialité, Intégrité, Disponibilité}

%-------------------------------------------------------
% Sources/EDU/acronymes DATABASE XLS Generation (cf. RESSOURCES DIR)
%-------------------------------------------------------
% DO NOT MODIFY DIRECTLY, USE XLS File (BEGIN)

\newacronym{a2FA}{2FA}{\textit{Two-Factor Authentication} : Authentification à deux facteurs}
\newacronym{aAAA}{AAA}{\textit{Authentication, Authorization, Accounting} : Authentification, autorisation, comptabilisation}
\newacronym{aAC}{AC}{\textit{Access Control} : Accès conditionnel}
\newacronym{aAES}{AES}{\textit{Advanced Encryption Standard} : Algorithme de chiffrement avancé}
\newacronym{aAPT}{APT}{\textit{Advanced Persistent Threat Group} : Groupe de cyber-espionnage avancé}
\newacronym{aASLR}{ASLR}{\textit{Address Space Layout Randomization} : Randomisation de l'espace d'adressage}
\newacronym{aAV}{AV}{\textit{Antivirus Software} : Logiciel antivirus}
\newacronym{aBOT}{BOT}{\textit{Robot (Malicious)} : Robot (malveillant)}
\newacronym{aCC}{C\&C}{\textit{Command And Control} : Commande et contrôle}
\newacronym{aCSRF}{CSRF}{\textit{Cross-Site Request Forgery Attack} : Attaque de confiance de site à site}
\newacronym{aDDOS}{DDOS}{\textit{Distributed Denial Of Service Attack} : Attaque de déni de service distribué}
\newacronym{aDES}{DES}{\textit{Data Encryption Standard} : Algorithme de chiffrement à clé de 64 bits}
\newacronym{aDLP}{DLP}{\textit{Data Loss Prevention} : Gestion de la protection des données sensibles}
\newacronym{aEDR}{EDR}{\textit{Endpoint Detection And Response} : Détection et réponse aux comportements anormaux d'un terminal}
\newacronym{aEK}{EK}{\textit{Exploit Kit} : Kit d'exploit}
\newacronym{aFTP}{FTP}{\textit{File Transfer Protocol} : Protocole de transfert de fichier}
\newacronym{aHTTPS}{HTTPS}{\textit{Protocol For Securing Communications On The Internet} : Protocole de sécurisation HTTP}
\newacronym{aIAM}{IAM}{\textit{Identity And Access Management} : Identité et Gestion des Accés}
\newacronym{aIDS}{IDS}{\textit{Intrusion Detection System} : Système de détection d'intrusion}
\newacronym{aIIS}{IIS}{\textit{Internet Information Services Server}}
\newacronym{aIPS}{IPS}{\textit{Intrusion Prevention System} : Système de prévention d'incident}
\newacronym{aISA}{ISA}{\textit{Internet Security And Acceleration Server} : Serveur de sécurité et d'accélération Internet (Microsoft)}
\newacronym{aLDAP}{LDAP}{\textit{Lightweight Directory Access Protocol} : Protocole de gestion d'annuaire léger}
\newacronym{aMD5}{MD5}{\textit{Message-Digest Algorithm 5}}
\newacronym{aMFA}{MFA}{\textit{Multi-Factor Authentication} : Authentification multi-facteurs}
\newacronym{aNAC}{NAC}{\textit{Network Access Control}}
\newacronym{aOAuth}{OAuth}{\textit{Open Authorization Protocol}}
\newacronym{aPAM}{PAM}{\textit{Privilege Access Management}}
\newacronym{aPKI}{PKI}{\textit{Public Key Infrastructure} : Infrastructure de clefs publiques}
\newacronym{aRSA}{RSA}{\textit{Rivest-Shamir-Adleman Encryption Algorithm}}
\newacronym{aSAML}{SAML}{\textit{Security Assertion Markup Language}}
\newacronym{aSMS}{SMS}{\textit{Short Message Service}}
\newacronym{aSQLI}{SQLI}{\textit{Sql Injection} : Injection de script de commande SQL}
\newacronym{aSSH}{SSH}{\textit{Secure Shell Protocol}}
\newacronym{aSSL}{SSL}{\textit{Protocol For Securing Communications On The Internet} : Protocole de sécurisation des communications sur Internet}
\newacronym{aTKIP}{TKIP}{\textit{Temporal Key Integrity Protocol}}
\newacronym{aTLS}{TLS}{\textit{Transport Layer Security}}
\newacronym{aURL}{URL}{\textit{Uniform Resource Locator}}
\newacronym{aVPN}{VPN}{\textit{Virtual Private Network} : Réseau privé virtuel}
\newacronym{aWAF}{WAF}{\textit{Web Application Firewall}}
\newacronym{aWPA}{WPA}{\textit{Wifi Protected Access}}
\newacronym{aXSS}{XSS}{\textit{Cross-Site Scripting Attack}}

% DO NOT MODIFY (END)
%-------------------------------------------------------





%-------------------------------------------------
\newpage
\ubody % MAIN BODY defined external
%-------------------------------------------------

	\printbibliography
 \printglossary[type=\acronymtype]
%	\printglossary
%-------------------------------------------------	
\newpage
\Ucontribute % Contribution TEXT
%-------------------------------------------------


%-------------------------------------------------
	\newpage
	\tableofcontents  % TOC
	\iftotalfigures
  			\listoffigures
	\fi
%-------------------------------------------------

\end{document}
}
	\else
	\providecommand{\umainload}{%-------------------------------------------------------------
%               FR CYBERDEF SECOPS COURSE
%                          ARTICLE MAIN FILE
%                             2020 eduf@ction
%-------------------------------------------------------------

\documentclass[10pt,fleqn,twoside]{\utemplatepath/LessonModel/Article/edxstyle} %edxstyle

\usepackage{beamerarticle}

\usepackage{ae}
\usepackage[french]{babel}
\usepackage[utf8]{inputenc}
\usepackage[T1]{fontenc}
\usepackage{lmodern}




\usepackage{titletoc} % Required for manipulating the table of contents

\input{\utemplatepath/commons.inc/CommonsPictures/\utemplatetheme.theme/colors.inc}

%-------------------------------------------------------------
%               FR CYBERDEF SECOPS COURSE
%                                  EDX PACKAGES
%.    Commons packages for BOOK, NOTES (.doc) 
%                            and BEAMER (.prz)
%                              2020 eduf@ction
%-------------------------------------------------------------

%\usepackage[hyphens,spaces,obeyspaces]{url}
%\usepackage[hyphens,spaces,obeyspaces]{url}



\PassOptionsToPackage{hyphens,spaces,obeyspaces}{url}

\usepackage[utf8]{inputenc}
\usepackage[french]{babel}
\usepackage[T1]{fontenc}
%\usepackage{markdown}



\usepackage{ccicons}

\usepackage{tikz} % Required for drawing custom shapes
\usetikzlibrary{tikzmark} 
\usetikzlibrary{mindmap,trees, backgrounds}


\usepackage{verbatim}
\usepackage{wrapfig}
\usepackage{xpatch}

\usepackage{hologo}

\usepackage[figure]{totalcount}

\usepackage{fontawesome}

\usepackage[skins]{tcolorbox} % don't forget SKINS option

\usepackage{listings}
\usepackage{upquote}

\lstset{
upquote=true,columns=flexible, basicstyle=\ttfamily,
language=HTML, 
frameround=tttt,
commentstyle=\color{gray},
identifierstyle=\color{blue},
keywordstyle=\color{uMainColor}\bfseries,
xleftmargin=2em,
xrightmargin=2em,
aboveskip=\topsep,
belowskip=\topsep, 
frame=single,
rulecolor=\color{uMainColor},
backgroundcolor=\color{uMainColor!5},
breaklines,
breakindent=1.5em,
showspaces=false,
showstringspaces=false,
showtabs=false,
xleftmargin=2em,
xrightmargin=2em,
aboveskip=1em,
belowskip=1em,
}



\usepackage{graphicx} % Required for including pictures
%\graphicspath{{Pictures/}} % Specifies the directory where pictures are stored

\usepackage{booktabs} % Required for nicer horizontal rules in tables

\usepackage{setspace}
\setstretch{1,2}

\usepackage{csquotes}

\usepackage{fancyhdr}

\usepackage{makeidx} 

\usepackage{xspace} 

\usepackage{floatflt} 

\usepackage{hyperref}
\usepackage{cleveref}

%pdfdisplaydoctitle={\utitle}, Attention le titre ne doit pas contenir de saut de ligne 

\hypersetup{
pdftitle={\uCoursetittle},
pdfauthor={\uauthor},
hidelinks,
backref=true,
pagebackref=true,
hyperindex=true,
pdfcenterwindow = true,
colorlinks=true,
breaklinks=true,
bookmarks=true,
allcolors=uMainColor,
pdfsubject = {\udescription},
pdfcreator = {MACTEX on TEXPAD},
bookmarksopen=true}


% urlcolor=cnam,
% runcolor=uMainColor,
%\pdfstringdefDisableCommands{%
%  \def\\{}%
%  \def\texttt#1{<#1>}%
%}

\usepackage[style=numeric,citestyle=numeric,sorting=nyt,sortcites=true,autopunct=true,hyperref=true,abbreviate=false,backref=true,backend=biber]{biblatex}

%\usepackage{cleveref} old position

\usepackage{multicol}

%--------------------------------------------------------
% pour la génération du QRCODE en page de garde d'article
\usepackage{qrcode} 

%--------------------------------------------------------
% POur la date d'impression en premiere page d'article
\usepackage[useregional]{datetime2}

%--------------------------------------------------------
% Gestion des glossaires et Sources/EDU/acronymes
% Un fichier acroymes.tex centralisé
\usepackage[acronym]{glossaries}

%\makeglossaries
\makenoidxglossaries 

\loadglsentries{Latex/Sources/commons/acronymes.tex}
\loadglsentries{Latex/Sources/commons/glossaire.tex}


%\renewcommand{\bibname}{R\'ef\'erences}

%--------------------------------------------------------
% Pour exclure des "environnement", ici techworkbox par exemple, travaux pour étudiants, non affiché dans le BOOK.
\usepackage{comment}


%--------------------------------------------------------
% pour filigramme sur les pages
\usepackage{eso-pic}

%--------------------------------------------------------
% for \gappto{\UrlBreaks}{\UrlOrds}
% rather than 
%\makeatletter
% \g@addto@macro{\UrlBreaks}{\UrlOrds}
% \makeatother

\usepackage{etoolbox}

\usepackage{breakurl}


\newcommand{\ucite}[1]{\cite{#1}}
%\newcommand{\ucite}[1]{REF BIB #1}
\newcommand{\ugls}[1]{\gls{#1}}
%\newcommand{\ugls}[1]{GLOSSAIRE (#1)}
%\usepackage[Rejne]{fncychap}

%\usepackage{pythontex} % need shell escape

%
%\makeatletter
%\DeclareRobustCommand{\afilename}[1]{%
% \begingroup
%  % \lstname seems to change hyphens into \textendash
%  \def\textendash{-}%
%  \afilename@parse{#1}%
%  \edef\afilename@base{\detokenize\expandafter{\afilename@base}}%
%  \texttt{\afilename@base.\afilename@ext}%
% \endgroup
%}
%\makeatother

%-------------------------------------------------------------
%               FR CYBERDEF SECOPS COURSE
%                                  EDX PACKAGES
%.    Commons COMMANDS for BOOK, NOTES (.doc) 
%                            and BEAMER (.prz)
%                              2020 eduf@ction
%-------------------------------------------------------------

\newcommand{\ucite}[1]{\cite{#1}}
%\newcommand{\ucite}[1]{REF BIB #1}
\newcommand{\ugls}[1]{\gls{#1}}
%\newcommand{\ugls}[1]{GLOSSAIRE (#1)}
\newcommand{\uac}[1]{\acrshort{#1}}



% \newcommand{\ExtractBeforeLastSlash}[1]{\directlua{extract_before_last_slash("#1")}}
% \newcommand{\ExtractAfterLastSlash}[1]{\directlua{extract_after_last_slash("#1")}}
% \newcommand{\RemoveExtension}[1]{\directlua{remove_extension("#1")}}
% \newcommand{\ExtractExtension}[1]{\directlua{extract_extension("#1")}}


\newcommand{\includepicturesubpath}[1]
{#1}

%{\ExtractBeforeLastSlash{#1}/\upicturethemepath\RemoveExtension{\ExtractAfterLastSlash{#1}}.\ExtractExtension{#1}}
    

%---------------------------------------------------
%										 	AUDIO NOTE 
%---------------------------------------------------

\newcommand\audionote[2]{%
  \note{#2}
  \immediate\write\audioexport{==AudioSlide=============}
  \immediate\write\audioexport{#1}
  \immediate\write\audioexport{\insertframenumber}
  \immediate\write\audioexport{---------------------------------}
  \immediate\write\audioexport{#2}
  \immediate\write\audioexport{---------------------------------}
}%

%---------------------------------------------------
%										 	TEXT STYLE SHORTCUT 
%---------------------------------------------------

\newcommand{\tcb}[1]{\textbf{\textcolor{uMainColor}{#1}}}

\newcommand{\voir}[1]{(Schéma~\ref{#1},~page~\pageref{#1}}

\newcommand{\graphicxextension}{pdf}

\newcommand\myemptypage{
   
    \thispagestyle{empty}
    \addtocounter{page}{-1}
    \newpage
  \null
    %CECI EST UNE PAGE VIDE
       \vfill
    }


\newcommand{\uNA}[2]
{
\newacronym{a#1}{#1}{#2}
}

\newcommand{\coauthor}[4]
{
\noindent (#2) \textbf{#1} - #4 : \textit{#3}
}


\newcommand\UKword[1]{\emph{#1}\xspace}
\newcommand\textUK[1]{\textit{\textbf{#1}}\xspace}
\newcommand\Ucenter[1]{\begin{center}{#1}\end{center}}

\newcommand\head[1]{\textbf{#1}}
\newcommand\tb[1]{\textbf{#1}}
\newcommand{\g}[1]{\og #1 \fg{}}
\newcommand{\uload}[1] {\input{../Latex/#1}}
\newcommand{\ucurcolor}{black}

\newcommand{\mak}[1]{\faicon{\aArrowCircleORight} #1}


%-----------------------------------------------------
%										 	\INFOGitHub 
%---------------------------------------------------


\newcommand{\GITfilename}{https://github.com/edufaction/CYBERDEF101/raw/main/Release/Builder/\jobname.pdf}

%\newcommand{\safeqrcode}[2][]{%
%\qrcode[#1]{\detokenize{#2}} }

\newcommand{\safeqrcode}[2][]{\qrcode[#1]{#2}} 

\newcommand{\INFOGitHub}
{
\begin{center} 
\setstretch{1.5}
 {Vérifiez la disponibilité d'une version plus récente de} \\
 {\ulink{\GITfilename}{\textbf{\jobname.pdf }sur GITHUB CYBERDEF} }   \\ 
{{\huge\ccbyncndeu}}  \\  
{Publication en \href{https://creativecommons.org/licenses/by-nc-nd/2.0/fr/}{Creative Common BY-NC-ND} by eduf@ction}    \\  % Copyright notice
{\safeqrcode[height=4 cm,padding]{\GITfilename}}  
\end{center} 
}

\newcommand{\INFOGitHubS}
{
\noindent{\safeqrcode[height=3 cm]{\GITfilename}} \\
\noindent {\tiny Téléchargez une version à jour} 
}



%-----------------------------------------------------
%										 	\upspicture 
%---------------------------------------------------
\newcommand{\upspicture}[3]
{
\renewcommand{\ucurcolor}{#3}
\resizebox{#2\textwidth}{!}{\input{../#1}}
}


%-----------------------------------------------------
%										 	\uindex 
%---------------------------------------------------


\newcommand{\uindex}[1]
{
#1\index{#1}\xspace
}

\newcommand{\ulindex}[1]{\index{#1}}

%-----------------------------------------------------
%										 	\includer 
%---------------------------------------------------

\newcommand{\includer}[1]
{
\input{../Latex/Sources/EDU/SRC1/Chapters/#1}
}

%-----------------------------------------------------
%										 	\edxdico 
%---------------------------------------------------

\newcommand{\edxdico}[2] {\textbf{#1}:\textit{#2} \index{#1}}

%-----------------------------------------------------
%  command DEF : InTextx
%-----------------------------------------------------
\newcommand\InTexta{no text}
\newcommand\InTextb{no text}
\newcommand\InTextc{no text}
\newcommand\rnc{\renewcommand}


%-----------------------------------------------------
%										 	ENV warningbox
%---------------------------------------------------

\newtcolorbox{warningbox}[2][]
{
  colframe = uMainColor!25,
  colback  = black!05,
  coltitle = uMainColor!20!black,
  fonttitle = \bfseries, 
  title    = #2,
  #1,
}


\newtcolorbox{techworkbox}[2][]{sharp corners,colframe = uMainColor!25,colback  = black!05,coltitle = uMainColor!20!black,fonttitle = \bfseries, title={\textcolor{uMainColor}{\faFile~\small {Sujet de mémoire SECOPS}}} : #2,#1,}


\newtcolorbox{notebox2}[2][]
{
sharp corners, 
colback = uMainColor!5!white, 
colframe = uMainColor!75!black,
fonttitle = \bfseries, 
colbacktitle= uMainColor!85!black,
title=#2,#1
}

\newtcolorbox{notebox}[2][]{colback=uMainColor!5!white,
colframe=uMainColor!75!black,fonttitle=\bfseries,
colbacktitle=uMainColor!85!black,enhanced,
attach boxed title to top right={yshift=+1mm},
title=#2,#1}


\newtcolorbox{toolsboxin}[2][]{sharp corners, colback=uMainColor!0!white,
colframe= black!25,fonttitle=\bfseries,
colbacktitle=black!60,enhanced,
attach boxed title to top left={yshift=+1mm},
title=#2,#1}


\newtcolorbox{techworkin}[2][]{sharp corners, colback=uMainColor!5!white,
colframe= uMainColor!100, fonttitle=\bfseries,
colbacktitle=black!5, enhanced,
 coltitle=black!50,
attach boxed title to top right ={yshift=+1mm},
title=#2,#1}



%-----------------------------------------------------
%										 	ENV ToolsBox
%---------------------------------------------------


\newcommand\tooleditor{No EDITOR in Datas.tex}
\newcommand\toolurl{No URL INSERTED in Datas.tex}
\newcommand\toolclass{No EDITOR in Datas.tex}
\newcommand\toolname{No NAME in Datas.tex}
\newcommand\toolanalyst{No ANALYST in Datas.tex}

\newcommand{\toolsbox}[2]
{
\input{../Latex/Sources/EDU/SRC1/SecComponents/#2.#1/datas.tex} 
\begin{toolsboxin}[sidebyside,righthand width=.7\textwidth]{#1}
\includegraphics[width=0.9\textwidth]{../Latex/Sources/EDU/SRC1/SecComponents/#2.#1/icon.png} \\
\tcblower
\input{../Latex/Sources/EDU/SRC1/SecComponents/#2.#1/description.tex}
\end{toolsboxin}
\noindent \small {\faCogs~Classe : \uppercase{\textbf{#2}}, Site de référence : \textbf{\ulink{\toolurl}{#1}}}
\newline
\small {\faGlobe~Editeur : \textbf{\tooleditor}}
\small {\faEye~Analyste : \textbf{\toolanalyst}}
}

%\renewcommand\toolsclass{SIEM}
%\renewcommand\toolsname{RSA Netwitness}
%\renewcommand\toolseditor{RSA}
%\renewcommand\toolsurl{https://www.rsa.com/fr-fr/products/threat-detection-response}

%-----------------------------------------------------
%										 	ENV  TechWork
%---------------------------------------------------


\newcommand{\techwork}[2]
{
\begin{techworkin} [righthand width=.8\textwidth]{\faCogs~\small {Fiche TECHNO} : #1}#2
\end{techworkin}
}

%-----------------------------------------------------
%										 	\uchap
%---------------------------------------------------
\ifdefined\INTERTITLE
\providecommand{\uchap}[1]
	{
	\begin{LARGE}
	\textbf{\Ucenter{#1 \\ -oOo-}}	
	\end{LARGE}
	}
\else
	\providecommand{\uchap}[1]{}
\fi

%-----------------------------------------------------
%											\uexpand
%---------------------------------------------------
\newcommand {\uexpand}[1]{
\mode<all>\input{../Latex/Sources/EDU/SRC1/Chapters/#1}
}


%-----------------------------------------------------
%											\spicture 
% direct include graphic with Scrivener
%-----------------------------------------------------

\newcommand{\spicture}[3]{
\begin{figure}[!hbtp] %h !hbtp
  \begin{center}
	 \includegraphics[width=0.8\textwidth]{#1}
  \end{center}
\caption{#3 - (slbl-#2)\label{slbl-#2}}
\end{figure}
}

%\StrBefore{#2}{.pdf}

%-----------------------------------------------------
%											\upicture 
%---------------------------------------------------

\ifdefined\PRZMODE

\providecommand{\upicture}[4]{
\framesubtitle{#2}
%\begin{figure}[!h] %h
  \begin{center}
	 \includegraphics[height=#3\textheight, width=#3\textwidth, keepaspectratio ]{\includepicturesubpath{#1.pdf}}
  \end{center}
%  \end{figure}
}

\else

\providecommand{\upicture}[4]{
\begin{figure}[!hbtp] %h !hbtp
  \begin{center}
	 \includegraphics[width=#3\textwidth]{\includepicturesubpath{#1.pdf}}
  \end{center}
\caption{#2 \label{#4}}
\end{figure}
}
	
\fi


\ifdefined\PRZMODE

\providecommand{\uwpicture}[4]{
\framesubtitle{#2}
%\begin{figure}[!h] %h
  \begin{center}
	 \includegraphics[height=0.8\textheight, width=0.8\textwidth, keepaspectratio ]{\includepicturesubpath{#1.pdf}}
  \end{center}
%  \end{figure}
}
\else
\providecommand{\uwpicture}[4]
{
\begin{wrapfigure}{O}{#3\textwidth} % inside / Outside
  \begin{center}
	 \includegraphics[width=#3\textwidth]{\includepicturesubpath{#1.pdf}}
\end{center}
\caption{#2 \label{#4}}
\end{wrapfigure}
}
\fi

%-----------------------------------------------------
%	\picframe for BEAMER (Image included in Beamer Slide)
%---------------------------------------------------

\ifdefined\PRZMODE

\providecommand{\picframe}[4]
{
\begin{frame}
\frametitle<presentation>{#2}
  \begin{center}

   \begin{tikzpicture}[remember picture, overlay]
    \node at (current page.center) {\includegraphics[width=#3\paperwidth, height=#3\paperheight,keepaspectratio]{\includepicturesubpath{#1.pdf}}};
    \end{tikzpicture}

  \end{center}
\end{frame}

}
\else
\providecommand{\picframe}[4]
{
\begin{figure}[!hbtp] % !hbtp ou !h 
\centering 
%{\includegraphics[width=#3\paperwidth, height=#3\paperheight,keepaspectratio]{\includepicturesubpath{#1.pdf}}}
\includegraphics[width=#3\textwidth]{\includepicturesubpath{#1.pdf}}
\caption{#2 \label{#4}}
\end{figure}
}
\fi

%-----------------------------------------------------
%.  for BEAMER (Image included in Beamer Slide)
%	\picframeS Simplified with Name and Label from One Param
%---------------------------------------------------


\ifdefined\PRZMODE

\providecommand{\picframeS}[2]
{
\begin{frame}
\frametitle<presentation>{#2}
  \begin{center}
	 \includegraphics[height=0.95\textheight, width=0.95\textwidth, keepaspectratio ]{\picpath/img-#1.\graphicxextension}
  \end{center}
\end{frame}

}
\else
\providecommand{\picframeS}[2]
{
\begin{figure}[!hbtp] % !hbtp ou !h 
\centering \includegraphics[width=0.9\textwidth]{\picpath/img-#1.\graphicxextension}
\caption{(s) #2 \label{lbl:#1}}
\end{figure}
}

\fi


%{figure}[!h]
%{wrapfigure}{O}{#3\textwidth}
%-----------------------------------------------------
%	FRAME			\wpicframe 
%---------------------------------------------------

\ifdefined\PRZMODE

\providecommand{\wpicframe}[4]
{
\begin{frame}
\frametitle<presentation>{#2}
  \begin{center}
	 \includegraphics[height=0.95\textheight, width=0.95\textwidth, keepaspectratio ]{#1.\graphicxextension}
  \end{center}
\end{frame}

}
\else
\providecommand{\wpicframe}[4]
{
\begin{wrapfigure}{O}{#3\textwidth}
\centering \includegraphics[width=#3\textwidth]{#1.\graphicxextension}
\caption{#2 \label{#4}}
\end{wrapfigure}
}
\fi

%-----------------------------------------------------
%	FRAME										\texframe 
%---------------------------------------------------
\ifdefined\PRZMODE

\providecommand{\texframe}[3]
{
\begin{frame}
\frametitle<presentation>{#1}
\framesubtitle<presentation>{#2}
#3
\end{frame}
}
\else
\providecommand{\texframe}[3]
{
#3
}

\fi

%-----------------------------------------------------
%											\uref 
%---------------------------------------------------
%\newcommand{\uref}[2]{(Voir~#1~\cref{#2} page~\pageref{#2})}

\newcommand{\uref}[2]{(Voir~#1~\ref{#2} page~\pageref{#2})}
%-----------------------------------------------------
%											\ubg
%---------------------------------------------------
\newcommand{\ubg}[1]{
\textbf{\g{\uppercase{#1}}}
}
%-----------------------------------------------------
%											\updfimage 
%---------------------------------------------------

\newcommand{\updfimage}[3]
{
\begin{wrapfigure}{R}{#3\textwidth}
  \begin{center}
	 \resizebox{#3\textwidth}{!}{\includegraphics{\includepicturesubpath{#1.pdf}}}
  \end{center}
\caption{#2}
\end{wrapfigure}
}


%************************PACKAGE***

%\usepackage{tikz}% Required for drawing custom shapes
\usetikzlibrary{shadows}


%-----------------------------------------------------
%											 		\rem 
% insertion d'une remarque avec R
%---------------------------------------------------
\newcommand\rem[1]{%
   \marginpar{%
   \tikzpicture[baseline={(title.base)}]
      \node[inner sep=5pt,text width=4cm,drop shadow={shadow yshift=-5pt,shadow xshift=5pt,uMainColor},fill=white] (box) {\vskip5pt \nointerlineskip #1};
      \node[right=10pt,inner sep=0pt,outer sep=10pt] (title) at (box.north west) {\bfseries\color{uMainColor}Remarque};
      \draw[draw=uMainColor,very thick](title.west)--(box.north west)--(box.south west)--(box.south east)--(box.north east)--(title.east);
      \fill[uMainColor]([yshift=-10pt]box.north west)--+(-5pt,-5pt)--+(0pt,-10pt);
   \endtikzpicture}%
}
%-----------------------------------------------------
%											 \utikzimage 
% insertion d'une image sur fichier tkz.tex
%---------------------------------------------------

\newcommand{\utikzimage}[3]
{
\begin{wrapfigure}{R}{#3\textwidth}
  \begin{center}
	 \resizebox{#3\textwidth}{!}{\input{../#1.tkz.tex}}
  \end{center}
\caption{#2}
\end{wrapfigure}
}

%-----------------------------------------------------
%														\link 
% Url et footnote du lien
%-----------------------------------------------------

% TT every character and hyphenate after it
\def\hyphenateAndTtWholeString #1{\xHyphenate#1$\wholeString\unskip}

\def\xHyphenate#1#2\wholeString {\if#1$%
    \else\transform{#1}%
    \takeTheRest#2\ofTheString\fi}

\def\takeTheRest#1\ofTheString\fi
{\fi \xHyphenate#1\wholeString}

\def\transform#1{\url{#1}\hskip 0pt plus 1pt}

% Define the \urlx command which works like \url, but with line brakes
\def\urlx #1{\href{#1}{\hyphenateAndTtWholeString{#1}}}
  

 %-----------------------------------------------------
%														\link 
%-----------------------------------------------------  
  
  \newcommand{\ulink}[2]
  {\href{#1}
  {#2~\raisebox{-0.2ex}{\faExternalLink}\footnote{\url{#1}}}\xspace}
  
  \newcommand{\wikipedia}[1]
  { \href{https://fr.wikipedia.org/wiki/#1}{\textcolor{grey}{#1}}\raisebox{0.5ex}{\tiny{\textcolor{uMainColor}{\faWikipediaW}}}}
  


 %-----------------------------------------------------
%														\utodo 
%-----------------------------------------------------

  \newcommand{\utodo}{\begin{center}{\noindent \textcolor{cyan}{\faGears~{En cours de rédaction, DRAFT non publiable} \faGears}}\end{center}}
  
    \newcommand{\utocomplete}{\begin{center}{\noindent \textcolor{blue}{\faClipboard~{à compléter, les éléments ne donnent qu'une vue trop réduite ou trop parcellaire du sujet}}}\end{center}}
  

 




%-------------------------------------------------------------
%               FR CYBERDEF SECOPS COURSE
%
%                                    EDX STYLE
%
%                              2020 eduf@ction
%-------------------------------------------------------------
\usepackage{enumitem}

\usepackage{parskip}
\setlength{\parindent}{0em}
\setlength{\parskip}{5pt} % 1ex plus 0.5ex minus 0.2ex}

\def\uCoverColor{}

%---------------------------------------------------------------
%	MARGINS
%---------------------------------------------------------------

\usepackage{geometry}

\geometry{
	paper=a4paper, 
%	top=1cm, 
%	bottom=1cm, 
%	left=1cm, 
%	right=1cm, 
%	headheight=20pt, 
%	footskip=1.4cm, 
%	headsep=10pt, 
%	showframe, 
	showcrop 
}

%---------------------------------------------------------------
%	FOOTNOTE (format)
%---------------------------------------------------------------

\makeatletter
\long\def\@makefntextFB#1{%
    \ifx\thefootnote\ftnISsymbol
        \@makefntextORI{#1}%
    \else
        \rule\z@\footnotesep
        \setbox\@tempboxa\hbox{\@thefnmark}%
            \ifdim\wd\@tempboxa>\z@
                \kern2em\llap{\@thefnmark.\kern0.5em}%
            \fi
        \hangindent2em\hangafter\@ne#1
    \fi}
\makeatother

%-----------------------------------------------------
%	COLOR
%-----------------------------------------------------
%\usepackage{xcolor} % Required for specifying colors by name
%\definecolor{uMainColor}{RGB}{160,0,0} 
\colorlet{edxcolorcover}{uCoverColor}
%\definecolor{edxcolorcover}{RGB}{160,0,0} 

%\definecolor{cnam}{RGB}{128, 0, 32} 
                 

  \definecolor{comments}{rgb}{0.7,0,0}    % rouge foncé
  \definecolor{link}{rgb}{0,0.4,0.6}      % ~RoyalBlue de dvips
  \definecolor{url}{rgb}{0.6,0,0}         % rouge-brun
  \definecolor{citation}{rgb}{0,0.5,0}    % vert foncé
  \definecolor{ULlinkcolor}{rgb}{0,0,0.3} % de ulthese.cls
  \definecolor{rouge}{rgb}{0.85,0,0.07}   % rouge bandeau identitaire
  \definecolor{or}{rgb}{1,0.8,0}          % or bandeau identitaire

%-----------------------------------------------------
%	FONTS
%-----------------------------------------------------

%%\usepackage{avant} % Use the Avantgarde font for headings
%\usepackage{times} % Use the Times font for headings
%\usepackage{mathptmx} % Use the Adobe Times Rgman as the default text font together with math symbols from the Sym­bol, Chancery and Com­puter Modern fonts

%%\usepackage{microtype} % Slightly tweak font spacing for aesthetics
%%\usepackage[utf8]{inputenc} % Required for including letters with accents
%\usepackage[T1]{fontenc} % Use 8-bit encoding that has 256 glyphs

% Using AVANT GARDE Family

%\renewcommand{\familydefault}{\sfdefault}

%\fontfamily{pag}\selectfont
%\renewcommand{\familydefault}{pag} % sfdefault pag lmss

%bch         Charter
%lmr         Latin Modern Roman
%lmss        Latin Modern Sans Serif
%lmssq       Latin Modern Sans Serif extended
%lmtt        Latin Modern Typewriter
%lmvtt       Latin Modern Typewriter proportional
%pag         Avant Garde
%pbk         Bookman
%pcr         Courier
%phv         Helvetica
%pnc         New Century Schoolbook
%ppl         Palatino
%ptm         Times
%put         Utopia
 
%\renewcommand{\rmdefault}{pag} % text 
%\renewcommand{\sfdefault}{pag} % titre


%\renewcommand{\rmdefault}{pag} % text 
%\renewcommand{\sfdefault}{phv} % titre


%-----------------------------------------------------
%	BIBLIOGRAPHY AND INDEX
%-----------------------------------------------------


\usepackage{calc} % For simpler calculation - used for spacing the index letter headings correctly

\usepackage{imakeidx}
%\makeindex % Tells LaTeX to create the files required for indexing

\makeindex[columns=2,intoc=true,options={-s \upath/indexstyle.ist}]

%-----------------------------------------------------
%	HEADERS AND FOOTERS
%-----------------------------------------------------


\usepackage{fancyhdr} % Required for header and footer configuration
\pagestyle{fancy} % Enable the custom headers and footers

\ifcsname chapter \endcsname   % verify if command name exists in the CLASS (Book, report, article ...)
% BEGIN IF
\renewcommand{\chaptermark}[1]{\markboth{\sffamily\normalsize\bfseries\chaptername\ \thechapter.\ #1}{}} % Styling for the current chapter in the header
  \else
%nothing
\fi
% END IF

\renewcommand{\sectionmark}[1]{\markright{\sffamily\normalsize\thesection\hspace{5pt}#1}{}} % Styling for the current section in the header

\fancyhf{} % Clear default headers and footers
\fancyhead[LE,RO]{\sffamily\normalsize\thepage} % Styling for the page number in the header
\fancyhead[LO]{\sffamily\normalsize\rightmark} % Print the nearest section name on the left side of odd pages
\fancyhead[RE]{\sffamily\normalsize\leftmark} % Print the current chapter name on the right side of even pages
\fancyfoot[RE,LO]{\sffamily\normalsize\uCoursetittle} % Uncomment to include a footer
\fancyfoot[LE,RO]{\hyperlink{toc}{\sffamily\normalsize\uMinilogo}} % Uncomment to include a footer
% \hyperlink{toc}
\renewcommand{\headrulewidth}{1.0pt} % Thickness of the rule under the header

\fancypagestyle{plain}{% Style for when a plain pagestyle is specified
	\fancyhead{}\renewcommand{\headrulewidth}{0pt}} 

% Removes the header from odd empty pages at the end of chapters
\makeatletter
\renewcommand{\cleardoublepage}{
\clearpage\ifodd\c@page\else
\hbox{}
\vspace*{\fill}
\thispagestyle{empty}
\newpage
\fi}


%----------------------------------------------------------------------------------------
%	MAIN TABLE OF CONTENTS
%----------------------------------------------------------------------------------------
%\usepackage{titletoc} % Required for manipulating the table of contents

\contentsmargin{0cm} % Removes the default margin

% Part text styling (this is mostly taken care of in the PART HEADINGS section of this file)
\titlecontents{part}
	[0cm] % Left indentation
	{\addvspace{20pt}\bfseries} % Spacing and font options for parts
	{}
	{}
	{}

% Chapter text styling
\titlecontents{chapter}
	[1.25cm] % Left indentation
	{\addvspace{12pt}\large\sffamily\bfseries} % Spacing and font options for chapters
	{\color{uMainColor!60}\contentslabel[\Large\thecontentslabel]{1.25cm}\color{uMainColor}} % Formatting of numbered sections of this type
	{\color{uMainColor}} % Formatting of numberless sections of this type
	{\color{uMainColor!60}\normalsize\;\titlerule*[.5pc]{.}\;\thecontentspage} % Formatting of the filler to the right of the heading and the page number

% Section text styling
\titlecontents{section}
	[1.25cm] % Left indentation
	{\addvspace{3pt}\sffamily\bfseries} % Spacing and font options for sections
	{\contentslabel[\thecontentslabel]{1.25cm}} % Formatting of numbered sections of this type
	{} % Formatting of numberless sections of this type
	{\hfill\color{black}\thecontentspage} % Formatting of the filler to the right of the heading and the page number

% Subsection text styling
\titlecontents{subsection}
	[1.25cm] % Left indentation
	{\addvspace{1pt}\sffamily\small} % Spacing and font options for subsections
	{\contentslabel[\thecontentslabel]{1.25cm}} % Formatting of numbered sections of this type
	{} % Formatting of numberless sections of this type
	{\ \titlerule*[.5pc]{.}\;\thecontentspage} % Formatting of the filler to the right of the heading and the page number

%\titlecontents{subsubsection}
%	[1.25cm] % Left indentation
%	{\addvspace{1pt}\sffamily\tiny} % Spacing and font options for subsections
%	{\contentslabel[\thecontentslabel]{1.25cm}} % Formatting of numbered sections of this type
%	{} % Formatting of numberless sections of this type
%	{\ \titlerule*[.5pc]{.}\;\thecontentspage} % Formatting of the filler to the right of the heading and the page number

\titlecontents{subsubsection}
[1.25cm]
  {\addvspace{1pt}\sffamily\tiny}
 {\contentslabel[\thecontentslabel]{1.25cm}}
  {}
  {}

% Figure text styling
\titlecontents{figure}
	[1.25cm] % Left indentation
	{\addvspace{1pt}\sffamily\small} % Spacing and font options for figures
	{\thecontentslabel\hspace*{1em}} % Formatting of numbered sections of this type
	{} % Formatting of numberless sections of this type
	{\ \titlerule*[.5pc]{.}\;\thecontentspage} % Formatting of the filler to the right of the heading and the page number

% Table text styling
\titlecontents{table}
	[1.25cm] % Left indentation
	{\addvspace{1pt}\sffamily\small} % Spacing and font options for tables
	{\thecontentslabel\hspace*{1em}} % Formatting of numbered sections of this type
	{} % Formatting of numberless sections of this type
	{\ \titlerule*[.5pc]{.}\;\thecontentspage} % Formatting of the filler to the right of the heading and the page number

%----------------------------------------------------------------------------------------
%	MINI TABLE OF CONTENTS IN PART HEADS
%----------------------------------------------------------------------------------------

% Chapter text styling
\titlecontents{lchapter}
	[0em] % Left indentation
	{\addvspace{15pt}\large\sffamily\bfseries} % Spacing and font options for chapters
	{\color{uMainColor}\contentslabel[\Large\thecontentslabel]{1.25cm}\color{uMainColor}} % Chapter number
	{}  
	{\color{uMainColor}\normalsize\sffamily\bfseries\;\titlerule*[.5pc]{.}\;\thecontentspage} % Page number

% Section text styling
\titlecontents{lsection}
	[0em] % Left indentation
	{\sffamily\small} % Spacing and font options for sections
	{\contentslabel[\thecontentslabel]{1.25cm}} % Section number
	{}
	{}

% Subsection text styling (note these aren't shown by default, display them by searchings this file for tgcdepth and reading the commented text)
%\titlecontents{lsubsection}
%	[.5em] % Left indentation
%	{\sffamily\footnotesize} % Spacing and font options for subsections
%	{\contentslabel[\thecontentslabel]{1.25cm}}
%	{}
%	{}


\usepackage[]{ccicons}


%----------------------------------------------------------------------------------------
%	THEOREM STYLES
%----------------------------------------------------------------------------------------

\usepackage{amsmath,amsfonts,amssymb,amsthm} % For math equations, theorems, symbols, etc

\newcommand{\intoo}[2]{\mathopen{]}#1\,;#2\mathclose{[}}
\newcommand{\ud}{\mathop{\mathrm{{}d}}\mathopen{}}
\newcommand{\intff}[2]{\mathopen{[}#1\,;#2\mathclose{]}}
\renewcommand{\qedsymbol}{$\blacksquare$}

\ifcsname chapter \endcsname   % verify if command name exists in the CLASS (Book, report, article ...)
	\newtheorem{notation}{Notation}[chapter]
 \else
	\newtheorem{notation}{Notation}[section]
\fi


% Boxed/framed environments
\newtheoremstyle{ocrenumbox}% Theorem style name
{0pt}% Space above
{0pt}% Space below
{\normalfont}% Body font
{}% Indent amount
{\small\bf\sffamily\color{uMainColor}}% Theorem head font
{\;}% Punctuation after theorem head
{0.25em}% Space after theorem head
{\small\sffamily\color{uMainColor}\thmname{#1}\nobreakspace\thmnumber{\@ifnotempty{#1}{}\@upn{#2}}% Theorem text (e.g. Theorem 2.1)
\thmnote{\nobreakspace\the\thm@notefont\sffamily\bfseries\color{black}---\nobreakspace#3.}} % Optional theorem note

\newtheoremstyle{ocrenumbox}% Theorem style name
{0pt}% Space above
{0pt}% Space below
{\normalfont}% Body font
{}% Indent amount
{\small\bf\sffamily\color{uMainColor}}% Theorem head font
{\;}% Punctuation after theorem head
{0.25em}% Space after theorem head
{\small\sffamily\color{uMainColor}\thmname{#1}\nobreakspace\thmnumber{\@ifnotempty{#1}{}\@upn{#2}}% Theorem text (e.g. Theorem 2.1)
\thmnote{\nobreakspace\the\thm@notefont\sffamily\bfseries\color{black}---\nobreakspace#3.}} % Optional theorem note

\newtheoremstyle{blacknumex}% Theorem style name
{5pt}% Space above
{5pt}% Space below
{\normalfont}% Body font
{} % Indent amount
{\small\bf\sffamily}% Theorem head font
{\;}% Punctuation after theorem head
{0.25em}% Space after theorem head
{\small\sffamily{\tiny\ensuremath{\blacksquare}}\nobreakspace\thmname{#1}\nobreakspace\thmnumber{\@ifnotempty{#1}{}\@upn{#2}}% Theorem text (e.g. Theorem 2.1)
\thmnote{\nobreakspace\the\thm@notefont\sffamily\bfseries---\nobreakspace#3.}}% Optional theorem note


\newtheoremstyle{blacknumboxS}% name of the style to be used
{5pt}% measure of space to leave above the theorem. E.g.: 3pt
{5pt}% measure of space to leave below the theorem. E.g.: 3pt
{\normalfont}% name of font to use in the body of the theorem
{}% measure of space to indent
{\small\bf\sffamily}% name of head font
{}% punctuation between head and body
{ }% space after theorem head; " " = normal interword space
{\color{black}\faEye~  \color{uMainColor}\small\sffamily\thmnote{#3 : }}

% Non-boxed/non-framed environments
\newtheoremstyle{ocrenum}% Theorem style name
{5pt}% Space above
{5pt}% Space below
{\normalfont}% Body font
{}% Indent amount
{\small\bf\sffamily\color{uMainColor}}% Theorem head font
{\;}% Punctuation after theorem head
{0.25em}% Space after theorem head
{\small\sffamily\color{uMainColor}\thmname{#1}\nobreakspace\thmnumber{\@ifnotempty{#1}{}\@upn{#2}}% Theorem text (e.g. Theorem 2.1)
\thmnote{\nobreakspace\the\thm@notefont\sffamily\bfseries\color{black}---\nobreakspace#3.}} % Optional theorem note
\makeatother

% Defines the theorem text style for each type of theorem to one of the three styles above



\ifcsname chapter \endcsname

\newcounter{dummy} 

\numberwithin{dummy}{section}
	\theoremstyle{ocrenumbox}
		\newtheorem{theoremeT}[dummy]{Proposition}
		\newtheorem{problem}{Problem}[chapter]
		\newtheorem{exerciseT}{Outillage}[chapter]
	\theoremstyle{blacknumex}
		\newtheorem{exampleT}{Example}[chapter]
	\theoremstyle{blacknumbox}
		\newtheorem{vocabulary}{Vocabulary}[chapter]
		\newtheorem{definitionT}{Definition}[section]
		\newtheorem{corollaryT}[dummy]{Concept}
	\theoremstyle{blacknumboxS}
		\newtheorem{remarqueST}[]{}
		\newtheorem{remarqueT}[dummy]{Remarque}
	\theoremstyle{ocrenum}
	\newtheorem{proposition}[dummy]{Proposition}
  
 \else
 
\newcounter{dummy} 
\numberwithin{dummy}{section}

	\theoremstyle{ocrenumbox}
		\newtheorem{theoremeT}[dummy]{Proposition}
		%\newtheorem{problem}{Problem}[section]
		\newtheorem{exerciseT}{Outillage}[section]
	\theoremstyle{blacknumex}
		\newtheorem{exampleT}{Example}[section]
	\theoremstyle{blacknumbox}
		\newtheorem{vocabulary}{Vocabulary}[section]
		\newtheorem{definitionT}{Definition}[section]
		\newtheorem{corollaryT}[dummy]{concept}
	\theoremstyle{blacknumboxS}
		\newtheorem{remarqueST}[]{}
		\newtheorem{remarqueT}[dummy]{Remarque}
	\theoremstyle{ocrenum}
		\newtheorem{proposition}[dummy]{Proposition}
\fi



%----------------------------------------------------------------------------------------
%	DEFINITION OF COLORED BOXES
%----------------------------------------------------------------------------------------

\RequirePackage[framemethod=default]{mdframed} % Required for creating the theorem, definition, exercise and corollary boxes

% Theorem box
\newmdenv[skipabove=7pt,
skipbelow=7pt,
backgroundcolor=black!5,
linecolor=uMainColor,
innerleftmargin=5pt,
innerrightmargin=5pt,
innertopmargin=5pt,
leftmargin=0cm,
rightmargin=0cm,
innerbottommargin=5pt]{tBox}

% Exercise box	  
\newmdenv[skipabove=7pt,
skipbelow=7pt,
rightline=false,
leftline=true,
topline=false,
bottomline=false,
backgroundcolor=uMainColor!10,
linecolor=uMainColor,
innerleftmargin=5pt,
innerrightmargin=5pt,
innertopmargin=5pt,
innerbottommargin=5pt,
leftmargin=0cm,
rightmargin=0cm,
linewidth=4pt]{eBox}	

% Definition box
\newmdenv[skipabove=7pt,
skipbelow=7pt,
rightline=false,
leftline=true,
topline=false,
bottomline=false,
linecolor=uMainColor,
innerleftmargin=5pt,
innerrightmargin=5pt,
innertopmargin=0pt,
leftmargin=0cm,
rightmargin=0cm,
linewidth=4pt,
innerbottommargin=0pt]{dBox}	

% Corollary box
\newmdenv[skipabove=7pt,
skipbelow=7pt,
rightline=false,
leftline=true,
topline=false,
bottomline=false,
linecolor=gray,
backgroundcolor=black!5,
innerleftmargin=5pt,
innerrightmargin=5pt,
innertopmargin=5pt,
leftmargin=0cm,
rightmargin=0cm,
linewidth=4pt,
innerbottommargin=5pt]{cBox}

% Creates an environment for each type of theorem and assigns it a theorem text style from the "Theorem Styles" section above and a colored box from above
%\newenvironment{theorem}{\begin{tBox}\begin{theoremeT}}{\end{theoremeT}\end{tBox}}
\newenvironment{exercise}{\begin{eBox}\begin{exerciseT}}{\hfill{\color{uMainColor}\tiny\ensuremath{\blacksquare}}\end{exerciseT}\end{eBox}}				  
%\newenvironment{definition}{\begin{dBox}\begin{definitionT}}{\end{definitionT}\end{dBox}}	
%\newenvironment{example}{\begin{exampleT}}{\hfill{\tiny\ensuremath{\blacksquare}}\end{exampleT}}		
%\newenvironment{corollary}{\begin{cBox}\begin{corollaryT}}{\end{corollaryT}\end{cBox}}	
\newenvironment{nota}{\begin{cBox}\begin{remarqueT}}{\end{remarqueT}\end{cBox}}	

%----------------------------------------------------------------------------------------
%	REMARK ENVIRONMENT
%----------------------------------------------------------------------------------------

\newenvironment{remark}{\par\vspace{10pt}\small % Vertical white space above the remark and smaller font size
\begin{list}{}{
\leftmargin=35pt % Indentation on the left
\rightmargin=25pt}\item\ignorespaces % Indentation on the right
\makebox[-2.5pt]{\begin{tikzpicture}[overlay]
\node[draw=uMainColor!60,line width=1pt,circle,fill=uMainColor!25,font=\sffamily\bfseries,inner sep=2pt,outer sep=0pt] at (-15pt,0pt){\textcolor{uMainColor}{i}};\end{tikzpicture}} % Orange R in a circle
\advance\baselineskip -1pt}{\end{list}\vskip5pt} % Tighter line spacing and white space after remark

%----------------------------------------------------------------------------------------
%	SECTION NUMBERING IN THE MARGIN
%----------------------------------------------------------------------------------------

\makeatletter
% \renewcommand{\@seccntformat}[1]{\textcolor{uMainColor}{\csname the#1\endcsname}\quad}   

%\makeatother

%----------------------------------------------------------------------------------------
%	PART HEADINGS
%----------------------------------------------------------------------------------------

% Numbered part in the table of contents
\newcommand{\@mypartnumtocformat}[2]{%
	\setlength\fboxsep{0pt}%
	\noindent\colorbox{uMainColor!20}{\strut\parbox[c][.7cm]{\ecart}{\color{uMainColor!70}\Large\sffamily\bfseries\centering#1}}\hskip\esp\colorbox{uMainColor!40}{\strut\parbox[c][.7cm]{\linewidth-\ecart-\esp}{\Large\sffamily\centering#2}}%
}

% Unnumbered part in the table of contents
\newcommand{\@myparttocformat}[1]{%
	\setlength\fboxsep{0pt}%
	\noindent\colorbox{uMainColor!40}{\strut\parbox[c][.7cm]{\linewidth}{\Large\sffamily\centering#1}}%
}

\newlength\esp
\setlength\esp{4pt}
\newlength\ecart
\setlength\ecart{1.2cm-\esp}
\newcommand{\thepartimage}{}%
\newcommand{\partimage}[1]{\renewcommand{\thepartimage}{#1}}%
\def\@part[#1]#2{%
\ifnum \c@secnumdepth >-2\relax%
\refstepcounter{part}%
\addcontentsline{toc}{part}{\texorpdfstring{\protect\@mypartnumtocformat{\thepart}{#1}}{\partname~\thepart\ ---\ #1}}
\else%
\addcontentsline{toc}{part}{\texorpdfstring{\protect\@myparttocformat{#1}}{#1}}%
\fi%
\startcontents%
\markboth{}{}%
{\thispagestyle{empty}%
\begin{tikzpicture}[remember picture,overlay]%
\node at (current page.north west){\begin{tikzpicture}[remember picture,overlay]%	
\fill[uMainColor!20](0cm,0cm) rectangle (\paperwidth,-\paperheight);
\node[anchor=north] at (4cm,-3.25cm){\color{uMainColor!40}\fontsize{220}{100}\sffamily\bfseries\thepart}; 
\node[anchor=south east] at (\paperwidth-1cm,-\paperheight+1cm){\parbox[t][][t]{8.5cm}{
\printcontents{l}{0}{\setcounter{tocdepth}{1}}% The depth to which the Part mini table of contents displays headings; 0 for chapters only, 1 for chapters and sections and 2 for chapters, sections and subsections
}};
\node[anchor=north east] at (\paperwidth-1.5cm,-3.25cm){\parbox[t][][t]{15cm}{\strut\raggedleft\color{white}\fontsize{30}{30}\sffamily\bfseries#2}};
\end{tikzpicture}};
\end{tikzpicture}}%
\@endpart}
\def\@spart#1{%
\startcontents%
\phantomsection
{\thispagestyle{empty}%
\begin{tikzpicture}[remember picture,overlay]%
\node at (current page.north west){\begin{tikzpicture}[remember picture,overlay]%	
\fill[uMainColor!20](0cm,0cm) rectangle (\paperwidth,-\paperheight);
\node[anchor=north east] at (\paperwidth-1.5cm,-3.25cm){\parbox[t][][t]{15cm}{\strut\raggedleft\color{white}\fontsize{30}{30}\sffamily\bfseries#1}};
\end{tikzpicture}};
\end{tikzpicture}}
\addcontentsline{toc}{part}{\texorpdfstring{%
\setlength\fboxsep{0pt}%
\noindent\protect\colorbox{uMainColor!40}{\strut\protect\parbox[c][.7cm]{\linewidth}{\Large\sffamily\protect\centering #1\quad\mbox{}}}}{#1}}%
\@endpart}
\def\@endpart{\vfil\newpage
\if@twoside
\if@openright
\null
\thispagestyle{empty}%
\newpage
\fi
\fi
\if@tempswa
\twocolumn
\fi}

%----------------------------------------------------------------------------------------
%	CHAPTER HEADINGS
%-------------------------------------------------%--------------------------------------

%
%
%% A switch to conditionally include a picture, implemented by Christian Hupfer
\newif\ifusechapterimage
\usechapterimagetrue
\newcommand{\thechapterimage}{}%
\newcommand{\chapterimage}[1]{\ifusechapterimage\renewcommand{\thechapterimage}{#1}\fi}%
\newcommand{\autodot}{.}

\def\@makechapterhead#1{%
{\parindent \z@ \raggedright \normalfont
\ifnum \c@secnumdepth >\m@ne
\if@mainmatter
\begin{tikzpicture}[remember picture,overlay]
	
\node at (current page.north west)
			{\begin{tikzpicture}[remember picture,overlay]
			\node[anchor=north west,inner sep=0pt] at (0,0) {\ifusechapterimage\includegraphics[width=\paperwidth]{\thechapterimage}\fi};
			\draw[anchor=west] (\Gm@lmargin,-9cm) node [line width=2pt,rounded corners=15pt,draw=uMainColor,fill=white,fill opacity=1,inner sep=15pt]{\strut\makebox[22cm]{}};
			
			\draw[anchor=west] (\Gm@lmargin+.3cm,-9cm) node {\fontsize{20}{30}\selectfont\sffamily\bfseries\color{uMainColor}~#1\strut};
			%{\Huge\sffamily\bfseries\color{uMainColor}\thechapter\autodot~#1\strut};
			\node at (current page.north east)  [xshift=-3.95cm, yshift=-4.05cm, text opacity=1]  {{\color{white}\centering\fontsize{200}{30}\selectfont \bfseries\sffamily\thechapter\strut} };
			\node at (current page.north east)  [xshift=-4cm, yshift=-4cm, text opacity=1]  {{\color{uMainColor}\centering\fontsize{200}{30}\selectfont \bfseries\sffamily\thechapter\strut} };
			
			\end{tikzpicture}};
	\end{tikzpicture}

\else

\begin{tikzpicture}[remember picture,overlay]
\node at (current page.north west)
{\begin{tikzpicture}[remember picture,overlay]
\node[anchor=north west,inner sep=0pt] at (0,0) {\ifusechapterimage\includegraphics[width=\paperwidth]{\thechapterimage}\fi};
\draw[anchor=west] (\Gm@lmargin,-9cm) node [line width=2pt,rounded corners=15pt,draw=uMainColor,fill=white,fill opacity=1,inner sep=15pt]{\strut\makebox[22cm]{}};

\draw[anchor=east] (\Gm@lmargin,-9cm) node [line width=2pt,rounded corners=15pt,draw=uMainColor,fill=white,fill opacity=1,inner sep=15pt]{\strut\makebox[22cm]{}};
\draw[anchor=west] (\Gm@lmargin+.3cm,-9cm) node {\huge\sffamily\bfseries\color{uMainColor}#1\strut};

\end{tikzpicture}};
 
\end{tikzpicture}
\fi\fi\par\vspace*{270\p@}}}

\makeatother
%-------------------------------------------

%\def\@makeschapterhead#1{%
%\begin{tikzpicture}[remember picture,overlay]
%\node at (current page.north west)
%{\begin{tikzpicture}[remember picture,overlay]
%\node[anchor=north west,inner sep=0pt] at (0,0) {\ifusechapterimage\includegraphics[width=\paperwidth]{\thechapterimage}\fi};
%\draw[anchor=west] (\Gm@lmargin,-9cm) node [line width=2pt,rounded corners=15pt,draw=uMainColor,fill=white, fill opacity=0.8,inner sep=15pt]{\strut\makebox[22cm]{}};
%\draw[anchor=west] (\Gm@lmargin+.3cm,-9cm) node {\huge\sffamily\bfseries\color{uMainColor}#1\strut};
%\end{tikzpicture}};
%\end{tikzpicture}
%\par\vspace*{270\p@}}




%----------------------------------------------------------------------------------------
%	LINKS
%----------------------------------------------------------------------------------------


\usepackage{bookmark}
\bookmarksetup{
open,
numbered,
addtohook={%
\ifnum\bookmarkget{level}=0 % chapter
\bookmarksetup{bold}%
\fi
\ifnum\bookmarkget{level}=-1 % part
\bookmarksetup{color=uMainColor,bold}%
\fi
}
}

%----------------------------------------------------------------------------------------
%	FANCY CHAPTER Header
%----------------------------------------------------------------------------------------





%\definecolor{uMainColor}{RGB}{255,128,0}   



%-------------------------------------------------
%	BIBLIOGRAPHY
%-------------------------------------------------

\addbibresource{../Latex/Sources/commons/bibliography.bib} % BibTeX bibliography file
\defbibheading{bibempty}{}

%-------------------------------------------------
%	COLORS & BORDERS
%-------------------------------------------------
\setlength{\fboxrule}{0.75pt} % Width of the border around the abstract
%\definecolor{uMainColor}{RGB}{255,128,0} % Color of the article title and section
\definecolor{color2}{RGB}{220,220,220} % Color of the boxes behind the abstract and headings

%-------------------------------------------------
% ITEMS DEFINITION
%-------------------------------------------------

\setlist [itemize,1]{label=\color{uMainColor}\faCaretRight }


\usepackage{environ}
\newwrite\audioexport



\begin{document}





  %-------------------------------------------------------
  % AudioNote Open File
  %-------------------------------------------------------
\mode<all>{\immediate\openout\audioexport=\jobname.txt}


\addtocontents{toc}{\protect\hypertarget{toc}{}}

%-------------------------------------------------
%	ARTICLE INFORMATION
%-------------------------------------------------

\Abstract {\uabstract%===========================
% COURS "INTRO CYBERDEF"
% Abstract général des articles
%===========================


Il fait partie du cours introductif aux fondamentaux de la sécurité des systèmes d'information vue sous deux prismes quelques fois opposés dans la littérature : la gouvernance et la gestion opérationnelle de la sécurité.
Le cours est constitué d'un ensemble de notes de synthèse indépendantes compilées en un document unique, mais édité par chapitre dans le cadre de ce cours.\\
C\edoc ne constitue pas à lui seul le référentiel du cours CYBERDEF101 (SEC101 du Cnam). Il compile des notes de cours mises à disposition de l'auditeur comme support pédagogique partiel à ce cours introductif à la cyberdéfense d'entreprise. }
\JournalInfo{\uJournalInfo} % Journal information
\Archive{Notes de cours éditées le  \DTMnow} % Additional notes (e.g. copyright, DOI, review/research article)
\PaperTitle{\utitle} % Article title
\Authors{\uauthor\textsuperscript{1,}\textsuperscript{2}*} % Authors
\affiliation{\textsuperscript{1}\textit{\uproa}} % Author affiliation
\affiliation{\textsuperscript{2}\textit{\uprob}} % Author affiliation
\affiliation{*\textbf{email}: \umaila\ -- \umailb} % Corresponding author
\Keywords{\ukeywords} % Keywords 
\newcommand{\keywordname}{Mots clefs} % Defines the keywords heading name

%-------------------------------------------------
%	MAKE TITLE
%-------------------------------------------------

\maketitle

\begin{tikzpicture}[remember picture,overlay]\node at (current page.north west)[xshift=3cm, yshift=-3cm, text opacity=1.0]{\includegraphics[width=0.1\paperwidth]{\utemplatepath/commons.inc/CommonsPictures/\utemplatetheme.theme/shield-20.pdf}}; 
\end{tikzpicture} 


\AddToShipoutPicture{%
\AtPageCenter{%
\begin{tikzpicture}[remember picture,overlay]
\node [rotate=60,scale=5,text opacity=0.1] at (current page.center) { \textit{\color{uAlertColor}\INFODistrib} };
\end{tikzpicture} 
%\begin{tikzpicture}[remember picture,overlay]
%\node [rotate=60,scale=10,text opacity=0.3]
%at (current page.center) {moi};
%\end{tikzpicture}
}%
}


%-------------------------------------------------	
\begin{center}
	\textcolor{uMainColor}{\textbf{\INFODistrib}} % Distribution information
\end{center}

\INFOGitHub % Link to download file and copyright
%-------------------------------------------------

%-------------------------------------------------
	\newpage
	\tableofcontents  % TOC
	\iftotalfigures
  			\listoffigures
	\fi
%-------------------------------------------------

%-------------------------------------------------
\newpage
\ubody % MAIN BODY defined external
%-------------------------------------------------

	\printbibliography
 \printglossary[type=\acronymtype]
%	\printglossary
%-------------------------------------------------	
\newpage
\Ucontribute % Contribution TEXT
%-------------------------------------------------

  %-------------------------------------------------------
  % AudioNote Open File
  %-------------------------------------------------------
  \mode<all>{\immediate\closeout\audioexport}

\end{document}
}
	\fi
\fi


 

%************************************************************
% Chargement des variables dédiées à l'article
%************************************************************
\lessonID{7.2}
\setauthor{EDU}

%========================================
\newcommand {\ukeywords}
%========================================
{%---------------------------------------------------------------------
Anticipation, veille, alerte, réponse, CERT, SOC
}%

%========================================
\newcommand {\utitle}
{%---------------------------------------------------------------------
Threat Management
}%---------------------------------------------------------------------

%========================================
\newcommand {\uabstract}
{%---------------------------------------------------------------------
C\edoc introduit la gestion de la menace de la partie cyberdéfense de la sécurité opérationnelle : Anticiper, Détecter, Réagir  et ceci sur les trois grands invariants des risques numériques : les vulnérabilités, les menaces et l'impact.
}
%---------------------------------------------------------------------

%************************************************************
%  variable définissant  le corps de l'article
%************************************************************

%========================================
\newcommand {\ubody}
{%---------------------------------------------------------------------
%$BEGIN_UBODY

% THREAT MAN

\mode<all>{%-------------------------------------------------------------
%               FR CYBERDEF SECOPS COURSE
%              $Chapitre : Threat Management
%                         $theme : Introduction
%                 $File : chap-Vulman-intro.tex
%                             2020 eduf@ction
%-------------------------------------------------------------
\uchap{chap-Vulman-intro.tex}
%-------------------------------------------------------------


\section{GERER les menaces}

Je vous propose d'aborder ce chapitre lié à la surveillance de l'évènement de sécurité avec les quelques points fondamentaux de la gestion de la menace.
En effet nous partons du principe qu'au sein de périmètre de surveillance donné, la gestion de la menace peut s'organiser autour de la gestion des évènements à risques détectés dans ce périmètre.

\mode<all>{\picframe{Latex/Sources/EDU/Pictures/img-cycle}{Cycle de vie de gouvernance Cyberdef}{0.8}{lblcycle1}}

Après quelques définitions et positionnement dans l'analyse de menace,  de la supervision et de l'analyse comportementale nous aborderons donc les grandes fonctions nécessaire à la \textbf{détection} comme :

% Planche PRZ =========================
\mode<all>{\texframe{Déroulement}{présentation Gestion de la menace}
{
\begin{itemize}
  \item \textbf{VOIR} :  capacité de voir et de capter le comportement d'un système d'information via des sources et capteurs avec le \textit{\uindex{LOG management}} (Systèmes et Applicatifs). En n'oubliant pas d'évoquer l'assurance sécurité des Logs (intégrité, \uindex{horodatage}, valeur probante ...)
\item \textbf{COMPRENDRE - PREVOIR} : Avec le \textit{\uindex{Threat Management}} : Veiller, surveiller la menace dans l'environnement digital de l'entreprises, modélisation de la menace et scénarios redoutés issus d'analyse de risque;
\item \textbf{DETECTER} : Surveiller le comportement des systèmes dans le périmètre défini, faire émerger les  évènements, anomalies, incidents pouvant révéler une attaque en cours, une suspicion de compromission par des menaces avancées (APT), où des attaques furtives et discrètes. Nous aborderons l'outillage avec les \uindex{SIEM} et l'organisation avec les \uindex{SOC};
\item \textbf{ALERTER} : mettre en place les mécanismes de remontée d'alerte et d'incident permettant de gérer les alertes adaptées au niveau d'impact d'une attaque. 
\end{itemize}
}}



% Planche PRZ ===========================
\mode<all>{\texframe{Menaces}{la veille}
{
\begin{description}
	\item \textbf{Menaces=Veille et recherche} : La gestion de la menace est au coeur des stratégies de cyberdéfense de l'entreprise. Comme pour les vulnérabilités, c'est la connaissance des menaces, de leur recherche et de leur découverte qui permet de réduire les risques;
\item  \textbf{Menaces=Évènements} : La détection d'une vulnérabilité ou d'une menace est un évènement, la question est de savoir à quel moment il est important de déclencher un mécanisme d'alerte, et comment cette alerte va devenir un incident déclenchant des mécanismes de réponse \uref{Cycle de gouvernance}{lblcycle1}.
\end{description}
}}
% end PRZ ===========================


%******************************************
\subsection{Modèles}

Ce sont généralement les attaques externes et massives qui font l’actualité dans les médias. 
Un grand nombre de risques cyber quotidiens sont issus de l'intérieur même de l'entreprise.
Des  fuites de la part d'employés qui, de façon intentionnelle ou non, révèlent des mots de passe ou des informations sensibles. Une opération initiée par des acteurs internes malveillants (salariés, partenaires, clients) qui peuvent utiliser les informations à leur portée afin d’exploiter ou de porter dommage aux systèmes d'information de l’entreprise, mais plus globalement à l'entreprise dans sa globalité.

\mode<all>{\picframe{Latex/Sources/EDU/Pictures/img-var}{Un modèle de gestion cyberdéfense}{0.7}{lblthreatman}}

N’importe quelle entreprise est exposée au risque. De nos jours, les connexions diverses d’une entreprise représentent de nombreuses voies pour des attaques informatiques, qui ciblent souvent les petites entreprises afin d’accéder à des plus grands acteurs (partenaires,  clients ou fournisseurs). Les grandes entreprises imposent de plus en plus  souvent à leurs fournisseurs et partenaires, quelle que soit leur taille, de mettre en place des mesures de cybersécurité. Ceci est généralement structuré autour de plan d'assurance sécurité, annexé au contrat.

Les attaquants externes sont, certes, des menaces croissantes car elles recherchent sans cesse des failles de sécurité afin d’accéder à vos systèmes, mais les menaces internes ne sont  pas à négliger sur le plan opérationnel. On verra d'ailleurs que pour couvrir ces menaces, la détection des évènements liés à des comportements de ses propres usagers ou salariés n'est pas chose facile (Législation sur le droit des correspondance, \gls{aCNIL} ...).

Il est important toutefois de ne pas distinguer des dynamiques internes et externes en terme d'agression cyber, c'est pour cela que les modèles de Cyberdéfense ne formalisent que très peu cette notion de provenance de l'attaque.

En outre, sur le triptyque \g{Vulnérabilités, Menaces, Incidents} \uref{Modèle de cyberdéfense}{lblthreatman}. la notion d'attaque interne et/ou externe n'a de sens qu'au titre de la responsabilité du périmètre, car rien n'interdit de positionner des capteurs en dehors de son périmètre technique pour anticiper la menace. Ainsi, la veille sur internet, le renseignement cyber sont nécessairement équipés de capteurs pouvant faire remonter des événements dans des chaînes d'alertes Cyber.
Par exemple, l'attaque d'une entreprise du même secteur peut être  en soit un incident pouvant générer une alerte.

Si on se focalise sur la gestion de la menace, il existe de nombreuses manières de présenter ce processus. Un modèle \uref{Threat Management Cycle}{lblthreatrosas} issu des travaux de l'\gls{aOTAN} est intéressant car il propose 4 axes d'analyses, toutefois Threat Mitigation et Threat Protection sont un peu ambigus.


\mode<all>{\picframe{Latex/Sources/EDU/Pictures/img-threatrosas}{les 4 axes de la gestion de la menace}{0.8}{lblthreatrosas}}

Je propose de continuer d'utiliser le modèle organisé autour de quatre volets (Surveiller, Détecter, Alerter, s'entrainer), \uref{Défendre les actifs}{lblvarmenaces} pour structurer la présentation.

%Différents types de menaces à la cybersécurité

La nature de ces menaces est en constante évolution.  On y trouve les plus courantes :

% Begin PRZ ==========================
\mode<all>{\texframe{les grandes menaces}{quelques éléments de la menace 1/2}
{
\begin{itemize}
\item Attaques par \textbf{déni de service distribuées} (DDoS).  Un réseau d’ordinateurs inonde un site Web ou un logiciel avec des informations inutiles. L'exemple, le plus classique est celui d'un serveur WEB. Quand la charge sur les services est trop importante et que le système n'est pas dimensionné ou filtré pour ce type de volume de demande, ce débordement de requêtes provoque une indisponibilité du système inopérant. 
\item \textbf{Codes malveillants} : Bots et virus. Un logiciel malveillant qui s’exécute à l'insu de l'utilisateur ou du propriétaire du système (bots), ou qui est installé par un employé qui pense avoir affaire à un fichier sain (cheval de Troie), afin de contrôler des systèmes informatiques ou de s’emparer de données.
 \end{itemize}
}}

\mode<all>{\picframe{Latex/Sources/EDU/Pictures/img-varmenaces}{la gestion active de la menace}{0.5}{lblvarmenaces}}

% ----Begin PRZ Frame------------
\mode<all>{\texframe{les grandes menaces}{quelques éléments de la menace 2/2}
{
\begin{itemize}
\item \textbf{Piratage}. Lorsque des acteurs externes exploitent des failles de sécurité afin de contrôler vos systèmes informatiques et voler des informations, en utilisant ou pas un code malveillant. Par exemple, un changement régulier des mots de passe et la mise à niveau des systèmes de sécurité est fondamentale pour limiter les impacts.
\item \textbf{Hameçonnage} ou dévoiement. Tentative d’obtenir des informations sensibles en se faisant passer frauduleusement pour une entité digne de confiance. Le hameçonnage se fait généralement par e-mail, mais il ne faut pas oublier les SMS et les services utilisant du message (Webmail, mail intégré comme Linkedin, ...), 
\end{itemize}
}}%----End Frame------------


C'est la combinaison d'actions élémentaires, d'attaques élémentaires qui font des scénarios de menaces. 

Il est important aussi de repositionner la définition des menaces par rapport à la notion d'attaque, mais aussi la notion de risques et de vulnérabilités \uref{Cycle du risque}{lblthreatwheel}.

%TODO modifier l'image ... est couvert et sont couverts / compromis

\mode<all>{\picframe{Latex/Sources/EDU/Pictures/img-threatwheel}{la gestion de la menace}{0.8}{lblthreatwheel}}

\begin{itemize}
  \item des attaques \textbf{matérialisent} des menaces,
  \item des menaces \textbf{exploitent} des vulnérabilités,
  \item des vulnérabilités \textbf{exposent} à des risques,
  \item des risques\textbf{ sont couverts }par des contre-mesures,
  \item des \textbf{contre-mesures }protègent des actifs,
  \item des actifs \textbf{sont  soumis }à  des menaces.
\end{itemize}

Nous voyons donc ici qu'il est important de ne pas séparer en terme de gouvernance et de pilotage opérationnel la sécurité la gestion des vulnérabilités, la gestion des menaces et la gestion des risques.


%******************************************


%*************************************
\subsection{Les processus de gestion de la menace}

% Begin PRZ ===========================
\mode<all>{\texframe{Gérer la menace}{Threat Intelligence et Detection}
{
Gérer la menace comporte deux donc domaines d'activités :
\begin{itemize}
  \item La veille, au sens renseignement sur la menace (Threat Intelligence);
  \item La détection d'attaque, ou de menaces potentielles au sein de l'environnement (Threat Detection).
\end{itemize}
}}% end PRZ ===========================

\mode<all>{\picframe{Latex/Sources/EDU/Pictures/img-logman}{la gestion de la menace}{0.8}{lbllogman}}

Ces deux domaines d'activités se basent sur la remontée d'information et l'automatisation des détections d'évènements à risques.  Pour automatiser la détection d'évènement, il donc imprimant de disposer de sources d'information de type \g{évènements} que des techniques historiques informatiques et réseaux  apportent grâce au \textit{LOG  Management} (gestion des traces et journaux). Les journaux informatiques des systèmes sont au coeur de la détection, mais il existe de nombreuses autres sources d'évènements (informations,  renseignements) formalisées qui peuvent apporter de l'information pertinente pour la détection ou l'anticipation d'attaques.

\subsection{Détecter, la surveillance du SI}

% FIN DE RELECTURE

\g{Détecter oui, mais détecter quoi et pourquoi} est la phrase maîtresse de la première étape de réflexion autour de la gestion de la détection d'incident de sécurité.
La première question à se poser est  de définir ce qu'est un incident de sécurité pour l'entreprise. S'il est vrai qu'il existe un certain nombre de menace \g{standard} que l'on considère très rapidement comme un incident, le déploiement d'outil de gestion d'incident de sécurité ne serait être limité qu'à cet usage standard.

Il y a de nombreuses manières de détecter des tentatives d'attaques dans un système. Les IPS/IDS (\textit{Intrusion Prevention System / Intrusion Detection System}), les Firewall réseaux et firewall applicatif, ...Toutefois l'imagination des attaquants est suffisamment grande, pour que des attaques complexes ne puisse être détectées par ces seuls outils et produits de sécurité protégeant les flux informationnels.

Nous pouvons en effet considérer par exemple que la détection d'un rançon-logiciel dans l'entreprise est un incident complexe, qu'un IPS/IDS ne détectera pas,  qui va par ailleurs nécessiter une alerte et une remédiation rapide si ce n'est immédiate.
Par ailleurs,  une fuite d'information sur un système métier par des mécanismes discrets sera souvent étudiée spécifiquement.
Globalement le déploiement d'une fonction d'alerte va nécessiter la définition des \g{menaces} redoutées par l'entreprise. 
Ces dernières sont généralement issues des analyses de risques.
En effet, il est important au delà des menaces dites standards de revenir aux origines du déploiement des fonctions de sécurité qui sont de gérer et couvrir les risques.

En premier lieu, il convient de chercher à détecter les menaces non couvertes par les mesures de sécurité, les fameuses menaces résiduelles.

Dans l'environnement de l'entreprise, les scénarios complexes issus de l'analyse de risques lors de l'étude des évènements redoutés vont donner les évènements corrélés à détecter.
On y trouvera l'application concrète des arbres d'attaques popularisés par un des plus célèbre cyber expert \wikipedia{Bruce Schneier} \ucite{schneier1999attack} qui est présentée de manière succincte  dans le  chapitre \ref{sec:arbreattaques}

\subsection{Attaques}

Le cycle de vie d’une cyberattaque, qu’elle soit complexe ou sophistiquée, reste le même depuis des années. Elle se déroule en 3 étapes :
\begin{itemize}
  \item la première : \textbf{la phase de reconnaissance}. Elle va permettre d’identifier sa cible et de rechercher l’ensemble des vulnérabilités. Contrairement à l’audit de sécurité, dans cette étape, l’assaillant n’a aucune contrainte de périmètre ni de cadre contractuel. Cette absence de contrainte va lui permettre d’exploiter tout type de vulnérabilité;
  \item la seconde : l’\textbf{attaque elle-même}. Les attaques, aujourd’hui plus sophistiquées, permettent aux assaillants, une fois entrés dans le système d’information de l’entreprise, d’effectuer diverses actions comme l’élévation de privilèges, la création d’une porte dérobée, la mise en sommeil des agents dormants et surtout l’effacement de toute trace de son passage;
  \item  la dernière : l’\g{atteinte de son objectif et l'action}. Cela peut se traduire par une simple perturbation des systèmes ou encore l’exfiltration d‘informations sensibles dans le but d’actes de manipulation.
\end{itemize}

Ces mécanismes sont souvent modélisés sous la forme d'arbres d'attaques ou de scénarios issus d'analyses de risques.

\subsubsection{Arbre d'attaques} 

\label{sec:arbreattaques}

%TODO a re writer

Détecter la menace dans un système d'information c'est aussi connaître les méthodes, stratégies des attaquants. Ces scénarios d'attaque ou d'opération peuvent être modélisés avec des outils au coeur des analyses de risques. Bien que très largement en arrière plan des méthodes et des outils de gestion de la menace, les arbres d'attaque restent au coeur des mécanismes de détection.

Les arbres d'attaques sont une représentation des scénarios d'attaques. La racine représente le but final de l'attaque, les différents noeuds sont les buts intermédiaires et les feuilles les actions élémentaires à effectuer. Ces actions seront évaluées par exemple avec les potentiels d'attaque des \textbf{critères communs} (Ensemble de normes \wikipedia{ISO 15408})

Globalement, ces arbres sont basés sur trois types de nœuds :

\begin{itemize}
  \item Nœud \textbf{disjonctif} OR : OU logique. Cela signifie que pour que le nœud soit réalisé, il faut qu’au moins un de ses fils soit réalisé.
  \item Nœud \textbf{conjonctif} AND : ET logique. Pour sa réalisation, il faut que l’ensemble de ses fils soit réalisé.
  \item Nœud \textbf{conjonctif séquentiel }SAND : Pour sa réalisation, il faut que l'ensemble de ses fils soit réalisé dans un ordre séquentiel c'est-à-dire les fils sont effectués les uns après les autres dans l’ordre indiqué.
\end{itemize}

En fonction de ces noeuds les valeurs des feuilles seront remontées pour obtenir le potentiel d'attaque de la racine.
C'est sur la base de ce type de technique que sont construits un certain nombre d'outils de détection.
 
%Vers ADTool http://satoss.uni.lu/members/piotr/adtool/
%ADTool est un logiciel qui permet la création, l'édition et l'affichage de ces arbres. Il permet de valuer les feuilles et faire remonter les valeurs à la racine selon différents algorithmes

\wikipedia{Gartner}, et \wikipedia{Lockheed Martin} ont dérivé le concept de ces arbres d'attaque dans des modèles dit de \g{Kill Chain} issus de modèles militaires établis à l’origine pour identifier la cible, préparer l’attaque, engager l’objectif et le détruire.

Ce modèle analyse une fragilité potentielle en dépistant les phases de l’attaque, de la reconnaissance précoce à l’exfiltration des données. Ce modèle de chaîne ou processus cybercriminel aide à comprendre et à lutter contre les ransomware, les failles de sécurité et les menaces persistantes avancées (APT). Le modèle a évolué pour mieux anticiper et reconnaître les menaces internes, l’ingénierie sociale, les ransomware avancés et les nouvelles attaques.

\subsubsection{Le déploiement d'une menace en 8 étapes}

%TODO  à revoir le texte

\begin{itemize}
  \item \textbf{Phase 1 : Reconnaissance}. Comme dans un « casse » classique, vous devez d’abord repérer les lieux. Le même principe s’applique dans un cyber-casse : c’est la phase préliminaire d’une attaque, la mission de recueil d’informations. Pendant la reconnaissance, le cybercriminel recherche les indications susceptibles de révéler les vulnérabilités et les points faibles du système. Les pare-feu, les dispositifs de prévention des intrusions, les périmètres de sécurité (et même les comptes de médias sociaux) font l’objet de reconnaissance et d’examen. Les outils de repérage analysent les réseaux des entreprises pour y trouver des points d’entrée et des vulnérabilités à exploiter.

 \item \textbf{Phase 2 : Intrusion.} Après avoir obtenu les renseignements, il est temps de s’infiltrer. L’intrusion constitue le moment où l’attaque devient active : les malwares (y compris les ransomwares, spywares et adwares) peuvent être envoyés vers le système pour forcer l’entrée. C’est la phase de livraison. Celle-ci peut s’effectuer par e-mail de phishing ou prendre la forme d’un site Web compromis ou encore venir du sympathique café au coin de la rue avec sa liaison WiFi, favorable aux pirates. L’intrusion constitue le point d’entrée d’une attaque, le moment où les agresseurs pénètrent dans la place.

 \item \textbf{Phase 3 : Exploitation}. Le hacker se trouve de l’autre côté de la porte et le périmètre est violé. La phase d’exploitation d’une attaque profite des failles du système, à défaut d’un meilleur terme. Les cybercriminels peuvent désormais entrer dans le système, installer des outils supplémentaires, modifier les certificats de sécurité et créer de nouveaux scripts à des fins nuisibles.

 \item \textbf{Phase 4 : Escalade de privilèges}. Quel intérêt y a-t-il à entrer dans un bâtiment si vous restez coincé dans le hall d’accueil ? Les cybercriminels utilisent l’escalade de privilèges pour obtenir des autorisations élevées d’accès aux ressources. Ils modifient les paramètres de sécurité des GPO, les fichiers de configuration, les permissions et essaient d’extraire des informations d’identification.

 \item \textbf{Phase 5 : Mouvement latéral}. Vous avez carte blanche, mais vous devez encore trouver la chambre forte. Les cybercriminels se déplacent de système en système, de manière latérale, afin d’obtenir d’autres accès et de trouver plus de ressources. C’est également une mission avancée d’exploration des données au cours de laquelle les cybercriminels recherchent des données critiques et des informations sensibles, des accès administrateur et des serveurs de messagerie. Ils utilisent souvent les mêmes ressources que le service informatique, tirent parti d’outils intégrés tels que PowerShell et se positionnent de manière à causer le plus de dégâts possibles.

 \item \textbf{Phase 6 :  Furtivité, camouflage, masquage}. Mettez les caméras de sécurité en boucle et montrez un ascenseur vide pour que personne ne voit ce qui se produit en coulisses. Les cyber-attaquants font la même chose. Ils masquent leur présence et leur activité pour éviter toute détection et déjouer les investigations. Cela peut prendre la forme de fichiers et de métadonnées effacés, de données écrasées au moyen de fausses valeurs d’horodatage (time-stamping) et d’informations trompeuses, ou encore d’informations critiques modifiées pour que les données semblent ne jamais avoir été touchées.

 \item \textbf{Phase 7 :  Isolation et Déni de service}. Bloquez les lignes téléphoniques et coupez le courant. C’est là où les cybercriminels ciblent le réseau et l’infrastructure de données pour que les utilisateurs légitimes ne puissent obtenir ce dont ils ont besoin. L’attaque par déni de service (DoS) perturbe et interrompt les accès. Elle peut entraîner la panne des systèmes et saturer les services.

 \item \textbf{Phase 8 :  Exfiltration}. Prévoyez toujours une stratégie de sortie. Les cybercriminels obtiennent les données. Ils copient, transfèrent ou déplacent les données sensibles vers un emplacement sous leur contrôle où ils pourront en faire ce qu’ils veulent : les rendre contre une rançon, les vendre sur eBay ou les envoyer à BuzzFeed. Sortir toutes les données peut prendre des jours entiers, mais une fois qu’elles se trouvent à l’extérieur, elles sont sous leur contrôle.

\end{itemize}

Différentes techniques de sécurité proposent différentes approches de la chaîne cyber-criminelle. De Gartner à Lockheed Martin, chacun définit les phases de manière légèrement différente.

C’est un modèle quelque peu critiqué pour l’attention qu’il accorde à la sécurité  périmètrique et focalisé sur la prévention des malwares. Cependant, quand elle est combinée à l’analyse avancée et à la modélisation prédictive, la chaîne cyber-criminelle devient essentielle à une sécurité complète.

\subsubsection{UBA : User Behavior Analytics}

% SOURCE : https://www.silicon.fr/avis-expert/uba-analyser-le-comportement-des-utilisateurs-pour-contrer-le-surmenage-des-equipes-it-et-les-cybermenaces

L’analyse du comportement des utilisateurs (UBA) apporte des informations détaillées sur les menaces liées à chaque phase de la chaîne criminelle. Et elle contribue à prévenir et arrêter les attaques avant que les dommages ne soient causés. En effet, le volume d’activités suspectes, dont des faux positifs, inhérents aux outils de sécurité traditionnels sont très chronophages à surveiller et donc sources d’erreurs.
Pour y pallier, les entreprises doivent être en mesure d’analyser le comportement des utilisateurs, souvent via des outils d'apprentissage automatique, afin de donner un sens aux informations remontées par la lecture des activités sur le réseau.

Cette analyse du comportement des utilisateurs (UBA : User Behavior Analytics) aide à  comprendre et hiérarchiser les alertes filtrant celles qui sont suspectes en comparaison avec des comportements habituel des utilisateurs.

\subsubsection{La surveillance des terminaux}

La surveillance des terminaux (EndPoint) est aujourd'hui un point important dans le prise en charge de la menace du côté l'utilisateur (mais aussi serveurs)

 Le  terme \g{Endpoint (Threat) Detection and Response} (EDTR ou EDR). Un système EDR met  l’accent sur la détection d’activités suspectes directement sur les hôtes de traitement du système d'information au delà de l'infrastructure. 

\subsubsection{La surveillance de l'IOT}

%https://medium.com/@dtembe/iot-and-siem-integration-pt-1-6645a012bdc

La surveillance dans l'internet des objets est en plein structuration. Les technologies s'intègrent dans des schémas d'API complexes qui nécessitent d'intégrer les sondes dans des architectures nécessitant de gérer contractuellement les périmètres de sécurité de surveillance. Par ailleurs, La diversité des protocoles utilisés sur les réseaux industriels ou spécifiques aux différentes technologies des IOT nécessitent pour les éditeurs de solutions de maitriser un large champs d'attaques.

\subsection{Gestion de la menace}

Nous avons évoqué dans le chapitre sur l'anticipation, la veille sur la menace. 
Opérer la détection d'attaques ou de menaces dormantes dans l'environnement de l'entreprise nécessite une connaissance précise des mécanismes d'exécution ou d'opération de ces menaces.
La connaissance de ces mécanismes d'action, de protection, de déploiement, de réplication, de survivabilité, de déplacement des codes malveillants par exemple est la base de leur détection. Il est en de même sur les scénarios mixant des actions sur les réseaux ou sur les systèmes informatiques ou numériques.
Ces connaissances sont généralement structurées dans des bases de connaissances dont les sources sont gratuites ou payantes.

\subsection{Bases de connaissance et menaces}

\subsubsection{Sources identifiées menaçantes}

% Begin PRZ ===========================
\mode<all>{\texframe{Threat Intelligence}{Sources identifiées menaçantes}
{
Nous parlerons ici de sources de menaces comme les indicateurs permettant d'identifier l'origine technique d'une menace. Cela peut être une adresse mail, un serveur/service de mail , une adresse IP de provenance d'un code malveillant, d'une attaque, ou d'un comportement anormal.
On peut citer par exemple  :
\begin{itemize}
  \item Une adresse mail connue pour envoyer des codes malveillants;
  \item des adresses IP ou des adresses de serveur Mail pour Spam.
\end{itemize}
}}% end PRZ ===========================


En face, il y a des attaquants qui bien entendu vont changer leur position pour émettre ou attaquer d'ailleurs, ou avec une autre forme (furtivité). Ces bases d'informations peuvent donc devenir rapidement obsolètes. Ceci dénote l'importance de disposer de base de connaissance sur les sources de menaces à jour et en temps réel.

\subsubsection{Cibles de menaces}

Les cibles de menace peuvent être connues à un instant T. Ces cibles peuvent être  sectorielles (Banques, sites étatiques ...).


\subsubsection{Threat Intelligence Database }

Dans la notion de partage de l'information sur la menace, le \ulink{https://www.misp-project.org}{projet MISP} (Open Standards For Threat Information Sharing) fournit les modèles de données et des indicateurs.

Il ne faut pas oublier que les données dans ces bases peuvent être \g{éphémères} (Adresse IP malveillante, machines infectées nettoyées...). Il est donc important de disposer de sources fiables et mise à jour en temps réel.

\begin{techworkbox}{Bases de Threat Intelligence}
Le marché des bases publiques et commerciales de \textit{Threat Intelligence} est un sujet intéressant pour \fichetech. On peut rechercher les sources les plus pertinentes, celles qui fournissent des informations techniques pour les SIEM, celles spécialisées dans des menaces sectorielles ou technologiques ...
\end{techworkbox}


\subsubsection{Exemple de TI}

On peut donc trouver dans ces base de threat intelligence des données du type suivant  :

% EXCEL 2 TAB https://www.tablesgenerator.com/latex_tables
% Please add the following required packages to your document preamble:
% \usepackage{graphicx}
\begin{table}[!hbtp]
	\resizebox{\textwidth}{!}{
	\begin{tabular}{ l|l l}
		\toprule
		\textbf{Nom }                 & \textbf{Command\&Control URL }                                        & \textbf{HASH : sha256} \\ \hline
		caracal.raceinspace.astronaut & http://api.lulquid.xyz                       & f1d32c17a169574369088...\\ \hline
		com.caracal.cooking           & http://api.namekitchen9.xyz/api/subscription & 46e41ef7673e34ef72fb3a9718...  \\ \hline
		com.leo.letmego               & http://api.leopardus.xyz/api/subscription    & b21cb5ebfb692a2db1c5cbbc20e00d90a... \\ \hline
		com.caculator.biscuitent      & http://api.lulquid.xyz                       & 734418efafd312e9b3e96adaac6f86cc1a... \\ \hline
		com.pantanal.aquawar          & http://api.pantanal.xyz                      & 8fec77c47421222cc754b32c60794e...\\ \hline
		com.pantanal.dressup          & http://api.pantanal.xyz                      & 64e2c905bcef400e861469e114bf...\\ \hline
		inferno.me.translator         & http://api.molatecta.icu                     & ebe3546208fd32d3f6a9e5daf21a7240...\\ \hline
		translate.travel.map          & http://api.nhudomainuong.xyz                 & f805e128b9d686170f51b1add35...  \\ \hline
		travel.withu.translate        & http://api.molatecta.icu                     & b7670b5d9a6643a54b800b4c...  \\ \hline
		allday.a24h.translate         & http://api.royalchowstudio.xyz               & 29f2fd6ccf0f632e45dd1f15ec72985... \\ \hline
		banz.stickman.runner.parkour  & http://api.lulquid.xyz/api/                  & e1027b6681e93d9763f19ea7e5ab2...  \\ 
		\bottomrule
	\end{tabular}
	}
\caption{Exemple de  données de TI (Threat Intelligence)}
\label{tab:ExempleTI}
\end{table}

avec le nom du package malveillant, l'url du centre de commande et de contrôle, et le Hash (empreinte) permettant de rechercher le fichier dans les données d'un système.

%TODO Ajouter : THREAT TELEGRAMME 

}
\mode<all>{%-------------------------------------
% Chapitre
% Threat Management
% Intro
% File : chap-ThreatMan-history.tex
%--------------------------------------
% sources 
% https://channellife.co.nz/story/a-brief-history-of-cyber-threats-from-2000-to-2020
\uchap{Menaces et histoires}
%_-_-_-_-_-_-_-_-_-_-_-_-_-_-_-_-_-_-_-_-_-_-_-_-_-_-_


\subsection{Quelques éléments d'histoire des menaces d'aujourd'hui}

« Pour prévoir l'avenir, il faut connaître le passé, car les événements de ce monde ont en tout temps des liens aux temps qui les ont précédés. Créés par les hommes animés des mêmes passions, ces événements doivent nécessairement avoir les mêmes résultats. » disait Nicolas Machiavel.

 Les événements passés peuvent éclairer les tendances futures, selon la sagesse courante - et la cybersécurité ne fait pas exception.
Les rapports annuels sur les menaces offrent aux équipes de sécurité l'occasion de réfléchir aux cyber-événements importants des 12 derniers mois, dans le but d'identifier les tendances pour le développement futur, idéalement se traduisant par une meilleure protection. Mais il peut être interessant de replonger un peu dans l'histoire.
De nombreux événements de cybersécurité importants se sont produits depuis le fameux passage de l’an 2000, il y a donc plus de 20 ans.

\subsubsection{Les années 2000 - L'ère du développement des vers}
Cette époque a vu certains des vers les plus prolifiques que l'industrie de la sécurité de l'information ait jamais connus, coûtant plus de 100 milliards de dollars en dommages et en coûts de réparation. Cela marque également le début du malware en tant que sensation médiatique grand public.

On peut noter  en particulier le ver ILOVEYOU, lancé en 2000, qui ciblait les utilisateurs de Microsoft Outlook et infectait au moins 10\% des hôtes connectés à Internet en quelques heures et causait alors jusqu'à 15 milliards de dollars de dommages.

En réponse, Microsoft a publié une mise à jour d'Outlook avec des modifications visant à lutter contre les pires symptômes de ILOVEYOU, notamment en empêchant les utilisateurs d'accéder à des pièces jointes dangereuses et en avertissant les utilisateurs si un programme tentait d'envoyer du courrier en leur nom.

Puis ce fut une véritable vague de vers dans les cinq années suivantes, exploitant les vulnérabilités du système d'exploitation Windows en particulier et  de l'infrastructure réseau. On peut noter en particulier

\begin{itemize}
  \item CodeRed (juillet 2001)
  \item Code Red II (août 2001)
  \item Nimda (septembre 2001)
  \item SQL Slammer (janvier 2003)
  \item Blaster (août 2003)
  \item Welchia (août 2003)
  \item Sobig.F (août 2003)
  \item Sober (octobre 2003)
  \item Bagle (janvier 2004)
  \item MyDoom (janvier 2004)
  \item Netsky (février 2004)
  \item Sasser (avril 2004)
\end{itemize}


Beaucoup de ces vers ont fortement utilisé des vulnérabilités de type "débordement de tampon/pile" dans diverses versions de Windows ou dans des applications telles que Internet Information Services (IIS) ou SQL Server les applications phares de Microsoft dans les années 2000.

La prolifération de ces codes malveillants de type vers ont poussé les éditeurs à mettre en place des mécanismes de corrections organisés. Microsoft a par exemple lancé le "Patch Tuesday", une approche structurée et cohérente de la distribution des correctifs corrigeant la distribution des correctifs ad hoc ou dans le cadre de service packs.

Bien structuré autour des firewall, le filtrage périmètrique évolue aussi vers la vérification des e-mails qui sont sous 

Le filtrage des e-mails s'est également amélioré à cette époque, les fournisseurs utilisant davantage de logiciels de détection pour leurs produits. Mais alors que les menaces par e-mail continuaient leurs attaques malgré tout, de nombreux cybercriminels passaient à leur prochain programme: gagner de l'argent.

\subsubsection{2005-2012 - L'ère du business malveillant}
Les cybercriminels, qui jusque-là provoquaient des perturbations principalement pour la notoriété ou la curiosité, ont commencé à penser à leur passe-temps comme une source de revenu. La cybercriminalité, soutenue par un internet portant la digitalisation de échange marchant, a trouvé aussi dans ces nouveaux "chimistes" une capacité à faire fructifier dans le monde numérique leurs techniques et méthodes criminelles classiques. C'est le début de la digitalisation du crime, en même temps que l'ouverture de monde numérique comme  nouvel espace de conflits des états, nouvel espace de bataille des militaires.

De nombreuses voies sont apparues pour transférer les bonnes idées de la filouterie et de la criminalité du monde réel au monde numérique, exploitant l'éventail des talents de l'écosystème de la cybercriminalité : Le malvertising, le spam, les botnets et les chevaux de Troie n'étaient que le début que la production de codes et de services malveillants Les cybercriminels, comme leurs homologues du monde réel, se sont organisés et ont créé une véritable industrie du cybercrime et du crime numérisé.

\textbf{Spam}

L'usage malveillant de la messagerie, s'est fortement basé sur les techniques liées au spam. Le SPAM est un bon moyen efficace de propager des vers; une lettre en chaîne d'un ami ou d'un parent involontaire était typique. Mais il est rapidement devenu monétisé par des escroqueries en ligne.

L'une des campagnes de spam coordonnées les plus connues a été la prolifération du spam  pour les  pharmacies en-ligne. Les «webmasters» généreraient du trafic vers les magasins en ligne gérés via des services de SPAM, pour présenter de nombreux médicaments divers et varié sur ordonnance à des prix largement réduits. Ce fut une époque ou le monde fut inondé par la plus grande campagne pour une nouvelle molécule miracle, le VIAGRA essentiellement vendu par ces pharmacie en-ligne.

Avec des hôtes à l'épreuve des balles en place, les cybercriminels ont probablement gagné des milliards grâce au spam dans les pharmacies, et la cybercriminalité à motivation financière était là pour rester.

\textbf{Botnets}

L'un des botnets les plus prolifiques de l'époque était le botnet Storm, autrefois surnommé le « supercalculateur le plus puissant» du monde, certains experts estimaient en 2007 que les millions d'hôtes infectés le rendait plus puissant que le supercalculateur BlueGene d'IBM de l'époque.

Storm a été conçu pour la furtivité et le profit. À son apogée, les estimations de la taille du botnet allaient d’un million à 10 millions d’ordinateurs infectés. Storm s'est écarté des standarts de ses prédécesseurs sur le bruit et l'agressivité pour privilégier une approche plus patiente et silencieuse.

Dans le cadre de sa tactique furtive, le botnet a utilisé un modèle peer-to-peer distribué. Il a utilisé un DNS à flux rapide et un polymorphisme pour échapper aux défenseurs et infecter un nombre croissant d'ordinateurs.

Cette notoriété et cette efficacité a rapidement fait de Storm, la norme pour tous les futurs botnets tant sur le plan de l'architecture que des des techniques utilisées.

\textbf{Chevaux de Troie}

L'un des pères du business de la malveillance par  cheval de Troie est l'auteur du cheval de Troie bancaire Zeus / Zbot qui ciblait les utilisateurs principalement par le biais du spam, du phishing, de la publicité ou de l'ingénierie sociale. Ce code est rapidement passé d'un simple cheval de Troie bancaire à un kit de de type crimeware à part entière et a marqué le début du cybercrime-as-a-service.

Le prix des licences Zeus ont commencé aux alentour de 1000 euros, mais l'auteur a rapidement proposé des versions personnalisées à un prix plus élevé, qui comprenait une assistance 24h / 24 et 7j / 7 et étaient distribuées par le biais de sociétés affiliées.

Le code source de Zeus a été divulgué en ligne en 2011, ce qui a permis à d'autres  cybercriminels moins techniques d'apprendre de l'un des kits de logiciels malveillants efficace les techniques efficaces et les plus répandues à cette époque. Cette fuite de code Zeus serait responsable de plusieurs variantes, notamment Citadel, Gameover Zeus, ICE-9, CIDEX, Ramnit, Dridex, Kronos, Tinba et Panda.

C'est la grande époque qui a fait exploser les besoins en services de filtrage des e-mails en raison de la flambée du spam et du phishing, tandis que les kits d'exploitation et le malvertising ont conduit à un filtrage supplémentaire du contenu Web.  De nouveaux produits et services de sécurité ont vu le jour. 

La coopération entre l’industrie de la sécurité de l’information, l’application de la loi et les processeurs de paiement a profondément affecté les opérations lucratives des cybercriminels, comme celle mise en œuvre par le cheval de Troie Reveton.

% https://www.cert-ist.com/public/fr/SO_detail?code=201301_article


\subsection{2013 à aujourd'hui - le Rançon-logiciel}

Ce n'est peut-être pas la seule caractéristique déterminante des codes malveillants de cette époque, mais les rançon-logiciels ont sans aucun doute eu l'impact le plus significatif sur les stratégies de cybersécurité des systèmes d'information.

À partir de 2020, les estimations des dommages causés par les attaques de rançon-logiciels se chiffrent en millions d'euros sur l''ensemble de l'économie mondiale. Ces rançon-logiciels exploitent les vulnérabilités des systèmes informatiques qui déploient très rapidement des 
 nouvelles technologies pas ou peu matures.

\subsubsection{CryptoLocker}
Lancée en 2013, l'attaque rançon-logiciels CryptoLocker a fourni aux criminels une formule gagnante en combinant deux technologies : le chiffrement et la crypto-monnaie. Cryptolocker et beaucoup de ses descendants ont également ressuscité un ancien vecteur de menace qui était en sommeil depuis plus d'une décennie : les macros malveillants dans les logiciels de documents.

Depuis lors, de nombreux groupes à l'origine de certaines des familles de rançon-logiciels les plus prolifiques ont perfectionné leurs compétences et se sont adaptés à un environnement en mutation.

Aucun individu ni aucune industrie n'est à l'abri d'une attaque de rançon-logiciel. Aucune technologie seule ne suffit pour arrêter ces typologies de menaces car elles utilisent des vulnérabilités multiples (HOT : Humaines (ingénierie sociale), techniques, Organisationnelles). 

\subsubsection{Double-extorsion}

Avec le rançon-logiciel Maze, les criminels ont popularisé un méthode devenue une norme commune aujourd'hui : la double extorsion, par laquelle les cybercriminels non seulement chiffrent et volent des données, mais menacent de les publier si les cibles ne paient pas.

Même si une entreprise pouvait se remettre complètement d'une attaque de rançon-logiciel sans payer la rançon, elle pourrait toujours se retrouver avec de grandes difficultés tant réglementaire que contractuelle quand des données volées contenant des informations de clients ou de partenaires se retrouvent  publiées

Sous la menace d'une violation de données , certaines entreprises peuvent aussi choisir de payer la rançon plutôt que les amendes officielles qui leur sont imposées. Payer la rançon pourrait être l'option la moins chère des deux, compte tenu des coûts de récupération et des dommages à la marque associés à cette violation. 

\subsubsection{La bataille continue}

Les rançon-logiciels ont bouleversé le monde de la cybersécurité peut-être plus que tout autre vecteur d'attaque.

En 2020, les opérateurs de Maze ayant pris un retraite surement bien méritée, c'est le système EGREGOR qui a pris le relais avec des temps de latence réduite pour déployer ses mécanismes de chiffrement. Les opérateurs ont par ailleurs bien industrialisées les demandes de rançons.

De nombreuses organisations ont revu leur politique de sécurité pour faire face à ces menaces. La notion de plan de cyberdéfense  vu le jour. 

Mais cela nous  conduit aussi à réévaluer les contrôles de sécurité, à créer ou à renforcer les cultures de sécurité au sein des organisations, mais aussi innover de nouveaux produits et services de cybersécurité.

La bataille contre la criminalité numérique continue, elle est un autre d'une bataille plus large incluant la bataille des états dans le Cyberespace. Tous ces champs s'interpénètrent Cyberdefense d'Entreprises, d'Etat, cybercriminalité,  cyberespionnage. La bataille continue. 



}
\mode<all>{%-------------------------------------------------------------
%               FR CYBERDEF SECOPS COURSE
%              $Chapitre : Threat Management
%                         $theme : Threat Intelligence
%        $File :  chap-ThreatMan-threatintel.tex
%                             2020 eduf@ction
%-------------------------------------------------------------
\uchap{chap-Vulman-threatintel.tex}
%-------------------------------------------------------------

\section{ANTICIPER  les menaces}

\subsection{Surveiller et anticiper : Cyber-Threat Intelligence (CTI)}

Une menace informatique ou cyber se matérialise par la combinaison de trois facteurs : 

\begin{itemize}
  \item une intention de nuire, 
  \item une capacité d’attaque et 
  \item une opportunité à exploiter, c’est-à-dire une vulnérabilité de nature technique ou humaine. 
\end{itemize}


Les motivations sont différentes et susceptibles de cibler tous types d’organisations : des hacktivistes poussés par une idéologie, les cybercriminels motivés par l’appât du gain ou  des groupes sponsorisés par un État. On peux se référer à des \ulink{https://attack.mitre.org/groups/}{origines d'APT (Mitre)} avec les groupes à l'origine de ces menaces.



La «threat intelligence» est un service de renseignement concernant les cyber-menaces. Les solutions SIEM par exemple, possèdent leurs propres sources, il n'y a pas que les SIEM qui peuvent utiliser ces sources,  il est possible de connecter de nombreux outils de détections à d’autres sources externes de «threat intelligence». De manière non exhaustive un service de «threat intelligence»  fournit des éléments comme :
\begin{itemize}
  \item \textUK{Malicious / Phishing IP / URL} : une liste d’URL utilisée pour  délivrer un fichier malicieux ou procéder à une attaque par hameçonnage;
  \item 	\textUK{Botnet C\&C URL}» : une liste d’URL utilisée pour héberger des serveurs de commande et contrôle de logiciel malveillant ou de réseaux de machine zombie;
  \item \textUK{Malicious Hash} : liste d’empreintes de logiciels malveillants connus et ayant déjà été analysés;
  \item \textUK{IP Reputation} : liste d’adresses IP suspectées dans des attaques informatiques ou en lien avec une cyber-menace (pouvant être utilisée en Black List sur les systèmes de filtrage, comme historiquement).
\end{itemize}


La Cyber Threat Intelligence est une activité possédant un double objectif  :  l’étude et la surveillance de ces cybermenaces. Pour se faire, la majorité des fournisseurs  de flux d’information de type CTI se basent sur deux approches :


\begin{itemize}
  \item La première approche est l’analyse des attaques passées qui permet de caractériser ces dernières par ces marqueurs techniques. Le but pour le défenseur est de se prémunir et de bloquer au plus vite une campagne qui se renouvellerait via l’utilisation de ces marqueurs techniques. (Voir MISP)
  \item La deuxième approche est la surveillance directe des attaquants. L’objectif est de se placer en amont d’une attaque et de détecter des éléments permettant d’identifier sa préparation (en sources ouverts, ou en infiltration dans le darkweb) : émergence de signaux faibles, objectifs , mode opératoire,  capacités, organisation de l'attaque. Cette approche nécessite un dynamique d’anticipation face à ces menaces,
\end{itemize}


% Begin PRZ ===========================
\begin{frame}
\frametitle<presentation>{Threat Intelligence}
% end header PRZ =======================
La surveillance et le renseignement de la menace au sens général du terme (Threat Intelligence) devraientt contenir les 2 niveaux :

\begin{itemize}
  \item Le renseignement à \textbf{vocation} cyber qui comprend toutes les analyses et informations permettant d'anticiper et de caractériser une menace qui pourrait s'exprimer dans le monde numérique;
  \item Le renseignement \textbf{d'origine} Cyber, dont les données techniques liées à des attaques, menaces qui permettent de configurer des systèmes de détection et de réponse.
\end{itemize}
\end{frame}

% end PRZ ===========================

Il est vrai qu'encore aujourd'hui parler de \g{threat intelligence} nous dirige systématiquement sur la deuxième assertion.


\mode<all>{\picframe{Latex/Sources/EDU/Pictures/img-veillecyber}{Veille cyber, une veille sur les risques}{0.7}{lblveillecyber}}

% Begin PRZ ===========================
\begin{frame}
\frametitle<presentation>{Veiller et surveiller}
\framesubtitle<presentation>{2 axes}
% end header PRZ =======================

Veiller et surveiller les menaces, détecter les attaques  nécessite d'analyser deux axes : 

\begin{itemize}
  \item Les menaces génériques, ou ciblant un domaine  particulier (Santé, Industrie, Banque ...) que l'on trouve généralement en utilisant des technologies de \g{threat Intelligence}; 
  \item Les menaces ciblées, dont les indices d'émergence peuvent être détecter en analysant la menace ou  en recherchant des indices de compromissions quand ces menaces sont actives dans le périmètre de l'entreprise. \g{threat Detection, Hunting ...} 
\end{itemize}
\end{frame}

% end PRZ ===========================



% Begin PRZ ===========================
\begin{frame}
\frametitle<presentation>{Veiller et surveiller}
\framesubtitle<presentation>{2 manières}
% end header PRZ =======================

et ceci de deux manières :

\begin{itemize}
  \item Surveillance de l'écosystème de la menace (IOC,  DarkWeb, Threat Intelligence...) 
  \item Recherche de compromission, ou d'infection (Threat Hunting, ...) 
\end{itemize}
\end{frame}
% end PRZ ===========================

Ce sont des sujets que nous aborderons dans le processus de gestion de la menace.
La surveillance des menaces génériques relève d'action de veille comme cela est fait pour les vulnérabilités. Les scénarios de message sont vus comme des éléments de signature d'une attaque ou d'une tentative ou de préparation  d'attaque.

Je vous propose de présenter  la gestion de la menace sous la forme de ces trois thèmes  (\uref{Gestion de la menace}{lbllogman}).

\begin{itemize}
  \item Log Management;
  \item Threat Detection;
 \item  et en dernier lieu des élémens Threat Intelligence (au sens renseignement);
\end{itemize}

\mode<all>{\picframe{Latex/Sources/EDU/Pictures/img-threatvigil}{Les sources}{0.7}{lblthreatvigil}}

\subsubsection{Quelques producteurs de CTI}
\toolsbox{FireEye}{cti}

\subsubsection{Quelques outils de CTI}
\toolsbox{MISP}{cti}
\toolsbox{OpenCTI}{cti}

%TODO Tester les SIEM et les SOC et le deuxième niveau de maturité

%TODO Niveau de maturité dans les processus

% TODO le renseignement


%*************************************

\subsection{Les surveillances} 


\subsubsection{Surveillance de la compromission}

Un des domaines de la surveillance est donc celui de la compromission. C'est à dire la surveillance dans le fameux Darkweb de l'émergence de données volées, \g{perdues} par une entreprise ou par un particulier.

Ce domaine, dénommé par certains \g{Leak Intelligence} ou \g{Leak Management}, correspond à la gestion des fuites de données, au sens de leur détection et la recherche de \g{la source} de fuite. C'est souvent dans cette dynamique de découverte d'information \g{interne} dans des espaces \g{malveillants} ou pas que l'on découvre des compromissions techniques ou non techniques issues d'attaques.

La compromission la plus connue reste encore de nos jours  la fuite du couple Utilisateur/mot de passe sur des sites hackés. Il existe des bases de données publiques qui publient ces données \g{compromises} comme par exemple : \ulink{https://haveibeenpwned.com}{Have i been pwned}, qui permettent à partir d'une adresse mail ayant servie d'identifiant de savoir si on a été compromis sur un site qui aurait été piraté.

\subsubsection{Surveillance des fragilités}

Comme nous l'avons vu dans le chapitre sur les vulnérabilités,  des \g{scans} de vulnérabilités sur des plages d'IP de l'entreprise permettent de déterminer les fragilités de services ouverts ou accessibles. Généralement organisées dans une dynamique d'audit, ces évaluations de sécurité sont conduites avec un cadre contractuel et légal. Il existe pourtant des entreprises qui fournissent des informations de fragilités sur des entreprises ou des plages d'IP. On trouvera pas exemple sur le site de \ulink{https://www.shodan.io}{SHODAN} des informations intéressantes (et payantes) sur des fragilités de systèmes dont une grande partie de systèmes d'objets connectés.

\begin{techworkbox}{Base de données de fragilitées publiques}
Les techniques et les sites qui propose des outils ou l'accès à des bases de données de \g{sites} vulnérables. C'est un sujet intéressant pour \fichetech de synthèse sur ce qui existe sur le marché. %sujets : https://alternativeto.net/software/shodan/
\end{techworkbox}


\toolsbox{Oval}{vmt}
\toolsbox{OpenSCAP}{vmt}


Les approches des analyses du renseignement peuvent se classer à différents niveaux :

\begin{itemize}
  \item \textbf{Niveau Stratégique} : ce sont en général des analyses de très haut niveau mais peu techniques et destinées à des décideurs. Cela peut-être par exemple des rapports d'analyses géopolitiques sur des adversaires qui ciblent un secteur donné.
  \item \textbf{Niveau Tactique} : ce sont souvent des documents qui donnent des informations sur les outils et méthodologies utilisés par les menaces (analyse de malware, contournement d’anti-virus, outils utilisés pour mener des attaques DDoS, etc.).
  \item \textbf{Niveau Opérationnel} : l’objectif est d’anticiper l'attaque en étant au plus près de l’attaquant (Ecoute, infiltration  des hacktivistes, cybercriminels ...).
  \item \textbf{Niveau Technique} : il est composé d’indicateurs de compromission (IPs, URLs, noms de domaine, listes de hashs, etc.) qui permettent d’identifier et donc de bloquer directement une attaque.
\end{itemize}


%*************************************
\subsubsection{Surveillance du ciblage}

% Threat Hunting 

La surveillance du ciblage, que les anglo-saxons appelle le \g{TARGETING} est aussi un élément d'anticipation. 
En effet, ces éléments sont souvent les premiers signaux d'une préparation d'un évènement \g{cyber} qui pourrait toucher l'entreprise.

On y trouve l'émergence de la collecte d'information sur une cible donnée, la mise en oeuvre dans les codes  malveillants de ciblage d'IP spécifique, etc...

% Begin PRZ ===========================
\begin{frame}
\frametitle<presentation>{Surveillance du ciblage}
\framesubtitle<presentation>{l'outillage du \g{targetting}}
% end header PRZ =======================
Il y a deux types d'outils pour se faire : 

\begin{itemize}
	\item La surveillance classique du web de type \g{cyberveille}, qui permet de découvrir des éléments compromis appartenant à l'entreprise (soit les données, soit des informations permettant de déduire que l'entreprise a été compromise).
	\item L'analyse en temps réel des codes malveillants qui peut permettre en regardant de manière détaillée l'évolution du code pour comprendre et connaître les modalités des attaques et les nouvelles cibles.
\end{itemize}
\end{frame}
% end PRZ ===========================


%*************************************
\subsubsection{Que faire des ces informations}

Disposer des fragilités de l'entreprise, et connaître les scénarios potentiels permet d'évaluer un niveau de risque. 

\subsection{de l'outillage sur la menace}

\subsubsection{la gestion des menaces CTI}

Le projet OpenCTI (Open Cyber Threat Intelligence), développé par l’ANSSI en partenariat avec le CERT-EU, est un outil de gestion et de partage de la connaissance en matière d’analyse de la cybermenace (Threat Intelligence). Initialement conçue pour structurer les informations de l’agence relatives à la menace informatique, la plateforme facilite aussi les interactions entre l’ANSSI et ses partenaires.

L’outil, intégralement libre, est  disponible à l’usage de l’ensemble des acteurs de la « threat intelligence ». L’application permet ainsi de stocker, organiser, visualiser et partager leurs propres connaissances en la matière. L'outil est structuré autour de des méodèle de description STIX 2 et utilise un base hypergraphe.

Le projet \ulink{https://www.ssi.gouv.fr/actualite/opencti-la-solution-libre-pour-traiter-et-partager-la-connaissance-de-la-cybermenace/}{OpenCTI} a été initié en septembre 2018 par l’ANSSI et co-développé avec le CERT-EU en l’absence de solutions complètement appropriées pour structurer, stocker, organiser, visualiser et partager la connaissance de l’ANSSI en matière de cybermenace, à tous les niveaux.

OpenCTI peut être vu comme un premier outil au centre d'un \g{fusion center}.




\subsubsection{STIX et TAXII}

%The structuration of the data is performed using a knowledge schema based on the STIX2 standards. It has been designed as a modern web application including a GraphQL API and an UX oriented frontend. Also, OpenCTI can be integrated with other tools and applications such as MISP, TheHive, MITRE ATT&CK, etc.

Les \textbf{modélisations} des attaques sont un large champ de recherche et d'outillage, elles sont au coeur de la compréhension des attaques mais surtout au coeur de la détection de celle-ci. Juste à titre d'illustration, nous pouvons parler d'un modèle comme STIX™ (Structured Threat Information Expression) langage et format de donnée permettant de modéliser et échanger des informations techniques sur les processus d'attaque cyber. 
Je vous propose d'explorer pour cela sur le site \ulink{https://oasis-open.github.io/cti-documentation/}{STIX sur GitHub}.

Le premier type de standard à avoir été proposé est le format de modélisation des informations. L’objectif est qu’un émetteur puisse communiquer à l’ensemble de ses destinataires les données collectées à partir d’un support et dans un format défini. Le plus connu d’entre eux, OpenIOC, standard historique créé par l'ex société MANDIANT  pour échanger des IOC. À la base, OpenIOC est un format (en XML) utilisé par les outils de la société mais qui a été rendu open-source pour permettre un usage par tous. Un fichier OpenIOC décrit les symptômes à rechercher pour identifier une menace.
D’autres formats comme comme SNORT  ou encore YARA qui se focalisent sur les caractéristiques intrinsèques d’un fichier et est particulièrement pertinent pour identifier des souches de familles de code malveillant.
Ces formats de modélisation rencontrent cependant des limites car centrés sur l'événement et n'offrent pas de vision globale de la cybermenace à l’origine de l’attaque. 
STIX (Structured Threat Information Expression) et TAXII (Trusted Automated eXchange of Indicator Information), et auparavant CybOX (Cyber Observable eXpression), ont été développés dès 2012 aux États-Unis par le CERT-US (United States Computer Emergency Readiness Team) du Department of Homeland Security  puis repris par les organisations MITRE et OASIS. Ces langages tentent de combler les lacunes.

STIX est un langage normalisé et structuré permettant de représenter les informations sur les menaces cyber, et en particulier des indicateurs de compromission (IOC) issus des analyses de Cyber Threat Intelligence (CTI).

TAXII est quand à lui un protocole conçu pour améliorer la qualité des échanges des informations de type CTI, formatées dans le langage STIX.

Les utilisateurs STIX peuvent en particulier décrire et modéliser les concepts  :

\begin{itemize}
  \item \textbf{Attack Pattern} : mode opératoire de l’attaquant ;
  \item \textbf{Threat Actor }: individus, groupes ou organisations agissant avec une intention malveillante ;
  \item \textbf{Malware} : logiciel malveillant utilisé pour compromettre la confidentialité, l'intégrité ou la disponibilité du système d’information de la cible (voir \ulink{https://www.virustotal.com/gui/}{VirusTotal};
  \item \textbf{Tool} : logiciel légitime utilisé par les cyber-menaces dans le cadre de leurs attaques ;
  \item \textbf{Vulnerability} : vulnérabilité présente dans un logiciel qui est exploitée directement par un attaquant afin de compromettre un système d’information ;
  \item  \textbf{Indicator} : indicateur utilisé pour détecter et bloquer une activité suspecte ou malveillante sur le système d’information.
\end{itemize}


\subsubsection{MISP}

La distribution d'informations caractérisant les attaques nécessite au delà de la modélisation des attaques des indicateurs des IOC (\UKword{Indice of compromission}) facilement détectables dans les plateformes techniques de détection que nous verrons par la suite.

MISP (Malware Information Sharing Platform and Threat Sharing) est une solution open-source permettant la collecte, le stockage, la distribution et le partage d’IOC liés aux malwares. Cet outil permet à divers organismes de partager les indicateurs de compromission identifiés lors des activités de SOC et campagnes de réponses à incidents (CERT et CSIRT).  D’un point de vue technique, MISP est une plateforme d’échange d’IOC où chaque acteur peut entrer et organiser ses IOC afin de les publier pour les partager aux autres acteurs présents sur ce portail.  MISP fournit des fonctionnalités pour faciliter les échanges d'informations, mais aussi l’intégration de l'information par les IDS (Intrusion Detection System) et les outils de défense comme ceux d'analyse de logs et les SIEMs (Security Information and Event Management).


\subsubsection{Vérifier une donnée par rapport à de la  CTI}

Disposer de Treat Intelligence donc de renseignement  est complémentaire de la connaissance des scénarios issus des analyses de risques, mais enrichisse aussi les analyses de risques par les modèles d'attaques existants.

\begin{itemize}
  \item Tester un fichier, une url (\textbf{Malveillance}) :  \ulink{https://www.virustotal.com/gui/}{VirusTotal};
  \item Tester une adresse email (\textbf{Compromission})  : \ulink{https://haveibeenpwned.com}{haveibeenpwned};
  \item Tester sans action une adresse IP publique ou un domaine (\textbf{Vulnérabilités Internet)} : \ulink{https://www.shodan.io}{Shodan};
\end{itemize}



}
\mode<all>{%-------------------------------------------------------------
%               FR CYBERDEF SECOPS COURSE
%              $Chapitre : Threat Management
%                         $theme : Threat Detection
%      $File :  chap-ThreatMan-threatdetect.tex
%                             2020 eduf@ction
%-------------------------------------------------------------
\uchap{chap-ThreatMan-threatdetect.tex}
%-------------------------------------------------------------


\section{DETECTER les attaques}

La détection et la surveillance de système d'information d'entreprise, passent souvent pas l'utilisation des informations qui remontent déjà des services de protections périmétriques comme les Firewall (Réseau, Web application...),  proxy,  routeurs, IPS/IDS. Tous les équipements de sécurité installés sur un environnement sont des sources d'évènements qui permettent pour peu qu'elles soient analysées, de détecter des attaques ou le déploiement de menaces variées (comportements déviants, flux d'APT, flux de scans, ...)
Ces sources peuvent être complétées par tout équipement ou système  produisant des journaux d'activité (LOGs).

Par ailleurs, certains équipements dédiés de \g{surveillance} appelés \g{sondes} sont aussi aptes à remonter des évènements de détection sur la base de \g{signatures} ou \g{algorithmie} dont des IOC (\g{indice of compromission}).

\begin{techworkbox}{Sondes souveraines}
La Loi de Programmation Militaire (LMP) appelle les OIV à s’équiper de sondes dites souveraines, qualifiées par l’ANSSI et  disposant d’un niveau de sécurisation élevé. C'est un sujet technique intéressant pour  \fichetech. 
\end{techworkbox}

Dans cette partie nous allons traiter donc de la détection de menaces sur la base  de différentes techniques de détection et de surveillance :

\begin{itemize}
  \item Logs et sources d'évènements; 
  \item Corrélation d'évènement avec les SIEM;
  \item Évènements provenant des terminaux.
\end{itemize}

et nous explorerons quelques éléments d'organisation autour de la détection et de la surveillance.

}
\mode<all>{%-------------------------------------------------------------
%               FR CYBERDEF SECOPS COURSE
%              $Chapitre : Threat Management
%                   $theme : Threat Detection
%               $File :  chap-ThreatMan-log.tex
%                             2020 eduf@ction
%-------------------------------------------------------------
\uchap{chap-ThreatMan-log.tex}
%-------------------------------------------------------------

\subsection{Log Management }

\subsubsection{Traces, journaux, logs}

Dans le domaine informatique et télécom, le terme log (ou ses synonymes traces, journaux) est généralement un fichier, une base de données ou tout autre moyen de stocker des informations.  On parle  donc de stockage d'un historique d'événements qu'un logiciel ou un système souhaite \g{tracer}. 
Ce mot qui est le diminutif  de \textit{logging}, est traduit en français par \g{journal}. Le log est donc un journal horodaté, qui stocke temporellement les différents événements qui se sont produits sur un  logiciel, un ordinateur, un serveur, etc. Il permet ainsi d'analyser avec une fréquence programmée (heure par heure,  minute par minute, etc) l'activité  d'un processus technique.  

\mode<all>{\picframe{Latex/Sources/EDU/Pictures/img-log}{Sources de log}{0.8}{lbllog}}

Vous trouverez des éléments très intéressants  dans le 
\ulink {https://www.microfocus.com/media/white-paper/the_complete_guide_to_log_and_event_management_wp_fr.pdf}{Guide LOG MANAGEMENT} édité par NetIQ.

Dans notre cas d'usage, les log sont les éléments techniques  bruts d'un capteur d'événements. L'objectif est d'assurer que l'ensemble des journaux d'évènements contient suffisamment d'informations pour assurer des corrélations permettant de  constituer des signatures d'attaques. 

\begin{notebox}{Sources choisies des logs }
Pour détecter des attaques en temps réel, il faut disposer des informations caractéristiques de ces attaques captées dans le système d'information. Cela nécessite donc de sélectionner les bonnes sources (équipements réseaux, ou informatiques), les bons journaux, les bonnes traces dans ces journaux.  Ces choix sont donc primordiaux. Pour faire ces choix, il est nécessaire de connaître la capacité des équipements et des systèmes logiciels de générer ces traces. Au coeur de ces évènements il sera alors possible par corrélation de détecter des scénarios complexes d'attaques
\end{notebox} 

%SOURCE  : https://www.microfocus.com/media/white-paper/the_complete_guide_to_log_and_event_management_wp_fr.pdf

Une grande majorité des équipements (réseau, serveurs, terminaux (endpoint), des bases de données ou des applications d’un systèmes d'information peuvent aujourd’hui générer des logs ou traces.  Ces fichiers contiennent, pour chaque équipe, la liste de tous les évènements  \g{traçables} qui se sont déroulés pendant l'execution : réussite ou échec d’une connexion, redémarrage, utilisation des ressources (mémoire, ...).


\mode<all>{\picframe{Latex/Sources/EDU/Pictures/img-logman}{Les logs au coeur de la détection}{0.8}{lbllogman2}}


L'exploitation de ces traces est souvent complexe car chaque équipement dispose de ses propres fonctions de gestion des traces, avec encore dans de nombreux cas un format  d'enregistrement et de stockage propriétaire. Il faut consulter ces logs équipement par équipement. Heureusement, il existe des outils qui permettent de centraliser et de \g{normaliser} ces traces.

On peut citer par exemple, SYSLOG, qui par ailleurs n'est pas le seul type d'outil pour assurer la collecte et la normalisation des traces. SYSLOG reste pourtant un outil de référence dans les architectures de collecte.

\subsubsection{Services et protocole SYSLOG}

Le protocole Syslog est un protocole réseau qui permet de transporter les messages de journalisation générés par les applications vers une machine hébergeant un serveur Syslog. 

Quand un système veut conserver les traces d'un événement, il est possible, d'utiliser syslog pour communiquer les détails de l'événement à un daemon syslog (serveur syslog) qui va le conserver dans une base de données.

Le protocole Syslog est structuré autour de la notion de périphérique, de relais et de collecteur dans une architecture Syslog.

\begin{itemize}
  \item Un \textbf{périphérique} est une machine ou une application qui génère des messages Syslog;
  \item Un \textbf{relais} est une machine ou une application qui reçoit des messages Syslog et les retransmet à une autre machine;
  \item Un \textbf{collecteur} est une machine ou une application qui reçoit des messages Syslog mais qui ne les retransmet pas.
\end{itemize}

Tout périphérique ou relais sera vu comme un émetteur lorsqu'il envoie un message Syslog et tout relais ou collecteur sera vu comme un récepteur lorsqu'il reçoit un message Syslog.

L'intérêt d'un serveur/collecteur Syslog est donc de permettre une centralisation de ces journaux d'événements, permettant de repérer plus rapidement et efficacement les défaillances de machines présentes sur un réseau. On trouvera par exemple, sur le site 
\ulink{https://homputersecurity.com/2018/03/01/comment-mettre-en-place-un-serveur-syslog/}{homputersecurity.com} des éléments pour déployer un serveur SYSLOG, et une description détaillée sur le site \ulink{https://ram-0000.developpez.com/tutoriels/reseau/Syslog/}{Developpez.com}.

 \subsubsection{L'usage des LOGs}
 
 % todo modifcation
 
 Dans un usage de cybersécurité, les traces, journaux et logs informatiques et réseaux structurent les sources d'évènements et de stockages des informations. Ils sont donc :
 
 \begin{itemize}

   \item Un outil indispensable au processus de \textbf{détection de menace}. La gestion des logs (ou d’événements) s’avère en effet un outil primordial pour les analyses à posteriori, mais peut aussi servir dans la détection en temps réel pour peu que les outils d'analyse puissent le faire; Nous verrons cela dans la partie sur la détection de la menace.

   \item  \textbf{Une couverture légale }. Confrontée à une plainte, une entreprise peut utiliser ces traces pour gérer un litige avec un tiers  en attestant de la non-implication de son système d’information ou, a contrario, assumer le litige tout en remontant jusqu’à l’utilisateur concerné. La société peut également utiliser ces traces pour fournir des éléments aux services de polices. La fourniture d'élément probant à valeur légale nécessite quelques précautions. 

   \item  \textbf{Le dépistage des malversations internes ou de comportements déviants}. Les flux illégaux, les flux de données déviants (copies de fichiers en masse avant qu'un salarié quitte l'entreprise par exemple)

\end{itemize}

 Évidement, un outil de gestion de LOG ne serait à lui seul et sans fonction de corrélation avoir la capacité de détecter des évènements liés entre eux. 
 
\begin{techworkbox}{Syslog et cybersécurité}
	L'environnement SYSLOG possède une richesse fonctionnelle qui nécessiterait une présentation détaillée pour en appréhender les capacités et la puissance d'usage. C'est un sujet pour \fichetech de référence. Il existe de nombreuses documentations sur internet, toutefois une présentation détaillée d'une architecture SYSLOG pour un usage de cybersécurité est un sujet à explorer. 
\end{techworkbox}

\subsubsection{Puits de logs}

La construction d’un « puits de log » est une première brique de réponse : il s’agit de collecter, à l’aide d’un outil automatisé du marché, l’ensemble des journaux d’équipements dans un espace de stockage unique. L’un des critères de sélection de cet outil est justement sa capacité à reconnaître différents formats de logs (syslog, traps SNMP, formats propriétaires…).

Le volume d’informations centralisées peut vite exploser : il est important d’éviter la collecte de données inutiles. Par ailleurs, le système peut également être gourmand en puissance de calcul en fonction des périmètres de recherches effectuées.

On parle de log management à partir du moment où les données contenues dans ces  puits sont traitées et exploitées, par exemple pour retrouver un élément dangereux (virus, problème de sécurité…), ou un comportement malveillant (fuite d’information, suppression de données…). Il est nécessaire de cadrer en amont les finalités du projet,  qui peuvent être multiples.

\subsubsection{Outils d'analyses}

Après avoir \textbf{collectés, stockés} des évènements dans un format compréhensible (structuré ou non), il est nécessaire de disposer d'outils de recherche et d'analyses de ces logs. Il existe de nombreux outils  dont beaucoup de codes open source pour se faire.


\mode<all>{\picframe{Latex/Sources/EDU/Pictures/img-log-level}{Les niveaux d'outillage}{0.8}{lblloglevel}}


Nous pouvons citer par exemple \ulink{https://www.graylog.org}{Graylog} basé sur MongoDB pour la gestion des métadonnées et Elasticsearch pour le stockage des logs et la recherche textuelle. Graylog  permet de mieux comprendre l’utilisation d'un système d'information tant dans l'amélioration la sécurité (comportements, évènements à risques, indicateurs de compromission (IOC)) que dans l'usage des applications et services.

%TODO Menaces LOG : Cas des end point
%\subsubsection{le cas des End Point}
%Syslog - osquery

\begin{techworkbox}{Analyse de logs}
	Les outils d'indexation et d'analyse de logs sont nombreux et chacun possède des avantages et des inconvénients, des facilités d'usage de déploiement et des fonctionnalités différentes. C'est un sujet de réalisation pour \fichetech
\end{techworkbox}






}
\mode<all>{%-------------------------------------------------------------
%               FR CYBERDEF SECOPS COURSE
%              $Chapitre : Threat Management
%                   $theme : Threat Detection
%               $File :  chap-ThreatMan-siem.tex
%                             2020 eduf@ction
%-------------------------------------------------------------
\uchap{chap-ThreatMan-siem.tex}
%-------------------------------------------------------------

% SOURCE https://www.logpoint.com/fr/comprendre/quest-ce-que-le-siem/
% SOURCE https://www.nomios.fr/type-de-solution-siem/
% SOURCE https://www.splunk.com/fr_fr/form/7-siem-trends-to-watch-in-2019.html

\subsection{SIEM, une technologie}

\subsubsection{SIEM et Logs}

% Begin PRZ ===========================
\begin{frame}
\frametitle<presentation>{SIEM et Log}
\framesubtitle<presentation>{Threat and leak}
% end header PRZ =======================
Puisque nous disposons maintenant d'une capacité de capter de l'information pertinente pour détecter des attaques, que nous avons une architecture de collecte,  ainsi qu'une architecture de stockage et de filtrage, nous pouvons injecter des informations dans des outils de recherche et de corrélation d'attaque.
Ceci pour peu que l'architecture et les outillages puissent suivre la charge d'analyse en temps réel.

Il ne faudra pas aussi oublier que les traces informatiques et réseaux ne sont pas les seules sources d'informations nécessaires à la détection d'attaques (en temps réel ou différé), il faut aussi connecter des sources de menaces : 
\begin{itemize}
  \item \textbf{Threat Intelligence Database }:  IOC et identifiants des sources malveillantes (IP, noms de domaine, serveur mail ...)
  \item  \textbf{Leak} : Fuite de données détectée par la surveillance du Web et du Darknet.
\end{itemize}

\end{frame}
% end PRZ ===========================


\subsubsection{Un peu d'histoire}

Le SIEM (\textit{Security Information and Event Management})  est aujourd'hui l'aboutissement d'un voeu  très ancien des responsables sécurité qui supervisent depuis bien des décennies des systèmes de contrôle périmétrique : Corréler tous les évènements de fonctionnement sur l'ensemble des équipements. 

Le SIEM se définit donc comme la collecte, la surveillance, la corrélation et l’analyse en temps réel d'événements fonctionels et techniques provenant de sources disparates. Les solutions SIEM d’aujourd’hui permettent à l'entreprise de réagir rapidement et avec suffisamment de traçabilité et d'éléments techniques pour  accélérer l'analyse et la réponse à incident.
Une solution SIEM assure la gestion, l’intégration, la corrélation et l’analyse généralement en seul lieu, ce qui facilite la surveillance et la résolution des problèmes des infrastructures hétérogènes, complexes et distantes. Sans ce type d'outillage, un analyste de sécurité doit passer en revue des millions de données non comparables, stockées dans des « silos » pour chaque matériel, chaque logiciel et chaque source de sécurité. Un SIEM est donc avant tout un outil de fédération des évènements de sécurité.

L'acronyme SIEM ou «gestion des informations de sécurité» fait référence à des technologies combinant à la fois la gestion des \textbf{informations} de sécurité et la gestion des \textbf{événements} de sécurité. Comme ils sont déjà très similaires, le terme générique plus large peut être utile pour décrire les outils récents  le conjuguant. Là encore, il est essentiel de différencier la surveillance des événements de la surveillance des informations générales. Un autre moyen essentiel de distinguer ces deux méthodes consiste à considérer la gestion des informations de sécurité comme une sorte de processus à long terme ou plus large, dans lequel des ensembles de données plus diversifiés peuvent être analysés de manière plus méthodique. En revanche, la gestion des événements de sécurité examine à nouveau les types d’événements utilisateur pouvant constituer des signaux d’alerte ou indiquer aux administrateurs des informations spécifiques sur l’activité du réseau.

C'est souvent l'usage d'un SIEM dans une ambiguïté de gestion long terme de la sécurité en tant que propriété d'un système d'une part, et la gestion court terme de l'urgence d'une attente à la sécurité qui pose problème dans les projets et dans les opérations.
 
 Ce genre d'outillage est passé par différentes étapes de maturation avec des SIM et SEM  et enfin des SIEM.
Il s’agit de combiner les fonctions de gestion des informations (SIM, Security Information Management) et des évènements  (SEM, Security Event Management) en un seul système de management de sécurité.

\begin{itemize}

	\item dans la gestion des informations de sécurité (SIM) , la technologie consiste  à collecter des informations à partir des journaux d'équipement de sécurité, qui peut consister en différents types de données. Globalement on peut dire qu'un SIM est un aimant important pour des équipes de supervision de la sécurité périmétrique. d'une part pour la traçabilité et le reporting de sécurité.
	\item technologies spécialement conçues pour rechercher des authentifications suspectes, des ouvertures de session sur un compte ou des accès de gestion de haut niveau à des heures précises du jour ou de la nuit.
\end{itemize}

Bien qu'outillant des processus très similaires mais distincts, les trois Sources/EDU/acronymes SEM, SIM et SIEM ont tendance à être confus ou à causer de la confusion chez ceux qui seraient peu familiarisés avec les processus de sécurité.
La similitude entre la gestion des événements de sécurité ou SEM et la gestion des informations de sécurité ou SIM est au cœur du problème.

Ces deux types de collecte d'informations concernent la collecte d'informations de journal de sécurité ou d'autres données similaires en vue d'un stockage à long terme, ou l'analyse de l'environnement de sécurité d'un réseau. Quoi qu'il en soit, de nos jours le terme SIEM est utilisé quelque soit d'ailleurs l'usage.

% SOURCE https://www.linkbynet.com/fr/recourir-a-un-siem

Plus concrètement, un système de type SEM centralise le stockage et l’interprétation des logs en temps réel et permet une analyse. Les experts en cyber sécurité peuvent ainsi prendre des mesures défensives plus rapidement. Un système de type SIM collecte pour sa part des données et les place dans un référentiel à des fins d’analyse de tendances. Dans ce cas, la génération de rapports de conformité est automatisée et centralisée.

Le SIEM, qui regroupe ces 2 systèmes, accélère donc l’identification et l’analyse des événements de sécurité, atténue les conséquences d’attaques et facilite la restauration qui s’ensuit. Pour y parvenir, il collecte les événements, les stocke (avec normalisation) et agrège des données pertinentes mais non structurées issue de plusieurs sources. L’identification des écarts possibles par rapport à la moyenne / norme nourrit la prise de décision. En outre, les tableaux de bord générés contribuent à répondre aux exigences légales de conformité de l’entreprise.

En d’autres termes, avec le SIEM, les équipes de sécurité opérationnelle industrialisent la surveillance tout en simplifiant l’analyse de multiples sources d’événements de sécurité (antivirus, proxy, Web Application Firewall…). La corrélation des événements provenant d’applications ou d’équipements très variés est aussi facilitée. De quoi détecter des scénarios de menaces avancées.

Dans la pratique, Il existe 3 grandes manières d'opérer ou de faire opérer un SIEM :

\begin{itemize}
  \item SIEM  déployé et intégré dans l'entreprise;
  \item SIEM basé dans le cloud;
  \item SIEM géré / managé en mode MSSP.
\end{itemize}

S'équiper d’une solution de type SIEM nécessite encore un investissement conséquent en raison de la complexité de la mise en œuvre des gros outils du marché. Toutefois, bien qu’initialement destiné aux grandes entreprises, le SIEM peut être déployé dans tous les types d’organisations, même les plus petits, avec des solutions de plus en plus automatisées. Toutefois, il ne faudra pas oublier que la configuration d'un SIEM (création de scénario de détection (USE CASE) demande de la ressource et des compétences spécifiques.
Par ailleurs, l'important sera par la suite de définir quel type de réaction à la suite de la détection d'un évènement à risque.


Un SIEM s’avère capable de détecter des incidents de sécurité qui seraient passés inaperçus. Pour une raison simple : les nombreux hôtes qui enregistrent des événements de sécurité ne disposent pas de fonctions de détection d’incidents.

Le SIEM dispose d'une fonction de détection grâce à sa capacité de corrélation des événements. Contrairement à un système de prévention d’intrusion qui identifie une attaque isolée, le SIEM regarde au-delà. Les règles de corrélations lui permettent d’identifier un événement ayant causé la génération de plusieurs autres (hack via le réseau, puis manipulation sur un équipement précis…).

Dans de tels cas de figure, la plupart des solutions ont la capacité d’agir indirectement sur la menace. Le SIEM communique avec d’autres outils de sécurité mis en place dans l’entreprise (Exemple pare-feu) et pousse une modification afin de bloquer l’activité malveillante. Résultat, des attaques qui n’auraient même pas été remarquées dans l’entreprise sont contrecarrées.

% Begin PRZ ===========================
\begin{frame}
\frametitle<presentation>{Fonction SIEM}
% end header PRZ =======================
\begin{itemize}
  \item la première fonction d'un SIEM est déjà de corréler les événements provenant des composants de sécurité;
  \item la deuxième  fonction de corréler des événements de comportement du SI;
  \item la troisième fonction de corréler avec des événements externes au SI sur la base de capteurs externes (threats intelligence de type renseignement).
\end{itemize}
\end{frame}
% end PRZ ===========================


\mode<all>{\picframe{Latex/Sources/EDU/Pictures/img-siem}{architecture d'un SIEM}{0.8}{lblsiem}}

Pour aller encore plus loin, une organisation peut choisir d’intégrer à son SIEM,  une Cyber Threat Intelligence (CTI ou Flux de renseignement sur les menaces). 

Selon la définition de Gartner, la Cyber Threat Intelligence (CTI) est la connaissance fondée sur des preuves, y compris le contexte, les mécanismes, les indicateurs, les implications et des conseils concrets, concernant une menace nouvelle ou existante ou un risque pour les actifs d’une organisation qui peuvent être utilisés afin d’éclairer les décisions concernant la réponse du sujet à cette menace ou à un danger.

La CTI consiste donc à collecter et organiser toutes les informations liées aux menaces et cyberattaques, afin de dresser un portrait des attaquants ou de mettre en exergue des tendances (secteurs d’activités visés, les méthodes d’attaque utilisées, etc.). Résultat, une meilleure anticipation des incidents aux prémices d’une attaque d’envergure.

 \subsection{la détection d'incidents}

\subsubsection {Collecter les logs }
\textbf{COLLECTE}  : Le premier acte est d'alimenter le système avec les «logs» des composants du système d’information (outils de sécurité comme les IDS, IPS, «HoneyPot», Pare-feu, etc… que ceux des autres composants comme les équipements réseaux, les serveurs mandataires, de fichier, de mail, les services dans le nuage, etc…);
\subsubsection {Normaliser et indexer les logs}
\textbf{NORMALISATION} : La normalisation des logs est un processus d’uniformisation des données collectées qui, quelle que soit la source, détecte le type de donnée pour intégrer chaque information dans une base de données. Par exemple plusieurs équipements différents écrivent chacun dans un de leur journal d’activité l’adresse IP source, la date, l’heure et l’activité. Chaque équipement aura son format particulier, il va donc falloir les analyser pour détecter et labelliser chaque information utile. Sur ce point, les constructeurs d’équipement et les éditeurs de solutions SIEM travaillent ensemble afin de proposer des modules de normalisation qui couvrent la plupart des usages clients. De plus de nombreux formats de solutions de log sont nativement pris en charge par les SIEM. Les informations sont ensuite indexées pour optimiser les temps de traitement de l’information.
\subsubsection {Corréler les logs}
\g{CORRELATION} : La valeur ajoutée du SIEM vis-à-vis d’un simple serveur de « log», se joue en grande partie sur la corrélation. En effet, ici il est question de pouvoir mettre en évidence des activités en analysant conjointement plusieurs évènements de sources différentes. Bien que ces évènements puissent paraître anodins de manière isolée, ensemble ils forment un indicateur de compromission. Cette corrélation se base sur des règles forgées à partir de cas d’usage (\textit{USE CASE})

\subsubsection {Analyser et alerter}
\textbf{ALERTE} : Ensuite il faut analyser le contenu corrélé pour faire le tri entre les faux positifs et les vrais positifs. Il y a aussi la détection de menace dont le but est de comparer des IOC à des listes de sources malveillantes provenant d’adresse IP, de nom de domaine, de signature de code malveillant, etc. Chaque solution propose ses propres algorithmes d’évaluation du risque lié aux évènements analysés. Une fois les analyses réalisées, des mécanismes peuvent déclencher des alertes si un seuil a été franchi. 
\subsubsection {Réponse graduelle}
Il est possible de proposer des actions d’isolation ou de réponse automatisée. Par exemple : Si les «logs» d’un Anti-virus analyse la présence d’un rançongiciel assimilable au logiciel malveillant «Wannacry» et que le poste en question fait des tentatives de connexions extérieures vers le port 445, il est possible de déclencher un isolement de la machine du réseau à l’aide d’un script ou de l’agent présent sur le poste. 





% TODO Voir article sur les problèmatique SIEM
% https://www.journaldunet.com/solutions/expert/39767/comprendre-les-raisons-de-l-echec-des-siem-et-sim.shtml

La cyber-protection d'une entreprise est principalement basée sur les outils de protection périmétriques que ceux ci soient des équipements physiques ou qu'ils soient dans le cloud : systèmes de détection d’intrusion (IDS), scanners de vulnérabilités, antivirus ainsi que systèmes de gestion et corrélation d’événements sécurité (SIEM). Lorsqu’il s’agit de superviser un système informatique à grande échelle réparti sur plusieurs sites, il devient vite très difficile de corréler et analyser toutes les sources d’information disponibles en temps réel afin de détecter les anomalies et les incidents suffisamment vite pour réagir efficacement. Cette complexité est due à la quantité d’information générée, au manque d’interropérabilité entre les outils ainsi qu’à leurs lacunes en matière de visualisation.


\subsection {De l'usage d'un SIEM pour la gouvernance}  

À l’heure où les normes et certifications de cyber-sécurité sont de plus en plus nombreuses, le SIEM devient un élément clé de tout système d’information. C’est un moyen relativement simple de répondre à plusieurs exigences de sécurité (Exemple : historisation et suivi des logs, rapports de sécurité, alerting, …) et de prouver sa bonne foi aux autorités de certification ou de suivi. D’autant que le SIEM peut générer des rapports hautement personnalisables selon les exigences des différentes réglementations. 

Ce seul bénéfice suffit à convaincre des organisations de déployer un SIEM. La génération d’un rapport unique traitant tous les événements de sécurité pertinents quelle que soit la source des logs (générés en outre dans des formats propriétaires) fait gagner un temps précieux. En outre la création de tableaux de bord issus des usages de surpervision ou de surveillance du SIEM offre des outils de pilotage de la sécurité. 

\subsubsection{Supervision}
En collectant, corrélant et analysant l’activité du système d’information, le SIEM peut présenter des indicateurs clés de performance (KPI) et des indicateurs clé de risque (KRI).
Les fonctions du SIEM ne sont pas totalement dédiées à la sécurité, on peut donc remonter des informations plus générales concernant l’usage du système d’information. Par exemple, l’outil de source ouverte Elastic Stack, peut être utilisé pour faire du «Business Intelligence» ou du «Big Data». A travers les indicateurs dynamiques et personnalisables, on peut développer un outil de gouvernance du système d’information. Par exemple, connaître l’utilisation de ses boîtes mails, la quantité de pièces jointes échangées comparativement à un service de partage de fichiers.
\subsubsection{ Surveillance}
Le SIEM est capable de faire de la surveillance générale. Le but est de définir des règles de corrélation en capacité de déclencher une alerte ou des règles spécifiques à destination des équipes techniques du SOC. Ces règles peuvent se baser sur l’analyse des menaces, sur des comportements anormaux, ou sur un indicateur de véracité d’une attaque. En termes de surveillance, l’atout du SIEM est de centraliser la gestion des alertes sur une seule interface. 
L’autre atout du SIEM sur les solutions de surveillance d’infrastructure classique, c’est qu’il ne nécessite pas directement d'agent spécifique. l'arrivée de l'analyse des terminaux pour compléter l'analyse réseau  modifie les cultures et offre de nouvelles manière de détecter des codes malveillants en action. 


\subsubsection{Les contreparties du SIEM}
Déployer un SIEM ne suffit pas pour autant à sécuriser complètement votre organisation. Les solutions SIEM présentent des limites qui les rendent inefficaces sans un accompagnement à la hauteur et sans solutions tierces. Contrairement à une solution de sécurité de type IDS ou Firewall, un SIEM ne surveille pas les événements de sécurité mais utilise les données de logs enregistrées par ces derniers. Il est donc essentiel de ne pas négliger la mise en place de ces solutions. 


\subsubsection{Une configuration pointue }
Les SIEM sont des produits complexes qui appellent un accompagnement pour assurer une intégration réussie avec les contrôles de sécurité de l’entreprise et les nombreux hôtes de son infrastructure. 

Il est important de ne pas se contenter d’installer un SIEM avec les configurations du constructeur et/ou par défaut, car elles sont souvent insuffisantes. Les configurations doivent être personnalisées et adaptées aux besoins des utilisateurs. De même concernant les rapports, mieux vaut créer ses propres rapports d’analyse, adaptés aux différentes menaces identifiées. À défaut, le risque est réel de ne pas pouvoir profiter des avantages d’une solution de SIEM. 

\subsubsection{Des investissements à bien anticiper }

La collecte, le stockage et l’analyse des événements de sécurité sont des tâches qui semblent relativement simples. Cependant, leur collecte, stockage et l’exécution des rapports de conformité, l’application des correctifs et l’analyse de tous les événements de sécurité se produisant sur le réseau d’une entreprise n’est pas trivial. Taille des supports de stockage, puissance informatique pour le traitement des informations, temps d’intégration des équipements de sécurité, mise en place des alertes… L’investissement initial peut se compter en centaines de milliers d’euros auquel il faut ajouter le support annuel. 

Intégrer, configurer et analyser les rapports nécessite la compétence d’experts. Pour cette raison, la plupart des SIEM sont gérés directement au sein d’un SOC souvent externalisé. Porteur de grandes promesses, le SIEM mal configuré peut apporter son lot de déceptions. Selon un sondage réalisé auprès de 234 entreprises (Source LeMagIT), 81 \% d’utilisateurs reprochent aux SIEM de produire des rapports contenant trop de bruit de fond et pour 63\% les rapports générés sont difficiles à comprendre. Faire appel à des prestataires externes disposant de l’expertise dans le domaine reste souvent la meilleure solution.  

%TODO Voir SIEM externalisation ou chapitre sur la compétence

\subsubsection{Un grand volume d’alertes à réguler }

Les solutions SIEM s’appuient généralement sur des règles pour analyser toutes les données enregistrées. Cependant, le réseau d’une entreprise génère un nombre très important d'alertes (en moyenne 10000 par jours) qui peuvent être positives ou non. En conséquence, l’identification de potentielles attaques est compliquée par le volume de logs non pertinents. 

La solution consiste à définir des règles précises (en général rédigées par un SOC) et le périmètre à surveiller que faut-il surveiller en priorité ? Le périmétrique ? L’interne ? Réseau/système/application ? Quelle technologie à prioriser ? etc. 

\subsubsection{Une surveillance à exercer 24h/24}

Pour fonctionner correctement, les solutions SIEM nécessitent une surveillance 24h/24 et 7j/7 des journaux et des alertes. Un personnel formé ou une équipe dédiée sont requis pour consulter les journaux, effectuer des examens réguliers et extraire les rapports pertinents.  
 Il s’agit à la fois de disposer des expertises requises, de gagner en lisibilité budgétaire et, aussi, de profiter d’engagements de services. Des conditions à réunir afin que l’investissement dans une solution SIEM marque une étape clé dans la protection de votre organisation contre les menaces avancées.

\subsubsection{Analyse d'impact}
% la complexiste de l'analyse d'impact (si on n'a pas fait le travail <risk vers scneario, il est complexe de remonter de scénario à risque et donc impact sur l'entreprise

% citer les travaux de Lagagec

Un autre problème majeur dans l'usage d'un SIEM est que l'action de comprendre l’impact réel d’une vulnérabilité ou d’une alerte IDS est généralement dévolue à un analyste cybersécurité humain, qui doit lui- même faire le lien entre toutes les informations techniques et sa connaissance de tous les services ou processus liés aux incidents de sécurité détectés sur les composants concernés (serveurs, PC, smartphone, IOT,...) . 

Le projet DRA est une étude complémentaire de CIAP qui vise à fournir une analyse de risque en temps réel, afin de déterminer automatiquement l’impact réel dû à la situation sécurité globale du système et du réseau. Pour cela une nouvelle méthodologie innovante a été développée en combinant un générateur automatique d’arbres d’attaque (attack trees/graphs) et un moteur d’analyse de risque \g{traditionnel}  similaire à \gls{aEBIOS}.

Les systèmes de gestion des informations et événements de sécurité (SIEM) font régulièrement l’objet de critiques acerbes. Complexité, besoins importants en ressources de conseil externes… de nombreuses entreprises ont été déçues par leur expérience du SIEM pour l’implémentation de la supervision de la sécurité.

Mais la technologie n’est plus, désormais, la raison pour laquelle des entreprises peinent à réussir leurs implémentations de SIEM. Les principales plateformes de SIEM ont reçu de véritables transplantations cérébrales, se transformant en entrepôts de données taillés sur mesure pour fournir les performances et l’élasticité requises. Les connecteurs système et les aggrégateurs de logs, autrefois complexes et peu fiables, sont aujourd’hui efficaces, rendant la collecte de données relativement simple.

Mais il y a une limite au SIEM, comme à toute technologie s’appuyant sur des règles : le SIEM doit savoir ce qu’il doit chercher. Aucun boîtier SIEM ne pourra identifier automatiquement, comme par magie, une attaque tirant profit d’une méthode ou d’une vulnérabilité inédite.

Le SIEM joue un rôle important dans la détection d’attaques. Mais pour qu’il puisse détecter les attaques connues et inconnues, l’entreprise qui le déploie doit construire des ensembles de règles qui lui permettront d’identifier des conditions d’attaques et des indicateurs spécifiques à son environnement. Et le tout de manière cohérente. Comment donc construire ces règles ?

\subsubsection{Tout collecter}

\mode<all>{\picframe{Latex/Sources/EDU/Pictures/img-siem-method}{Cadre méthodologique}{0.8}{lblsiemmethod}}

Sans disposer de suffisamment de données collectées, le SIEM n’a pas grand chose à analyser. Mais la première étape est de collecter les bonnes données. Et celles-ci sont notamment les logs des équipements réseau, de sécurité  et des serveurs. Ces données sont nombreuses et faciles à obtenir. Ensuite, il faut s’intéresser aux logs de l’infrastructure applicative (bases de données, applications). Les experts du SIEM ajoutent à cela les données remontées par de nombreuses autres sources, comme celles des systèmes de gestion des identités et des accès, les flux réseau, les résultats des scans de vulnérabilités et les données de configurations.

Avec les SIEM, plus il y a de données collectées, mieux c’est. Si possible, autant tout collecter. S’il est nécessaire de définir des priorités, alors mieux vaut se concentrer sur les actifs technologiques critiques, à commencer par les équipements installés dans les environnements sensibles et ceux manipulant des données soumises à régulation, ou encore ceux touchant à la propriété intellectuelle. Le soucis principal de ces outils est souvent le modèle économique et de licence. En effet, le prix est majoritairement basé sur le nombre d'évènement.


\subsection{Construire les règles (UseCase)}

%Exemple de scénario d’attaque : Plusieurs tentatives infructueuses de connexions à un service VPN de l’entreprise. Une connexion réussie sur le VPN avec le même nom d’utilisateur. Un programme suspect installé sur la ressource depuis l‘équipement connecté au VPN. Accès continue depuis cette ressource vers un serveur de base de données. Copie des données de la base de données avec un déversement dans un fichier. Séparément les évènements sont anodins, cependant une fois corrélés, cela ressemble à une tentative d’attaque par «brute force» réussie et un vol de donnée stockée dans la base de donnée.

%TODO MENACE SIEM : Relecture

% Begin DOC`

\mode<all>{\picframe{Latex/Sources/EDU/Pictures/img-siem-usecase}{Construction des UseCase}{1}{lblsiemusecase}}


Construire une règle pour SIEM est un processus itératif. Cela signifie qu’il est relativement lent et qu’il doit être affiné, précisé au fil du temps. De nombreuses personnes sont atteintes de la « paralysie de l’analyste » en début de processus, parce qu’il existe des millions de règles pouvant être définies. Ainsi, il est  conseiller de se concentrer sur les menaces les plus présentes pour déterminer les règles à définir en premier.

Dans le cadre du processus de modélisation, il convient de commencer par un actif important. Pour cela, il faut adopter le point de vue de l’attaquant et chercher ce que l’on pourrait vouloir voler.

Modéliser la menace. Il faut se mettre à la place de l’attaquant et imaginer comment entrer et voler les données. C’est la modélisation de l’attaque, avec énumération de chaque vecteur avec le SIEM. Et il convient de ne pas oublier l’exfiltration car sa modélisation offre une opportunité supplémentaire de détecter l’attaque avant que les données ne se soient envolées. Dans ce processus, il s’agit d’adopter des attentes réalistes car le modèle d’attaque ne peut pas par essence être parfait ni complet. Mais il convient toutefois d’engager le processus de modélisation. Et il n’y a pas de mauvais point de départ.

Affiner les règles. Il convient ensuite de lancer l’attaque contre le SI, telle que modélisée. Les outils pour cela ne manque pas. C’est l’occasion de suivre ce que fait le SIEM. Déclenche-t-il les bonnes alertes ? Au bon moment ? L’alerte fournit-elle suffisamment d’informations pour assister les personnes chargées de la réaction ? Si l’alerte n’est pas adéquate, il convient de revoir le modèle et d’ajuster les règles.

Optimiser les seuils. Avec le temps, il deviendra de plus en plus clair que certaines alertes surviennent trop souvent, et d’autres pas assez. Dès lors, il convient d’ajuster finement les seuils de déclenchement. C’est toujours une question d’équilibre… un équilibre délicat.

Laver, rincer, recommencer. Une fois l’ensemble initial de règles pour ce modèle d’attaque spécifique implémenté et optimisé, il convient de passer au vecteur d’attaque suivant, et ainsi de suite, en répétant le processus en modélisant chaque menace.

Ce processus ne s’arrête jamais. Il y a constamment de nouvelles attaques à modéliser et de nouveaux indicateurs à surveiller. Il est toujours important de suivre les informations de sécurité pour savoir quelles attaques sont en vogue. Les rapports tels que celui de Mandiant sur le groupe APT1 intègrent désormais des indicateurs clairs que chaque organisation peut surveiller avec son SIEM. Armé de ces renseignements sur les menaces et d’un environnement de collecte de données complet, il n’y a plus d’excuse : il est temps de commencer à chercher les attaques avancées qui continuent d’émerger.

Mais avec le temps, il sera nécessaire d’ajouter de nouveaux types de données au SIEM, ce qui impliquera de revoir toutes les règles. Par exemple, le trafic réseau, s’il est capturé et transmis au SIEM, fournira quantité de nouvelles informations à étudier. Mais comment ce regard sur le trafic réseau sera-t-il susceptible d’affecter la manière dont certaines attaques sont traitées ? Quelles autres règles faudrait-il ajouter pour détecter l’attaque plus vite ? Ce ne sont pas des questions triviales : il convient de revoir les règles du SIEM chaque fois qu’est ajoutée une nouvelle source de données (ou retirer, le cas échéant) ; cela peut faire la différence sur la rapidité avec laquelle une attaque est détectée… si elle l’est.

Le plus important aspect de ce processus est la cohérence. Le SIEM n’est pas une technologie du type « installe et oublie ». Il requiert du temps, de l’attention, et d’être alimenté, tout au long de sa vie opérationnelle.

% End DOC


\subsection{Quelques défis des SIEM	}

La problématique globale des SIEM est de corréler de l'événement, la question de fond est la collecte de ses évènements. La collecte de LOG est la principale source d'événements, toutefois, toute les sources d'événements sont susceptibles d'enrichir la corrélation, en particulier les vulnérabilités, les IOC, les infos de end-point ...
La collecte des informations d'opérations et de renseignements nécessite de la FUSION de capteurs.  Cette fusion chère aux militaires est un premier pas qui prend compte aussi de l'information économique, politique ou sociale de l'entreprise, car ces événements peuvent \g{matcher} avec des attaques complexes.
  
\subsubsection{L'intelligence artificiel}

%Le traitement de masse permet 
%Mais l'important est que l'IA (ex: vectra) puisse explicquer ses propositions et décisions. 

L’Intelligence artificielle est devenue en quelques années un objetif marketing chez les éditeurs de cybersécurité. Avec un certain succès puisque de nombreuses entreprises disent  utiliser des solutions de sécurité basées sur de l’IA. Certaines solutions s’appuient davantage sur des moteurs de règles sophistiquées que sur de réelles fonctionnalités d’IA. Pour parler d’intelligence artificielle, il faut en effet que la technologie inclut : 

\begin{itemize}
  \item  une capacité de perception de l’environnement au moyen d’un apprentissage supervisé ou non,
  \item  une capacité d’analyse et de résolution de problème, 
  \item   une capacité de proposition d’action, voire de décision autonome.
\end{itemize}

Sur le plan théorique, les apports de l’IA en matière de cybersécurité sont donc nombreux, qu’il s’agisse de prévention, d’anticipation, de détection ou de réaction. Dans la pratique, la détection de vulnérabilités ou de menaces internes ou externes apparaît aujourd’hui l’un des usages les plus matures.  Les systèmes actuels basés sur des signatures montrent encore leurs limites : nombre élevé de faux positifs, incapacité à s’adapter aux dernières menaces, notamment aux APT, lourdeur des bases de signature, ce qui a un impact sur les performances. 

\begin{techworkbox}{SIEM et IA}
L'évolution des SIEM est orientée par le traitement de masse d'évènement. L'IA du BigData offre des possibilités nouvelles.  L'IA dans la détection d'attaque est un bon sujet pour \fichetech.
\end{techworkbox}

\subsubsection{Quelques SIEMs}

On peut citer  ainsi quelques SIEM non pas pour en faire une publicité particulière mais simplement pour donner quelques indications sur la provenance ...

% Begin PRZ ===========================
\begin{frame}
\frametitle<presentation>{Quelques SIEM}
\framesubtitle<presentation>{en provenance du Gartner}
% end header PRZ =======================
Le Gartner positionne régulièrement des produits et services dans son magic Quadrant. en 2019, Splunk, IBM QRadar et LogRythm NextGen SIEM sont toujours bien positionnés.
Dell Technologies (RSA NetWitness), Exabeam (Security Management Platform), McAfee (Enterprise Security Manager) et Securonix complètent le carré des Leaders.
Toutefois des entreprises comme Microsoft challengent ces acteurs.
\end{frame}
% end PRZ ===========================


\toolsbox{RSA-Netwitness}{siem}
\toolsbox{Qradar}{siem}


\subsection{Threat Hunting}

\subsubsection{la chasse aux menaces}

La chasse aux menaces dormantes ou aux compromissions (Threat Hunting) fait partie intégrante de la gestion de la menace. 
C'est souvent les équipes \g{Hunters} qui assurent le maintien du contact entre la défense et les attaquants lors d'une attaque en cours. Souvent issus d'équipe de compétences liées à réponse à incident.


% Begin PRZ ===========================
\begin{frame}
\frametitle<presentation>{Threat Hunting}
% end header PRZ =======================
La chasse aux menaces est une tactique permettant de connaître avec plus d'acuité l'environnement de la menace et donc  le degré de risque de cyberattaques auquel est soumise une entreprise. 

La terminologie threat hunting regroupe plusieurs types d'action et la définition n'est pas totalement stabilisée. Globalement on y trouve deux grandes classes de \g{\textit{threat hunting}} :

\begin{itemize}
  \item Celles travaillant autour de l'environnement, de la surface d'attaque et qui orientent ses actions sur des méthodes de \g{recherche} permettant de débusquer des menaces latentes ou des menaces dormantes, de les réveiller, de les suivre , de les comprendre pour établir le contact avec l'attaquant.

  \item Et une autre plus active ou proactive dont l'objectif est de rester, conserver le contact avec l'attaquant lors d'une réaction à une alerte.
\end{itemize}


\end{frame}
% end PRZ =========================== 


\subsubsection{Etablir le contact}

Quand on parle d'établir contact, nous parlons d'aller au contact au sens martial du terme, c'est à dire en direct de suivre, caractériser la sources de la menace et jouer avec elle.

La méthode de \g{hunter} consiste en premier à dresser un portrait global de la surface d’attaque, tout en identifiant les attaquants potentiels, leurs motifs et leurs façons de faire. Plus précisément, le \g{threat hunting} consiste en une analyse détaillée de différents éléments :

\begin{itemize}
	\item la position de  l’entreprise,  notoriété, popularité sur internet, en analysant en particulier,  les médias traditionnels et les médias sociaux;
	\item l'environnement économique de l’entreprise dont ses fournisseurs, ses clients, ses partenaires, ses employés;
	\item le corpus technologique et physique de l'entreprise, dont les architectures techniques et les mécanismes informatiques avec l'environnement économique ainsi que l'environnement sécuritaire de ses relations.
\end{itemize}

Sur la base de cette analyse globale, des SPOF (\textit{Sigle Point Of Failure}) peuvent être trouvés. 

Grace à la visualisation  globale des liens il est possible  de comprendre où, comment, pourquoi et potentiellement par qui (activistes, anciens employés, fournisseurs, etc.) la prochaine attaque pourrait être perpétrée. 
Les « threat hunters », ne sont pas simplement en attente de répondre aux alertes du système de défense, ils cherchent activement des menaces dans leurs propres réseaux afin de prévenir ou de minimiser les dommages. Cette méthode s’avère l’une des plus proactives. 


}
\mode<all>{%-------------------------------------------------------------
%               FR CYBERDEF SECOPS COURSE
%              $Chapitre : Threat Management
%                   $theme : Threat Detection
%               $File :  chap-ThreatMan-edr.tex
%                             2020 eduf@ction
%-------------------------------------------------------------
\uchap{chap-ThreatMan-siem.tex}
%-------------------------------------------------------------


\subsection{la détection de menace sur des terminaux}

La mise en place de mesures de remontées de LOGs pour les terminaux fait partie généralement des politiques de cyberdéfense, et les journaux d'évènement au sein des systèmes d'exploitation ou des applications sont utilisables pour les SIEM.
La sécurité du poste de travail est un axe à part entière des politiques de sécurité. Dans cette sécurité, on y trouve évidement la sécurité périmètrique avec les \textit{Firewalls} personnels, les anti-virus. La détection de menace au plus prêt du terminal et donc au plus prêt de l'utilisateur, permet d'être, bien des fois, plus pertinent pour contenir des attaques ou des déviances. Cette détection peut être réalisée par un EDR (\textit{Endpoint Detection and Response}) qui est caractérisée par ses capacités de détection, d’investigation et aussi de remédiation.

Etre au plus prêt de l'utilisateur, donc du terminal, est toujours un axe majeur à pour un responsable sécurité.

Chaque jour des milliers de terminaux (\textit{Endpoints}) sont compromis par des attaques ciblées. Parmi ces attaques, les menaces avancées et persistantes (APT : Advanced Persistent Threats) compromettent des machines et avec les risques associés. 

Des techniques d'attaques permettent de by-passer les solutions locales d'antivirus et de firewall personnel ou d'utiliser plus simplement la crédulité des utilisateurs pour infecter des machines ou tout simplement utiliser les faiblesses de configuration du systèmes pour ex-filtrer des données.

Dans ce contexte l’EDR est un outil de sécurité complémentaire aux outils personnels de sécurité avec lequel il travaille afin de bloquer les menaces (connues ou nouvelles (zero-days). C’est un outil qui se place au niveau du système d'exploitation en complément du niveau réseau ...  Un EDR fait de l’analyse comportementale, et permet de monitorer les actions de l'utilisateur ou des applications au niveau  terminal et de réagir.

\subsubsection{Détection}

L’EDR est capable de surveiller l’exploitation de failles de sécurité en surveillant les appels noyaux et les différents services habituellement ciblés notamment chez Windows. Cette capacité de surveillance et la corrélation d’évènements lui permettent de reconnaître des méthodes et habitudes qu’ont les hackers et dont il est plus difficile de se prémunir.

L’analyse comportementale (UBA - User Behavior Analytic) est un autre point supporté par les EDR et qui permet de reconnaître des comportements déviant d’un cadre normatif, ou technique souvent après une phase d’apprentissage. Grâce à ces analyses, l’EDR peut émettre des alertes vérifiable qui renforceront l’apprentissage. 

 L’intérêt de cette technique est qu’elle permet de stopper un attaquant dans son élan : si un PDF contient un script qui ouvre powershell et ouvre une connexion sur un port classique d’un serveur extérieur au SI alors cette suite d’action sera considérée comme anormale et va être bloquée par l’EDR. Cette visibilité est une grande force de l’EDR car elle permet une remédiation à la source de l’infection.


\subsubsection{Investigation}

Comme mentionné plus haut, l’EDR permet d’observer des suites d’actions dont le résultat est des plus douteux. Cette visibilité sur les processus est d’une très grande aide pour faire de l’investigation : les actions sont corrélées et remontées dans une plateforme centralisée qui permet d’étendre l’apprentissage observé d’un poste à tous les autres. 

Grâce à ces plateformes, un \textit{Security Operation Center} (SOC) est capable de savoir immédiatement combien de postes sont touchés et d’en remonter la piste : c’est un outil d’investigation qui peut s'interfacer à  un \textit{Security Information and Event Management} (SIEM) en offrant de la visibilité sur les terminaux. D’autant plus que votre SOC pourra se servir de l’EDR pour récupérer des artefacts de l’attaques à distance.

\subsubsection{Remédiation}

En ce qui concerne la remédiation, l’EDR a des capacités similaires à celles d’un antivirus de nouvelle génération et peut notamment bloquer, supprimer et mettre en quarantaine des fichiers. Les équipes de sécurité pourront aussi s’appuyer sur l’outil pour faire du nettoyage de clé de registre par exemple, ou patcher la mémoire en direct pour contrôler un malware.  De plus, certains EDR permettent aux analystes du SOC d’Orange Cyberdefense de prendre la main à distance sur un terminal qui nécessiterait une investigation plus poussée encore.

L’EDR n’est généralement pas une solution stand-alone. C’est un  complément qui s’intègre à d'autres outils de sécurité locale, de SIEM ou de sécurité réseau. Il permet d’étendre la visibilité sécurité du SI jusqu’aux terminaux et permet d'améliorer la pertinence et la précision de détection des scénarios.

\subsubsection{Quelques EDR de références}

\toolsbox{microsoft}{edr}
}
\mode<all>{%-------------------------------------------------------------
%               FR CYBERDEF SECOPS COURSE
%              $Chapitre : Threat Management
%                         $theme : Caracterization
%              $File : chap-ThreatMan-caract.tex
%                             2020 eduf@ction
%-------------------------------------------------------------
\uchap{chap-ThreatMan-caract.tex}
%-------------------------------------------------------------

\subsection{Caractérisation de la menace}

Nous avons évoqué les mécanismes de surveillance et de détection constitués en particulier d'outils comme le SIEM, ou l'EDR.
Il est toutefois important de disposer de compétences et de processus structurés pour caractériser la menace.
Cette caractérisation permet de lever les doutes, de caractériser les éléments d'une attaque dont les impacts mais surtout les raisons et les intentions des attaquant. Ces éléments sont évidement primordiaux pour définir une réponse adaptée.

\subsection{De l'usage d'un CSIRT}

Nous verrons l'usage d'un CSIRT dans le chapitre sur la réponse à incident toutefois nous pouvons donner ici les fonctions ou compétences clefs pour effectuer cette caractérisation :

\begin{itemize}
  \item Rétro-conception : afin d'analyser un code malveillant;
  \item Liaisons police/gendarmerie : afin de lancer au plus tôt les investigations via les services de l'état;
  \item Base \g{enrichie} d'IP \& Url malveillantes : afin de caractériser non seulement la provenance, mais aussi l'intentionalité de l'attaque;
  \item Etudes de cyber-intelligence sur le secteur d'activité de l'entreprise : afin de corréler une agression avec des risques dans un environnement politique ou économique particulier.
\end{itemize}
}
\mode<all>{%-------------------------------------------------------------
%               FR CYBERDEF SECOPS COURSE
%              $Chapitre : Threat Management
%                         $theme : Caracterization
%              $File : chap-ThreatMan-caract.tex
%                             2020 eduf@ction
%-------------------------------------------------------------
\uchap{chap-ThreatMan-csoc.tex}
%-------------------------------------------------------------

\section{(SOC) Security Operation Center}

Le SOC (Security Operation Center) est au coeur du système de \g{Veille Alerte et réponse}. C'est une tour de contrôle sécurité de l'espace Cyber.

Il est constitué généralement d'une équipe d'analystes, et d'outils permettant de surveiller l'environnement. Cette surveillance s'effectue sous la forme de l'exploitation de différents outillages (SIEM, EDR....).

%TODO : réaliser un schéma de la fonction SOC

% Begin PRZ ===========================
\begin{frame}
\frametitle<presentation>{SOC}
\framesubtitle<presentation>{Security Operation Center}
% end header PRZ =======================

Il intègre l'ensemble des fonctions liées à la menace :

\begin{itemize}
  \item Veille sur la menace
  \item Détection d'évènements à risques et gestion de ceux ci
  \item Détection d'attaques ou de comportements critiques
  \item Réaction aux incidents et remédiation
\end{itemize}

\end{frame}
% end PRZ ===========================

Malheureusement, dans encore beaucoup de cas, les équipes SOC et les équipes liées à la gestion des vulnérabilités sont cloisonnées, ce qui ne couvre pas de manière intégrée l'ensemble des fonctions de cyberdéfense d'entreprise.

On peut aussi Intégrer dans  le SOC des fonctions de  \textit{Threat Hunting }.

Les grands principes de réussite d’un SOC sont : 

\begin{itemize}
  \item Une veille Cyber efficace et à large ouverture en terme de menaces;
  \item Une capacité à identifier d’une cartographie détaillée des ressources de l’infrastructure et de correctement identifier les menaces avec des analyses des risques;
  \item Réaliser une collecte des évènements de sécurité, pour nourrir une corrélation temps réel;
  \item Contextualisation et amélioration continue afin de limiter le nombre de faux-positifs;
  \item Faciliter la communication entre les niveaux opérationnels
\begin{itemize}
\item Niveau 1 : \textbf{réception} des alertes en temps réel;
\item Niveau 2 : \textbf{corrélation} et analyse multi-alertes pour déclenchement de l’incident;
\item Niveau 3 : \textbf{investigations} poussées, forensiques et découverte des indicateurs de compromission.
\end{itemize}
\end{itemize}


\subsubsection{Le SOC de demain}
On peut par ailleurs s'interroger sur le fait qu'un tel système peut et doit opérer d'autres missions que les missions de sécurité pures. Si la supervision des réseaux a été longtemps aux outils au services des techniciens, la supervision de l'environnement digital c'est à dire l'environnement informationnel de l'entreprise est un axe fondamental. Le SOC peut devenir \textUK{Cybersecurity Operational Center} (CSOC) opérant le suivi des risques digitaux au sens large, incluant les réseaux sociaux et leur cohorte de fausses informations et d'information pouvant être des indicateurs de crise à venir pour l'entreprise.

Aujourd’hui les entreprises organisent donc leur environnement de gestion de la menace dans la sécurité informatique avec ce centre opérationnel de sécurité CSOC. Ses principales missions se structurent donc autour de  : 

\begin{itemize}
  \item la supervision de la sécurité;
  \item du management du risque ;
  \item de l’analyse des menaces;
  \item de l’audit;
  \item de l’investigation numérique;
  \item de la prévention.
\end{itemize}

Chacune de ces missions nécessite des spécialistes, des outils et une gouvernance. Cependant au centre de ces activités  le CSOC s’appuie entre autres sur des données de référence
\begin{itemize}
  \item Les journaux d'activité;
  \item des scénarios menaces; mis à jour avec des d'outils de surveillance d'émergence de la menaces, en plus de scénarios spécifiques à l'entreprise;
  \item des bases de vulnérabilités et des accès à des services permettant;
  \item des outil de surveillance des fuites (Data Leak Detection) dans le Darkweb par exemple;
  \item des outils de remédiation, de notarisation (enregistrement des données à valeur probante).
\end{itemize}

Le SIEM et l'EDR répondent à une partie des besoins des CSOC et doit pour être efficaces s'intégrer à des services connexes.
Le SIEM existe depuis les années 90, sous la forme de SIM et de SEM. Le rôle du SIM est de centraliser tous les logs dans une seule base de données pour permettre des analyses et de l’archivage. Celui du SEM est de surveiller en temps réel les évènements sur le système d’information, les corréler et alerter si les conditions sont remplies. Ces deux outils de sécurité ont fini par fusionner dans un seul produit, permettant de lier les évènements aux informations, mais il y a encore des convergences à venir.

\subsubsection{Evaluation d'un SOC}

L'efficacité d'un SOC peut être évaluée. A l'image d'équipe de Pentest qui testent la résistance d'un système, des équipes de tests de SOC peuvent être déployées pour auditer le niveau d'efficacité d'un SOC.
Les équipes qui testent des OSC sont nommées des \textbf{\textit{Purples Team}}.
%SOURCE https://www.globalsecuritymag.fr/SOC-comment-en-mesurer-l,20151005,56416.html

Un CSOC est efficace s’il arrive à détecter avec pertinence les attaques en cours,  Cependant, il est important, de comprendre qu’un CSOC ne détectera jamais  des attaques dont le scénario n'a pas été pensé / \g{programmé}. On trouvera dans \ulink{https://clusif.fr/publications/reussir-deploiement-dun-soc/}{une publication du  CLUSIF}, les critères pour réussir le déploiement d'un SOC.

\subsubsection{Les outils connexes d'un SOC}

Au delà des SIEM, il semble important d'ajouter à l'outillage d'un SOC un ensemble de systèmes permettant de mesurer et d'évaluer l'impact des attaques. Un travail intéressant autour de la notion d'Echelle de RICHTER (Voir  \ulink{https://observatoire-fic.com/prendre-la-mesure-des-cyberattaques-peut-on-definir-une-echelle-de-richter-dans-le-cyber}{un article du FIC 2014}) d'une attaque afin de définir des indicateurs \g{de cotation}.

% Begin PRZ ===========================
\begin{frame}
\frametitle<presentation>{Cotations connexes	}
\framesubtitle<presentation>{gérer le niveau de gravité de l'alerte}
% end header PRZ =======================
\begin{itemize}
  \item \textbf{l’origine} de l’attaque qui mesure la puissance potentielle de la source de menace : du hacker de base à la menace étatique ;
  \item Le type de \textbf{cible} qui mesure la précision de la diffusion de la menace : de la cible au hasard à la menace ciblée ;
  \item Le \textbf{vecteur} d’attaque qui mesure le niveau de sophistication de la menace : du malware « sur étagère » à l’APT élaborée ;
   \item Le \textbf{préjudice} qui mesure l’impact subit par la cible : d’une perte faible à une mise en péril de la résilience même de l’organisme ;
  \item La \textbf{visibilité} de la menace qui mesure de nombreux éléments comme la motivation ou durée de l’attaque : d’un DDOS immédiatement constaté à une attaque invisible ;
  \item La \textbf{persistence} qui mesure la fréquence de l’attaque sur sa cible : d’une fréquence forte de type robotisée (Bots) à  une fréquence unitaire visant un but précis, ou la furtivité.
\end{itemize}
\end{frame}
% end PRZ ===========================

\subsection{Les outillages d'un SOC}
% $sources : SOURCE https://observatoire-fic.com/prendre-la-mesure-des-cyberattaques-peut-on-definir-une-echelle-de-richter-dans-le-cyber/

Au delà des SIEM, des sondes, des EDR, des XDR, l'orchestration est au coeur de l'efficacité des fonctions d'un CSOC en particulier pour l'automatisation de la réponse à incidents.

L'orchestration et l'automatisation de la sécurité  permettent de réduire les délais de réponse, de  limiter l'exposition aux attaques et offrir une cohérence des processus cyberdéfense. Ces  outils d’automatisation et d’orchestration,  appelés SOAR, sont conçus pour améliorer la productivité et l’efficacité de centres des opérations de sécurité et des analystes.

Ces outils automatisent les tâches de routines chronophages,  ils aident à coordonner les cycles de vie de réponse aux incidents et de gestion des incidents. Outils de  cohérence, ils permettent d'assurer reproductibilité de la discipline des opérations de cybersécurité et permet de réduire le temps nécessaire pour détecter et traiter les incidents.

\subsubsection{Qu'est-ce que la détection-réponse des Endpoints (EDR) ?}

La détection et la réponse Endpoints (EDR) est une solution de cybersécurité qui capture toute l'activité des Endpoints et s'appuie sur des analyses avancées pour fournir une visibilité en temps réel de la santé de tous ces terminaux ; détecter les activités anormales ; alerter l'équipe de sécurité (SOC par exmple) des événements ; et fournir des suggestions et des capacités de remédiation pour répondre, arrêter une attaque en cours ou limiter sa propagation.

Les solutions de détection et de réponse des Endpoints ont les capacités suivantes :

\begin{itemize}
  \item Surveillance des points d'extrémité et enregistrement des événements
  \item Recherche de données, investigation et chasse aux menaces
  \item Triage des alertes ou validation des activités suspectes
  \item Détection d'activités suspectes
  \item Analyse des données
  \item Renseignements exploitables à l'appui de la réponse
  \item Remédiation
\end{itemize}


\subsubsection{Qu'est-ce que la détection-réponse étendues (XDR)}


La détection-réponse dite étendue (XDR) rationalisent la collecte, l'analyse et les flux de données de sécurité sur l'ensemble de l'infrastructure de sécurité d'une entreprise, améliorant ainsi la visibilité des menaces de sécurité cachées et avancées et unifiant la réponse.

Une plateforme XDR collecte et met en corrélation les données de l'ensemble de l'infrastructure afin d'améliorer la visibilité des menaces dans toute l'entreprise, d'accélérer les opérations de sécurité et de réduire les risques. XDR analyse, hiérarchise et rationalise ces données, afin qu'elles puissent être transmises aux équipes de sécurité dans un format normalisé via une console unique et consolidée.

Les plateformes XDR offrent généralement les fonctionnalités suivantes :

\begin{itemize}
  \item Télémétrie de sécurité diversifiée et multi-domaine
  \item Analyse des événements axée sur les menaces
  \item Détection des menaces et hiérarchisation de la fidélité des données
  \item Recherche de données, investigation et chasse aux menaces dans la télémétrie multi-domaine.
  \item Réponse pour atténuer la menace et y remédier
\end{itemize}


\mode<all>{\texframe{Automatisation, SOAR}{Orchestration}
{
On y trouve par exemple dans ces outils de  \g{\textit{Security Orchestration, Automation, and Response}} (SOAR) :
\begin{itemize}
  \item l'introduction de sources de menaces de manière automatique au base SIEM (abonnement de threat-intelligence);
  \item la production de règles sur la base de déviances relevées;
  \item le pilotage automatique des composants de sécurité (modification de règles, passage en mode dégradé ...);
  \item l'exécution de tâche de conservation de traces légales (notarisation);
  \item la gestion automatisée de \g{patchs} critiques (intégration au DEVSECOPS) ...
\end{itemize}
}} % end frame

\toolsbox{Resilient}{soar}
\toolsbox{Demisto}{soar}

% TODO https://www.splunk.com/fr_fr/form/the-soar-buyers-guide.html

\subsubsection{DEVSECOPS}
% $source : https://www.redhat.com/fr/topics/devops/what-is-devsecops

Dans de nombreuses entreprises les processus liés à la sécurité sont encore isolés et confiés à une équipe spécifique sans s'intégrer totalement avec les chaînes de développement ou dans les équipes opérationnelles et encore moins dans les chaînes intégrées DEVOPS. 
 Si une approche DevOps efficace garantit des cycles de développement rapides et fréquents, les équipes sécurité doivent de plus en plus s'intégrer dans les processus DEVOPS tant pour y apporter un volet \textit{security by design}, mais surtout intégrer des mécanismes de sécurité opérationnelle.
 Ce sont ces dynamiques de sécurité tant de conception, que de sécurité opérationnelle que nous appelons \g{DEVSECOPS}.
 
\utocomplete
%TODO DESECOPS à compléter

 
\begin{techworkbox}{DEVSECOPS}
Le domaine du DEVSECOPS est un sujet à part entière de la gestion de la sécurité dans les chaînes de prise en compte de la sécurité  en DEVOPS, et l'utilisation des techniques DEVOPS dans les chaînes de gestion du dynamique du changement dans la détection de menace et la réponse sur incident. Un sujet intéressant pour \fichetech.
\end{techworkbox}
 
 \subsection{l'efficacité du CSOC}
 % $source : https://www.excube.fr/ressources/soc/mesurer-efficacite-soc
 
\mode<all>{\texframe{Efficacité du CSOC}{métrologie}
 {
 Pour mesurer l’efficacité d'un CSOC, il existe plusieurs moyens de mesures  :
 \begin{itemize}

  \item La \textbf{couverture fonctionnelle et technique }du CSOC pour estimer l'efficacité  de l'articulation entre les stratégies de cyberdéfense, de cyber-protection et les stratégies de surveillance détection, 
  \item La \textbf{performance de la détection } pour évaluer l’efficacité des règles de corrélation en place, basée sur les indicateurs de services (Nombre de détections, nombre d'évènement ...);
  \item La \textbf{maturité du service CSOC}, mesurée sur le niveau d'organisation des services (ITIL par exemple), les coûts, les compétences, les services connexes ...
\end{itemize}
}} % end frame


}
\mode<all>{%-------------------------------------------------------------
%               FR CYBERDEF SECOPS COURSE
%              $Chapitre : Threat Management
%                         $theme : Leak Management
%              $File : chap-ThreatMan-leak.tex
%                             2020 eduf@ction
%-------------------------------------------------------------
\uchap{chap-ThreatMan-leak.tex}
%-------------------------------------------------------------

% SOURCE : LECTURE : https://www.groupjoos.com/gdpr-17-que-faire-en-cas-de-fuite-de-donnees-gestion-des-incidents/?lang=fr


\subsection{Leak  : surveiller les fuites}

J'ai ajouté un chapitre spécial sur les fuites de données pour deux grandes raisons :

\begin{itemize}
  \item La détection des fuites de données peut simplement se révéler par l'apparition de tout ou partie de ces données dans le Darkweb. 
  \item Les fuites de données étant souvent des fuites de données de type \g{données personnelles}, elles impliquent le déroulement de processus de déclaration au titre de la GDPR.
\end{itemize}


Je ne rentrerai pas ici dans la présentation du RGPD avec son cortège d'exigence et d'organisation à mettre en place (
Liste de traitement, déclaration, nomination de responsable, etc).
Je ne vous propose que de regarder rapidement, la partie détection et partie réponse à incident.

Le terme « fuite de données », ou « data breach » en anglais, est utilisé pour toute situation impliquant la perte, la modification injustifiée ou la publication par accident, par malveillance, de données considérées ou marquées comme confidentielles. 

Il est important dans la mise en place de scénario dans les SIEM, et dans le traitement de SOC que l'évènement de fuites de données personnelles puissent être traitées avec un mécanisme précis et documenté, car ces évènements sont très contraints par la réglementation.
A titre de remarques, les évènements touchant la fuite de données liées à la protection du secret de défense (Secret Défense) puissent aussi être traitées dans un processus particulier car les ces fuites peuvent aussi faire l'objet de procédure au pénal.

% TODO MENACE LEAK : A relire  (passage en incidents ?)
%*************************************************

Le GDPR prévoit que le responsable du traitement des données à caractère personnel signale au plus tôt les fuites de données  pouvant constituer une atteinte à la vie privée des personnes concernées. Cette information à la \gls{aCNIL} et aux personnes concernées en cas d'impact important sur ces personnes.

La méthodologie est assez simple pour peu que le constat de l'incident puisse être fait le plus vite possible. Cela peut se faire sur la base d'évènement provenant des équipements de sécurité (via un SIEM par exemple) ou par l'utilisation de services de veille, ou simplement par l'avertissement d'un tiers qui découvre cette fuite. 

\begin{itemize}
  \item Détection,
  \item Enrayer la fuite, limiter l'impact,
  \item Analyser les sources de menaces,
  \item réagir de manière juridique.
\end{itemize}


Deuxièmement, vous devez entreprendre dès que possible les démarches pour enrayer l’incident ou en limiter l’impact. Tous les collaborateurs doivent respecter plusieurs règles. S’ils trouvent des informations à un endroit inapproprié, ils doivent les supprimer ou en informer un responsable. Il peut s’agir de supports physiques, mais aussi de fichiers sur le réseau. Ils doivent également donner l’alerte s’ils rencontrent des étrangers non accompagnés dans une zone sécurisée. Et ainsi de suite. Si des alarmes indiquent un piratage ou une infection des systèmes, les gestionnaires de ces systèmes devront les examiner au plus vite et peut-être les désactiver de manière préventive.

En cas de doute, il est préférable d’arrêter un traitement ou d’empêcher le transport des données traitées jusqu’à ce que vous sachiez clairement s’il y a effectivement un problème, et dans quelle mesure les données traitées sont encore correctes. Cela permet souvent d’éviter qu’un incident ne se transforme en fuite de données. Tant que des données traitées à mauvais escient ne sont pas diffusées ou rendues publiques, il n’y a pas d’infraction, et donc pas d’impact. Au sens strict, il n’est pas encore question d’une fuite de données.

Ensuite, et éventuellement en parallèle, vous pouvez lancer une analyse des faits. D’une part, il faut établir la cause du problème. Vous pourrez ensuite réfléchir aux améliorations dans l’organisation, les systèmes ou les applications, et dans le mode de travail de vos collaborateurs, pour éviter que l’incident ne se reproduise. D’autre part, il faut examiner l’impact réel ou éventuel de l’incident. Y a-t-il des risques pour la confidentialité et l’intégrité des données ? S’agit-il (en partie) de données à caractère personnel ? Quelles peuvent-être les conséquences de cette infraction ? Dans de nombreux cas, il vous faudra du temps pour savoir quelle quantité de données a été impactée et combien de personnes sont concernées. Souvent, vous ne saurez pas non plus d’emblée s’il y a véritablement un risque d’impact, ni quelle peut être l’ampleur des dommages.

Ce n’est que lorsque vous aurez une réponse à toutes ces questions qu’il vous sera possible de faire le bon choix quant à la nécessité de signaler la fuite de données à la Commission de protection de la Vie Privée ou aux personnes concernées. Le quand et le comment de ce signalement seront abordés dans le prochain article.

% TODO MENACES LEAK : A relire

%Par ailleurs, chaque incident doit être consigné dans un registre interne. Qu’il s’agisse d’une véritable infraction ou d’un quasi-accident, il faut toujours analyser l’incident. Ces informations sont importantes pour évaluer les procédures et les directives existantes, et vérifier si les mesures prises offrent une protection suffisante contre les risques éventuels. Les causes d’un incident doivent être consignées, au même titre que les actions visant une amélioration. En assurant un suivi systématique, vous améliorerez systématiquement la sécurité de votre organisation.

%Dans les cas extrêmes, une fuite de données peut être catastrophique. Une organisation peut être confrontée à des problèmes de communication dantesques suite à une fuite de données très sensibles à propos d’un grand nombre de personnes. Il arrive que la fuite de données sorte des murs de l’organisation et que la presse en soit informée. En pareil cas, il est bon de pouvoir retomber sur des scénarios de communication de crise préalablement établis. Si votre organisation est couverte par une assurance couvrant les risques de cyber-sécurité, votre compagnie d’assurance devrait pouvoir vous aider.

%Si vous soupçonnez que l’incident est d’origine criminelle, vous devez veiller à constituer un dossier juridique à temps. Il est parfois important de réaliser un back-up rapide des systèmes au moment de la découverte de l’incident ou de conserver les fichiers log, avant que ces informations ne soient perdues ou modifiées par les démarches entreprises pour résoudre l’incident. Il est évident qu’une telle étape ira parfois à l’encontre de ce qu’il convient de faire rapidement pour limiter le problème existant. Si la police ou la justice intervient, ne perdez jamais de vue ce que vous pouvez et ne pouvez pas faire de votre propre chef, surtout si vous êtes en charge du traitement des données. Faites appel au responsable dès que possible. Si les autorités vous obligent à fournir des informations, vous devez toujours veiller à les protéger au mieux et à ne pas exposer de données (par exemple d’autres personnes concernées) si cela n’est pas nécessaire à l’enquête.

%Il est judicieux de bien documenter ces démarches successives, afin que chacun dans l’organisation les connaisse et agisse en fonction. Cela peut également s’avérer utile pour démontrer que vous prenez le respect des obligations du GDPR très au sérieux.

%*********************************************************************


}
\mode<all>{%-------------------------------------------------------------
%               FR CYBERDEF SECOPS COURSE
%              $Chapitre : Threat Management
%                         $theme : Divers
%              $File : chap-ThreatMan-divers.tex
%                             2020 eduf@ction
%-------------------------------------------------------------
\uchap{chap-ThreatMan-divers.tex}
%-------------------------------------------------------------

\section{Technologies et Organismes connexes}

\subsection{Stratégie de cyberdéfense et de surveillance}

Ce chapitre en construction donne quelques éléments complémentaires sur les stratégies et techniques de cyberdéfense d'entreprise. 

%\subsection{Techniques de défense cyber}

\subsubsection{Technique de déception (Deceptive Defense)}

 D'une protection statique en profondeur encore trop souvent périmètrique, vers une détection résiliante pour désormais envisager des logiques de contre-attaques dynamiques, les stratégie de cybersécurité évoluent. L'échelle de temps entre l'occurrence d'une attaque, sa détection et son élimination est un marquant significatif du risque : de plusieurs mois à quelques jours,  de jours à quelques heures, l'enjeu est  de nos jours  d'agir en temps réel contre l'attaquant. 
 Le leurrage numérique (deceptive security) est au cœur de nouvelle stratégie de  cybersécurité. Il s'agit de retourner la dynamique  de l'attaquant en cherchant à le tromper pour tenter de le démasquer, ou de le dissuader de continuer de part le risque d'être découvert qu'il prend. Le leurrage numérique peut relever d'une forme de dissuasion cyber.

\subsubsection{Honeypots}

Les pièges « honeypots » sont un leurre pour les attaquants, en imitant une ressource de calcul réel (par exemple, un service, une application, un système ou des données). Toute entité entrée en connexion à un « honeypot » est alors considérée comme suspecte, et son activité est surveillée pour détecter une malveillance.
L'arsenal du leurrage défensif est historiquement basé sur les pots de miel (honeypots). Ceux-ci reposent sur l'analyse statique d'écart de composants par rapport à un comportement connu et sain. Ces technologies se heurtent à deux problèmes : le passage à l'échelle pour couvrir la diversité et la complexité des systèmes numériques, et la génération excessive de faux positifs. Les honeypots évoluent pour devenir des pièges actifs qui sont disséminés dans l'environnement réel pour mieux cerner les stratégies de l'attaquant. Les architectures de déploiement des leurres se spécialisent selon le domaine d'application (systèmes d'information, systèmes industriels, finance, médical...) ou en fonction des composants ciblés par les attaques (serveurs, pare-feux, antivirus...) ou encore par rapport à la charge offensive (malwares...). Les leurres tendent à générer de vrais positifs en temps réel. Leur efficacité repose sur deux propriétés, l'une inhérente aux composants de sécurité, la non-compromission, et l'autre caractéristique de l'attaque : la furtivité.  
Cette nouvelle génération de leurres numériques enrichit les stratégies d'investigations au sein des centres opérationnels de sécurité (SOC).  Ainsi, des logiques de raisonnements déductifs (déterministes) ou inductifs (hypothétiques) se confrontent pour caractériser finement le mode opératoire des attaquants en le resserrant si possible jusqu'à l'attribution de l'attaque. 
Le caractère actif de ces nouvelles technologies de leurre soulève toutefois de nombreuses  interrogations sur le plan éthique et  réglementaire (respect de la vie privée en particulier).

Le leurrage numérique devient une composante essentielle d'une lutte informatique défensive et contribue de plus en plus au processus des scénarios de réponses, ripostes et d'escalade.

\begin{techworkbox}{\g{Deceptive Defense} en cyberdéfense}
Les techniques de déception en cyberdéfense sont en pleine évolution. C'est un sujet parfait pour \fichetech avec les différents thèmes  :
\begin{itemize}
  \item le leurrage numérique : honeypots, leurres, pièges
  \item 	les architectures de déploiement de leurres selon les domaines d'application
  \item 	la spécialisation de leurres pour les services, pour la sécurité, contre les malwares...
  \item 	les propriétés du leurrage numérique :  non-compromission,  furtivité...
  \item 	l'apport à l'investigation numérique : raisonnements déductifs/inductifs, caractérisation des attaques, attribution...
  \item 	la contribution à la lutte informatique défensive : scénarios de ripostes et d'escalade.
  \item 	le positionnement du leurrage dans les modèles d'attaques (MITRE ATT@CK,...), par rapport à la caractérisation des attaques (CAPEC...) et plus généralement son apport à la connaissance du risque cyber (cyber threat intelligence – CTI)
  \item 	le leurrage et la réglementation (NIS, RGPD,...).
\end{itemize}

\end{techworkbox}

%\subsection{Data Lake}
%\subsection{La SecOps par secteur industriel}








}
\mode<all>{%-------------------------------------------------------------
%               FR CYBERDEF SECOPS COURSE
%                                 CONTRIBS
%
%.    		   David BATTANY 2019/2020
%                     Architecture BOTNET
%-------------------------------------------------------------

\section{Quelques éléments techniques des attaques}

Afin d'illustrer les menaces et les attaques, je vous propose de présenter  avec un peu de détails quelques éléments techniques utilisés dans le cadre d'attaques informatiques classiques comme.

\begin{itemize}
  \item Botnet;
  \item Banker;
  \item Rançon-logiciel;
  \item  Techniques des malware ;
  \item Command and control.
\end{itemize}


toutefois, il est important de connaitre les grandes stratégies des attaquants.
En particulier les mécanismes de Kill Chain. Issu de la terminologie militaire ces éléments caractérisent et décrivent les mécanismes d'attaques.


%\section{Botnet, des codes malveillants organisés}.

A titre d'illustration, je vous propose de découvrir les botnets avec les travaux de David Batany.
Les Botnets font partis de ce que nous pourrions appeler des architectures techniques malveillantes. En effet, un Botnet en lui même est un code malveillant qui fonctionne dans une architecture informatique complexe avec des mécanismes de pilotage, de stockage, de réplication et d'activation spécialisées. Un Botnet est peut être considéré comme un système d'arme utilisable pour des attaques.

\section{MITRE Attack}

\utocomplete

%\newcommand{ \localpath}{../Contribs/david.batany@contribs.cyberdef/}
%\input{\localpath botnet.proxy.tex}


%Historique-motivations-Architecture_de_BOTNET

\subsection{Définition}
\par
Le terme \textbf{botnet}, contraction de l'anglais \textbf{robot+net}, se définit par l'ensemble des programmes, machines, serveurs connectés à internet ayant un ou plusieurs processus commun de communication.
Placé sous le contrôle d'un opérateur humain, appelé botmaster, le botnet recrute des machines en exploitant les vulnérabilités, failles, infections afin d'étendre son réseau à travers l'utilisation de canaux de Command and Control(C\&C).
\newline Avec l'IoT\footnote{Internet of Things}, et ses appareils connectés, le réseau s'étend de plus en plus au sein de notre société. L'actuelle faiblesse en terme de sécurité liée aux objets connectés représente une menace majeure et croissante dans notre environnement.
\head{Historique}
\par
Le concept, inventé en 1988 à l'université de Oulu en Finlande, fut développé à l'origine pour gérer les services associés au protocole IRC\footnote{Internet Relay Chat, un protocole de communication textuel}.
\newline
Le premier bot \g{<GM}> assistait ainsi l'utilisateur dans la gestion des connections IRC.
Cette gestion automatisée, permettant via un accès à distance, de contrôler et de réaliser des opérations a très vite montré un haut pouvoir malveillant.
\newline
En Mai 1999, Pretty Park, un malware de forme trojan horse se propageant sur le net permettait de voler les mots de passe.
\newline
Les premières dérives furent notamment l'affrontement de botnet IRC (\ulink{https://www.eggheads.org/}{Eggdrop} en décembre 1993, puis GTbot en avril 1998).
\head{Motivations liées à la menace botnet}
\begin{itemize}
	\item L'aspect lucratif représente l'intérêt majeur pour l'utilisateur de botnet. L'automatisation d'une tâche contrôlée à distance permettant de rapporter facilement des revenus (revente d'information, fraude au clic, spam), surtout si celle-ci est réalisée de manière anonyme(réseau TOR\footnote{The Onion Routing, un réseau d'anonymisation}).
	\item La motivation idéologique, comme par exemple, lors du conflit entre la Georgie et la Russie en 2008 ou de nombreux sites étatiques faisant l'objet de cyberattaques massives paralysaient les infrastructures.
	\item La motivation personnelle, à travers la vengeance ou le chantage, est également une finalité grâce notamment au caractère anonyme de l'attaque.
\end{itemize}


%Attaque-Architecture_de_BOTNET

\subsection{Les utilisations des Botnets}
On peut citer les attaques classiques utilisant des techniques de botnets:
\begin{itemize}
	\item Déni de service distribué (DDoS\footnote{Distributed Denial of Service}) ;
	\item DDoS contre paiement ;
	\item Cryptominage ;
	\item Récupération des identifiants;
	\item Exposition médiatique ou démonstration de force ;
	\item Dissimulation d’une autre attaque;
	\item Création d’un avantage concurrentiel ;
	\item Censure par attaque de serveurs;
	\item Vengeance par cryptolocker;
	\item Infrastructure d’anonymisation ;
	\item Recherche de vulnérabilités;
	\item Infrastructure d’anonymisation des communications;
	\item Contournement de mesures de limitation ou blocage ;
	\item Envoi de pourriels (Spam);
	\item Diffusion de codes malveillants;
	\item Exécution de codes malveillants sur les machines-zombies;
	\item Hébergement de codes malveillants ;
	\item Fraude aux clics ;
	\item Compromission d’accès ;
	\item Brute force hors-ligne;
	\item Brute force direct ...

\end{itemize}

%CycleDeVie-Architecture_de_BOTNET

\subsection{Cycle de vie d'une attaque:}
\par Pour la compréhension, il est nécessaire de comprendre les différentes étapes depuis l'infection jusqu'au fonctionnement complet du botnet.

\head{Infection de la machine}  \head{}
Cette première étape a généralement pour but de télécharger la charge virale sur un serveur.
Elle peut être initiée via les vecteurs suivants:
	 \begin{itemize}
		 \item Par spam (existence de spambot)
		 \item Exploitation de faille lié à la navigation sur un site web (malvertising\footnote{exploitation de pop-up publicitaires}, waterholing\footnote{ciblage de sites web fréquentés})
		 \item P2P
		 \item Spear phishing\footnote{hameçonnage ciblé pour récupérer données ou identifiants}
		 \item SMS, MMS
		 \item Bluetooth
		 \item TDS\footnote{Traffic Distribution Service, outil et service de redirection de trafic}
		 \item Exploit kits\footnote{plate-forme d'exploitation supporté par un site web permettant de tester une liste d'exploits}
	 \end{itemize}
	
\head{Activation}
Après téléchargement, l'installation du malware peut établir un premier contact avec le botnet (serveur dédié, servant-bot) ayant une fonctionnalité de C\&C.
Le téléchargement de \gls{gROOTKIT} d'installation ou de \gls{aDLL} complémentaire finalise la mise en place du botnet sur la machine infectée.	
	
\head{Mise à jour}
Les échanges permettent l'ajout de fonctionnalité, de configurations afin que le botnet puisse identifier et s'adapter à son environnement.
Il peut, par exemple, vouloir modifier son hash\footnote{signature numérique, ici on parle de signature virale} afin de conserver une certaine furtivité pour la continuité de l'attaque.

\head{Auto-protection}
La persistance et la dissimulation sont les facteurs clés de cette étape. L'installation de rootkit de protection, la modification du système, etc permettent de masquer l'action du botnet.
	
\head{Propagation}
Cette phase d'extension est à la fois locale par du scan et distante par diffusion virale (mail avec lien ou pièce jointe).
	
\head{Phase opérationnelle}
Cette dernière phase vise à accomplir les actions souhaitées de l'attaquant. 
Déclenchées, synchronisées ou persistantes ces attaques s'adaptent aux cibles désignées.Ordonné par le C\&C elles peuvent être activées ou mises en sommeil afin de ne pas attirer l'attention.




%Menace-Architecture_de_BOTNET

\subsection{Les type de menaces}
Les botnets représentent les outils de diffusion des attaques.
Cet outil permet aux cybercriminels de disposer d'un grand nombre de services développés dans un environnement collaboratif. Vendus sur le web, ils instrumentalisent l'attaque quelque soit le but recherché.
\newline \textbf{Liste des menaces possibles:}
\begin{itemize}
	\item Relayer du spam pour du commerce illégal ou pour de la manipulation d'information (par exemple des cours de bourse) 
  \item Réaliser des opérations d'hameçonnage 
  \item Identifier et infecter d’autres machines par diffusion de virus et de programmes malveillants (malwares) 
  \item Participer à des attaques groupées de déni de service (DDoS)
  \item Générer de façon abusive des clics sur un lien publicitaire au sein d’une page web (fraude au clic) 
  \item Capturer de l’information sur les machines compromises (vol puis revente d'information) ;
  \item Exploiter la puissance de calcul des machines ou effectuer des opérations de calcul distribué notamment pour cassage de mots de passe 
  \item Voler des sessions utilisateurs par credential stuffing ;
  \item Mener des opérations de commerce illicite en gérant l'accès à des sites de ventes de produits interdits ou de contrefaçons via des techniques de fast flux, simple ou double-flux ou RockPhish 
  \item Miner des cryptomonnaies, telles que le bitcoin1.
\end{itemize}

%Section-archi-aleatoire-Architecture_de_BOTNET

\newpage\subsection{Architecture aléatoire}

\upicture{Latex/Sources/EDU/Pictures/img-reseau_aleatoire}{Architecture aléatoire}{0.75}{lbl:netall}

%\begin{figure}[h]
%	\centering
%		\includegraphics[width=0.75\textwidth]{reseau_aleatoire}
%	\label{fig:réseau_aléatoire}
%	\caption[Architecture aléatoire]{Architecture aléatoire}
%\end{figure}

\subsubsection{Définition}
Ce concept représente une variante de l'architecture centralisée et peut se retrouver dans une variété de malware connue sous le nom de RATs\footnote{Remote Access Trojan}.Ces chevaux de Troie fonctionnent sur un principe de client/serveur parfois mis en place avec des techniques de social engineering (exemple: un fichier d'installation récupéré sur un site douteux). Ils exécutent la partie client à l'insu de l'utilisateur pour se connecter au serveur.
\newline Cette architecture tire profit de l'exploitation de plate-formes existantes supportant le protocole HTTP comme Facebook, Twitter, Yahoo, Evernote, Google, etc. et de réseau permettant de camoufler les échanges comme TOR, Hornet, etc.

\subsubsection {Liste des protocoles utilisés par le botnet}
\begin{itemize}
	\item HTTP
	\item protocole propriétaire (exemple XMPP\footnote{Extensible Messaging and Presence Protocol de MSN Messenger})
	\item IRC
\end{itemize}

\subsubsection{Avantages}
\begin{itemize}
	\item Architecture déjà existante
	\item Réseau important en terme d'utilisateurs
	\item Utilisation des canaux existants (exemple: messagerie instantanée)
	\item utilisation des fonctionnalités du réseau (exemple :Le réseau anonymisation TOR)
	\item Difficulté de démanteler son propre réseau
\end{itemize}

\subsubsection{Inconvénients}
\begin{itemize}
	\item Vulnérabilité du botnet face aux mécanismes de défense
	\item Connexion en permanence
	\item bloqué par le filtrage de liens de la plate-forme
\end{itemize}

%--------------------------------------------------------------------
%ajouter un botnet
%glisser le .csv dans le dossier exemples de la section 2
%inclure la ligne \csvautotabular[respect all]{Section/2-Ssection/exemples/nom-du-botnet}
%--------------------------------------------------------------------


\subsubsection{exemples}


%\noindent
%\resizebox{\textwidth}{!}{
%\csvautotabular[respect all]{Section/2-Ssection/exemples/Cutwail.csv}
%}
%\resizebox{\textwidth}{!}{
%\csvautotabular[respect all]{Section/2-Ssection/exemples/Carna.csv}
%}




%Section-archi-centralise-Architecture_de_BOTNET


\subsection[Architecture  centralisée]{Architecture  centralisée}

%\begin{figure}[h]
%	\centering
%		\includegraphics[width=0.75\textwidth]{archi_centralise}
%	\label{fig:archi_centralisée}
%	\caption[Architecture  centralisée]{Architecture  centralisée}
%\end{figure}


%\begin{figure}[h]
%	\centering
%		\includegraphics[width=0.75\textwidth]{archi_centralise_2}
%	\label{fig:Archi_centralisé_2}
%	
%	\caption[Architecture  centralisée répartie par utilisation d’une batterie de serveurs relais]{Architecture  centralisée répartie par utilisation d’une batterie de serveurs relais}
%\end{figure}

\subsubsection{Définition}
Un ou plusieurs nœuds de communication permettent aux bots d'échanger des données via un canal de communication.
Les nœuds représentent un serveur ou serveur relais avec comme fonction le C\&C.

\subsubsection{Liste des protocoles utilisés par le botnet}
\begin{itemize}
	\item IRC\footnote{Internet Relay Chat}
	\item HTTP\footnote{HyperText Transfer Protocol}
	\item IRC modifié
\end{itemize}

\subsubsection{Avantages}
\begin{itemize}
	\item Architecture  centraliséee
	\item Simplicité de mise en oeuvre (mIRC, ...)
	\item Utilisation des canaux IRC (topics, messages) pour l’envoi des commandes vers les botsPerformance (non gourmant en bande passante)
	\item Connexions régulières entre les bots et le C\&C (non-permanente)
	\item Recherche des ordres dans des forums, avec des mots clés ou même dans des images (stéganographie)
\end{itemize}

\subsubsection{Inconvénients}
\begin{itemize}
	\item Vulnérabilité du botnet (serveur central)
	\item Connexion en permanence
	\item Facile à détecter (filtrage du flux IRC)
\end{itemize}

%--------------------------------------------------------------------
%ajouter un botnet
%glisser le .csv dans le dossier exemples de la section 2
%inclure la ligne \csvautotabular[respect all]{Section/2-Ssection/exemples/nom-du-botnet}
%--------------------------------------------------------------------


%\subsubsection{exemples}
%\noindent
%\resizebox{\textwidth}{!}{
%\csvautotabular[respect all]{Section/2-Ssection/exemples/BEBLOH.csv}
%}
%\noindent
%\resizebox{\textwidth}{!}{
%\csvautotabular[respect all]{Section/2-Ssection/exemples/BOBAX-KRAKEN.csv}




%Section-archi-decentralise-Architecture_de_BOTNET


\subsection{Architecture décentraliséee}

%\begin{figure}[h]
%	\centering
%		\includegraphics[width=0.75\textwidth]{P2P_non_structure}
%	\label{fig:P2P_non_structure}
%	\caption{P2P non structuré}
%\end{figure}
%\begin{figure}[h]
%	\centering
%		\includegraphics[width=0.75\textwidth]{P2P_avec_super_pairs}
%	\label{fig:P2P_avec_super_pairs}
%	\caption{P2P avec super-pairs}
%\end{figure}

\subsubsection{Définition}
\par En utilisant des réseaux peer-to-peer, on s’affranchit d'un point central de communication.Chaque bot, selon ses caractéristiques, apporte les ressources pour élaborer le système de C\&C.
Il existe plusieurs typologies de réseau overlay\footnote{réseau logique de recouvrement}:
\begin{itemize}
	\item \textbf{overlay P2P non-structuré} : les topologies sont aléatoires (loi de puissance, aléatoire uniforme,...)
  \item \textbf{overlay P2P par super-pairs} : tous les pairs du réseau ne sont pas égaux, certains d’entre eux étant automatiquement sélectionnés pour servir temporairement le rôle de serveur pour les recherches ou le contrôle du réseau (comme FastTrack ou Gnutella)
  \item \textbf{overlay P2P structuré} : une cartographie établissant le lien entre le contenu et son emplacement ; ce type de réseau implémente en général – mais pas systématiquement une table de hachage distribuée (DHT) ; on retrouve dans cette catégorie les protocoles P2P Chord, Tapestry et Kademlia (utilisé par le logiciel eMule).
\end{itemize}

\subsubsection{Liste des protocoles utilisés par le botnet}
\begin{itemize}
	\item TCP/IP
	\item UDP
\end{itemize}

\subsubsection{Avantages}
\begin{itemize}
	\item Architecture décentralisée
	\item Indépendant de l’architecture DNS
	\item difficile à repérer
	\item Connexions régulières entre les bots et le C\&C (non permanente)
	\item Le botmaster donne les informations comme un bot faisant partie du réseau
  \item L’information transite de voisin en voisin
	\item très difficile à neutraliser
\end{itemize}

\subsubsection{Inconvénients}
\begin{itemize}
	\item Pas de vision globale du réseau par un bot
	\item Connexion en permanence
	\item Facile à détecter (filtrage du flux IRC)
\end{itemize}

%--------------------------------------------------------------------
%ajouter un botnet
%glisser le .csv dans le dossier exemples de la section 2
%inclure la ligne \csvautotabular[respect all]{Section/2-Ssection/exemples/nom-du-botnet}
%--------------------------------------------------------------------

%\subsubsection{exemples}
%\noindent
%\resizebox{\textwidth}{!}{
%\csvautotabular[respect all]{Section/2-Ssection/exemples/Storm.csv}
%}
%\resizebox{\textwidth}{!}{
%\csvautotabular[respect all]{Section/2-Ssection/exemples/QAKBOT.csv}
%}

%Section-archi-hybride-Architecture_de_BOTNET
\subsection{Architecture hybride}

%\begin{figure}[h]
%	\centering
%		\includegraphics[width=0.75\textwidth]{structure_hybride_P2P}
%	\label{fig:structure_hybride_P2P}
%	\caption[Structure hybride P2P]{Structure hybride P2P}
%\end{figure}

\subsubsection{Définition}
Une architecture hybride intègre une solution de repli, comme à l'aide d'un DGA\footnote{Domain Generation Algorithm} ou de plusieurs niveaux successifs entre le P2P et le C\&C.Cette hiérarchie permet de masquer une partie des adresses IP utilisées afin de rendre plus complexe l'analyse du botnet.
L'association de bot clients et de bots servent\footnote{serv-\textbf{eur} et cli-\textbf{ent}} met en évidence l'organisation du réseau par le botmaster.
Ces niveaux intermédiaires permettent de solidifier l'architecture du réseau.

\subsubsection{Liste des protocoles utilisés par le botnet}
Reprenant les protocoles présentés dans les architectures précédentes, L'efficacité, la furtivité et la complexité des méthodes de communication ont pour objectif de nuire aux efforts de démantèlement.
Ainsi, les coopérations internationales entre acteurs institutionnels et
privés sont généralement nécessaires pour permettre de démanteler les botnets les plus sophistiqués.

\subsubsection{Avantages}
\begin{itemize}
	\item Nombre important de domaines
  \item masquage de l'adresse IP
\end{itemize}

\subsubsection{Iconvénients}
\begin{itemize}
	\item tributaire de la bande passante
\end{itemize}

%--------------------------------------------------------------------
%ajouter un botnet
%glisser le .csv dans le dossier exemples de la section 2
%inclure la ligne \csvautotabular[respect all]{Section/2-Ssection/exemples/nom-du-botnet}
%--------------------------------------------------------------------

\subsubsection{exemples}

%\noindent
%\resizebox{\textwidth}{!}{
%\csvautotabular[respect all]{Section/2-Ssection/exemples/Gameover-Zeus.csv}
%}
%\resizebox{\textwidth}{!}{
%\csvautotabular[respect all]{Section/2-Ssection/exemples/Mirai.csv}
%}

%$$$$$$$$$$$$$$$$$$$$$$$$$$$$$$$$$$$$$$$$$$$$$$$$

%Analyse-Architecture_de_BOTNET

\subsection{L'analyse} 
\subsubsection{L'analyse statique}
Réalisée en sandbox, sur une machine virtuelle ou une machine dédiée avec des outils pré-installés comme InetSim,FakeNet ou Mozzle, elle débute par une analyse statique pour identifier les éléments et la composition du malware.
L'examen du code et des fonctions appelées permettent d'évaluer les capacités du botnet.
\subsubsection{L'analyse dynamique} 
Cette étape est relativement utile pour la compréhension de la menace car elle présente la machine infectée sous plusieurs états à l'aide de snapshots\footnote{copie des données/modifications apportées à un système}.
Ces captures instantanées situent l'avancement de l'infection lors de l'attaque.
Il est souvent nécessaire en présence d'algorithme chiffré d'utiliser cette méthode pour désobfusquer le code et comprendre la structure du malware.


%Defense-et-blocage-Architecture_de_BOTNET

\subsection{La défense et le blocage}
Au même titre que la protection contre les malwares, les recommandations en terme de SSI\ulink{https://www.ssi.gouv.fr/administration/bonnes-pratiques/}{recommandations ANSSI}, applicables localement, doivent s'inscrire dans nos habitudes et augmentent ainsi la probabilité de bloquer l'étape initiale de l'infection (spam, navigation non-sécurisée,etc.).
\newline Les mises à jour logicielles et système sont essentielles pour bloquer l'exploitation de CVE\footnote{Common Vulnerabilities and Exposures} \ulink{https://www.cvedetails.com}{Common Vulnerabilities and Exposures}.
\newline Au niveau du FAI, les notifications en cas de connexions malveillantes et la surveillance des adresses IP sont un frein à l'extension du botnet.
\newline Enfin la détection d'un appel de fonctions anormales par l'antivirus, l'autorisation et l’identification des flux sortants par le firewall permettent le blocage de l'activité malveillante.
Les fonctionnalités recherchées de l'antivirus dans ce cadre sont un firewall bidirectionnel, une protection contre le phishing, la vérification de la certification, la lutte contre le tracking, la vérification du téléchargement, le blocage des pop-ups et pages WEB malveillantes,etc. 

%Demantelement-Architecture_de_BOTNET
\subsection{Le démantèlement}
Les méthodes de détection impliquent parfois des actions offensives visant à entraver le développement du botnet.
Il est cependant nécessaire d'avoir un support juridique et judiciaire pour mener à bien des actions adaptées au type et à la taille du botnet.
Celles-ci sont généralement menées en coopération avec les industriels (Microsoft,Level 3, Cisco,etc.) et la communauté scientifique.

%Detection-Architecture_de_BOTNET

\subsection{La détection}
Recherche d'anomalie, comparaison de signatures, pots-de-miel, toutes ces techniques basées sur l'activité du réseau reposent sur l'inspection des flux et des paquets.

\subsubsection{La détection passive}
L'identification et l'analyse passive des flux (adresses IP,port source et destination, étiquette MPLS\footnote{MultiProtocol Label Switching}, etc.) permettent de classifier les protocoles suspects et les serveurs de C\&C.
\newline BotFinder, par exemple, permet de décomposer un flux (durée moyenne des connexions, nombre d'octets transférés, etc.) et de le comparer à l'activité normale du réseau.
\newline L'observation des DNS\footnote{Domain Name Server} permet aussi d'identifier les domaines malveillants afin de caractériser le botnet suspecté.
BotGad, un système d'exploitation permettant d'analyser le trafic DNS sur un réseau local, utilise un algorithme basé sur l'apprentissage afin de définir la stratégie de groupe du botnet.
\newline EXPOSURE, un autre système d'exploitation déployé au sein de l'ISP\footnote{Internet Service Provider}, permet d'analyser à large échelle mais sur une durée de plusieurs mois le trafic DNS. Produisant une liste de domaines malveillants, il permet, par exemple, d'identifier un volume conséquent de requêtes pour un même domaine.
\newline Enfin le recours aux pots-de-miel et l'analyse des journaux d'activité sont les éléments de base d'une recherche d'activité liée aux botnets.
Suivant cette idée, le SIEM\footnote{Security Information and Event Management}, une solution de gestion de la sécurité, représente un outil précieux et novateur afin d'optimiser la veille du trafic et d'automatiser les processus de sécurité en cas de comportement anormal.

\subsubsection{La détection active}
Différentes techniques existent comme le sinkholing\footnote{également appelé serveur gouffre, gouffre Internet ou Blackhole DNS } redirigeant le trafic vers des serveurs afin de simuler le comportement du C\&C et de diminuer la puissance du réseau du botnet.
\newline L'infiltration, fonction de l'architecture du botnet, consiste à simuler le comportement d'un botnet contrôlé à l'aide de drones IRC ou de script afin de capturer le trafic et de remonter jusqu'au botmaster.
Le projet Pebbletrace reprend cette idée d’identification du botmaster en piégeant les équipements infectés avec une charge défensive afin de retourner le trafic contre le botnet.


%synthese
Les botnets font maintenant partie d'une économie souterraine sous forme de services payants.
Le harcèlement, le vol, les attaques par déni de service, etc figurent comme des produits de vente accessibles, moyennant finances, pour n'importe quel criminel.
\newline D'après l'ENISA\footnote{European Union Agency for Cybersecurity, anciennement European Network and Information Security Agency} les prix varient suivant la fiabilité, la durée et le type de service requis.
Par exemple une heure de DDoS est disponible pour 38\$.
Les différentes architectures permettent de mieux comprendre l'organisation du botnet et l'importance du C\&C.
\vspace{5mm}
\newline Les notions de veille technologique et de partage d'informations sur la menace sont essentielles du fait de l'implication du botnet dans les réseaux publiques et privés.
\newline Dans le cadre du renseignement lié aux menaces, les constructeurs de smartphone fournissent les informations (recherches, cibles des attaques, menaces associés aux mobiles, vulnérabilités des objets IoT) issues de leurs Threat Intelligence Center\footnote{Centre de renseignements liés aux menaces cyber}.
%ouverture
\vspace{5mm}
\newline Selon Nokia (\ulink{https://www.nokia.com/about-us/news/releases/2018/12/04/nokias-threat-intelligence-report-2019-warns-on-the-fast-growing-and-evolving-threat-of-malicious-software-targeting-internet-of-things-iot-devices/}{Nokia's Threat Intelligence Report}, l'activité des botnet sur l'IoT représente 78\% des événements de détection de malware en 2018.
\newline La menace omniprésente de ces objets connectés (santé, domotique, médias, électroménager,...) n'est aujourd'hui pas assez prise en considération par notre société de consommation.
Le manque de sécurité accrue de ces appareils fait apparaître ces objets connectés comme des acteurs potentiels constituant le réseau d'un botnet.
\newline L'arrivée prochaine de la 5G\footnote{prévisions en moyenne 100Mbit/S en download et 50 Mbit/s en upload} sera, dans ce domaine, un vecteur majeur pour la diffusion de l'infection du malware et le nombre d'attaques associés au botnet (exemple attaques DDoS).


%\begin{figure}[h!]
%\begin{center}
%	\includegraphics[width =0.6\textwidth, keepaspectratio]{coverback}
%	\caption{Figure 1}\label{mon image}
%	\end{center}
%\end{figure}

\section{Malware}

\subsection{Malware industriel : citadel}

Les créateurs du malware Citadel ont adopté un modèle de développement Open Source qui leur permet de corriger collectivement les bugs et d'ajouter plus rapidement des fonctionnalités à leur logiciel malveillant. Citadel est basé sur ZeuS, l'un des plus anciens et des plus populaires Cheval de Troie mis au point pour pirater les services bancaires en ligne. Zeus a été abandonné par son créateur fin 2010, mais quelques mois plus tard, le code source de son malware a été diffusé sur le Net. Depuis la publication de son code source, Zeus a servi de base au développement d'autres Trojan comme Ice IX,  et Citadel. « Le 17 décembre 2011, pour la première fois, un laboratoire de recherche (Seculert) a mis en évidence l'existence d'un botnet Citadel, » a déclaré le spécialiste de la sécurité dans un blog. « Le taux d'adoption et le développement de ce malware est en pleine expansion. » Seculert dit avoir identifié plus de 20 réseaux de zombies utilisant des versions différentes du Cheval de Troie. « Chaque version comporte de nouveaux modules et de nouvelles fonctionnalités, certaines proposées par les clients de Citadel eux-mêmes, » a indiqué l'entreprise de sécurité. L'aspect le plus intéressant de ce malware est sans aucun doute le processus adopté pour son développement. Les modalités de développement adoptées pour Citadel ressemblent à celles des communautés supportant des projets Open Source. « L'organisation est calquée sur le développement des logiciels courants : les créateurs de Citadel fournissent à leurs clients un manuel utilisateur, des avis donnant des détails sur les mises à jour et même un contrat de licence.}


%$END_UBODY
}%---------------------------------------------------------------------

%************************************************************
% Chargement  du MODELE
%************************************************************

\umainload



